\input{preamble}

% OK, start here.
%
\begin{document}

\title{Topologie}


\maketitle

\phantomsection
\label{section-phantom}

\tableofcontents

\section{Introduction}
\label{section-introduction}

\noindent
Nous exposerons dans ce chapitre les éléments de topologie générale dont nous
aurons besoin. On pourra consulter \cite{Engelking} comme ouvrage de référence.

\section{Notions de base}
\label{section-topology-basic}

\noindent
Voici une liste de notions fondamentales de topologie. Certaines seront
étudiées plus en détail dans la suite, certaines sont définies dans la liste,
mais les autres sont tenues pour acquises et ne seront pas définies. Si la
plupart des notions en italique ne vous sont pas familières, nous vous
conseillons de consulter un ouvrage d'introduction à la topologie avant de
poursuivre.

\begin{enumerate}
\item
\label{item-space}
$X$ est un {\it espace topologique},
\item
\label{item-point}
$x\in X$ est un {\it point},
\item
\label{item-locally-closed}
$E \subset X$ est une partie {\it localement fermée},
\item
\label{item-closed-point}
$x\in X$ est un {\it point fermé},
\item
\label{item-dense}
$E \subset X$ est une partie {\it dense},
\item
\label{item-continuous}
$f : X_1 \to X_2$ est {\it continue},
\item
\label{item-upper-semi-continuous}
une fonction à valeurs réelles étendues
$f : X \to \mathbf{R} \cup \{\infty, -\infty\}$
est {\it semi-continue supérieurement} si $\{x \in X \mid f(x) < a\}$ est
ouvert pour tout $a \in \mathbf{R}$,
\item
\label{item-lower-semi-continuous}
une fonction à valeurs réelles étendues
$f : X \to \mathbf{R} \cup \{\infty, -\infty\}$
est {\it semi-continue inférieurement} si $\{x \in X \mid f(x) > a\}$ est
ouvert pour tout $a \in \mathbf{R}$,
\item une application continue d'espaces $f : X \to Y$ est {\it ouverte} si
$f(U)$ est ouvert dans $Y$ pour toute partie ouverte $U \subset X$,
\item une application continue d'espaces $f : X \to Y$ est {\it fermée} si
$f(Z)$ est fermé dans $Y$ pour toute partie fermée $Z \subset X$,
\item
\label{item-neighbourhood}
un {\it voisinage de $x \in X$} est toute partie $E \subset X$ qui contient un
ouvert contenant $x$,
\item
\label{item-induced-topology}
la {\it topologie induite} sur une partie $E \subset X$,
\item
\label{item-covering}
$\mathcal{U} : U = \bigcup_{i \in I} U_i$ est un {\it recouvrement ouvert de}
$U$ (remarquons que chaque $U_i$ peut être vide et que, lorsque $U$ est vide,
nous admettons même que $I$ soit l'ensemble vide),
\item
\label{item-subcover}
un {\it sous-recouvrement} d'un recouvrement comme en
(\ref{item-covering}) est un recouvrement ouvert
$\mathcal{U}' : U = \bigcup_{i \in I'} U_i$, où $I' \subset I$,
\item
\label{item-refinement}
le recouvrement ouvert $\mathcal{V}$ est un {\it raffinement} du recouvrement
ouvert $\mathcal{U}$ (si $\mathcal{V} : U = \bigcup_{j \in J} V_j$ et
$\mathcal{U} : U = \bigcup_{i \in I} U_i$, cela signifie que chaque $V_j$ est
contenu dans l'un des $U_i$),
\item
\label{item-fundamental-system}
{\it $\{ E_i \}_{i \in I}$ est un système fondamental de voisinages de $x$
dans $X$},
\item
\label{item-Hausdorff}
un espace topologique $X$ est dit {\it séparé}, ou {\it de Hausdorff}, si et
seulement si, pour toute paire de points distincts $x, y \in X$, il existe des
ouverts disjoints $U, V \subset X$ tels que $x \in U$, $y \in V$,
\item le {\it produit} de deux espaces topologiques,
\label{item-product}
\item
\label{item-fibre-product}
le {\it produit fibré $X \times_Y Z$} de deux applications continues
$f : X \to Y$ et $g : Z \to Y$,
\item
\label{item-discrete-indiscrete}
la {\it topologie discrète} et la {\it topologie grossière} sur un ensemble,
\item etc.
\end{enumerate}



\section{Espaces séparés}
\label{section-Hausdorff}

\noindent
La catégorie des espaces topologiques admet les produits finis.

\begin{lemma}
\label{lemma-Hausdorff}
Soit $X$ un espace topologique. Les assertions suivantes sont équivalentes :
\begin{enumerate}
\item $X$ est séparé,
\item la diagonale $\Delta(X) \subset X \times X$ est fermée.
\end{enumerate}
\end{lemma}

\begin{proof}
Supposons que $X$ soit séparé. Soit $(x, y) \not \in \Delta(X)$,
c'est-à-dire $x \neq y$. Il existe deux ouverts disjoints $U$ et $V$ de $X$
tels que $x \in U$ et $y \in V$. Il s'ensuit que $U \times V \subset
(X \times X) \setminus \Delta(X) $. Cela montre que $ (X \times X) \setminus
\Delta(X)$ est un ouvert de $ X \times X$, ce qui revient à dire que
la diagonale $\Delta(X) \subset X \times X$ est fermée dans $X \times X$.
La réciproque est analogue :
le complémentaire $(X \times X) \setminus \Delta(X)$ est formé
précisément des $(x, y) \in X\times X$ tels que $ x \neq y$. Ainsi, si
$ \Delta(X) \subset X \times X$ est fermée, alors, par définition de la
topologie produit, pour tout tel couple $ (x, y)$, il existe des ouverts
$U, V \subset X$ tels que $(x, y) \in U \times V$ et
$(U \times V)\cap \Delta(X) = \emptyset$. Autrement dit, $x \in U$ et
$y \in V$, avec $U \cap V = \emptyset$.
\end{proof}

\begin{lemma}
\label{lemma-graph-closed}
\begin{slogan}
Les graphes des applications à valeurs dans un espace séparé sont fermés.
\end{slogan}
Soit $f : X \to Y$ une application continue d'espaces topologiques.
Si $Y$ est séparé, alors le graphe de $f$ est fermé dans $X \times Y$.
\end{lemma}

\begin{proof}
Le graphe est l'image réciproque de la diagonale par l'application
$X \times Y \to Y \times Y$. Le lemme résulte donc du
lemme \ref{lemma-Hausdorff}.
\end{proof}

\begin{lemma}
\label{lemma-section-closed}
Soit $f : X \to Y$ une application continue d'espaces topologiques.
Soit $s : Y \to X$ une application continue telle que
$f \circ s = \text{id}_Y$.
Si $X$ est séparé, alors $s(Y)$ est fermé.
\end{lemma}

\begin{proof}
Cela résulte du lemme \ref{lemma-Hausdorff}, puisque
$s(Y) = \{x \in X \mid x = s(f(x))\}$.
\end{proof}

\begin{lemma}
\label{lemma-fibre-product-closed}
Soient $X \to Z$ et $Y \to Z$ des applications continues d'espaces
topologiques. Si $Z$ est séparé, alors $X \times_Z Y$ est fermé dans
$X \times Y$.
\end{lemma}

\begin{proof}
Cela résulte du lemme \ref{lemma-Hausdorff}, car $X \times_Z Y$ est l'image
réciproque de $\Delta(Z)$ par $X \times Y \to Z \times Z$.
\end{proof}




\section{Applications séparées}
\label{section-separated}

\noindent
Voici simplement la définition et quelques lemmes élémentaires.

\begin{definition}
\label{definition-separated}
Une application continue $f : X \to Y$ d'espaces topologiques est dite
{\it séparée} si et seulement si la diagonale $\Delta : X \to X \times_Y X$
est une application fermée.
\end{definition}

\begin{lemma}
\label{lemma-separated}
Soit $f : X \to Y$ une application continue d'espaces topologiques.
Les assertions suivantes sont équivalentes :
\begin{enumerate}
\item $f$ est séparée,
\item $\Delta(X) \subset X \times_Y X$ est une partie fermée,
\item étant donnés deux points distincts $x, x' \in X$ ayant la même image
dans $Y$, il existe des voisinages ouverts disjoints de $x$ et de $x'$.
\end{enumerate}
\end{lemma}

\begin{proof}
Si $f$ est séparée, la définition \ref{definition-separated} dit que $\Delta$
est une application fermée. Le fait que $X$ soit fermé dans $X$ entraîne que
$\Delta(X)$ est fermé dans $ X\times_Y X $. Ainsi, (1) implique (2).

\medskip\noindent
Supposons que $\Delta(X) \subset X \times_Y X$ soit une partie fermée et
notons $U$ son complémentaire ouvert. Il existe donc un ouvert
$W \subset X \times X$ tel que $W \cap (X \times_Y X) = U$. Or, par
définition de la topologie produit, si
$(x, x') \in W \cap (X \times_Y X)$,
il existe des ouverts $V$ et $V'$ de $X$ tels que $x \in V$, $x' \in V'$ et
$V \times V' \subset W$. Si $V \cap V' \neq \emptyset$, il existerait
$z \in V \cap V'$. Mais $(z, z) \in X \times_Y X$, donc
$(z,z) \in (V \times V') \cap (X\times_Y X) \subset U$, ce qui est absurde.
Ainsi $V \cap V' = \emptyset $, et nous avons obtenu deux voisinages ouverts
disjoints de $x$ et $x'$. Cela montre que (2) implique (3).

\medskip\noindent
Supposons enfin que, pour tous points distincts $x, x' \in X$ ayant la même
image dans $Y$, il existe des voisinages ouverts disjoints de $x$ et de $x'$.
Soit $F$ un fermé de $X$. Montrons que $\Delta(F)$ est une partie fermée de
$X \times_Y X$. Soit $(x, x') \in X \times_Y X$ un point qui n'appartient
pas à $\Delta(F)$. Alors soit $x \not = x'$, soit $x \not \in F$.
Dans le premier cas, choisissons des voisinages ouverts disjoints
$V, V' \subset X$ de $x, x'$ ; alors $(V \times V') \cap X \times_Y X$
est un voisinage ouvert de $(x, x')$ qui ne rencontre pas $\Delta(F)$.
Dans le second cas, $((X \setminus F) \times X) \cap X \times_Y X$
est un voisinage ouvert de $(x, x')$ qui ne rencontre pas $\Delta(F)$.
Nous avons montré que (3) implique (1).
\end{proof}

\begin{lemma}
\label{lemma-from-hausdorff}
Soit $f : X \to Y$ une application continue d'espaces topologiques.
Si $X$ est séparé, alors $f$ est séparée.
\end{lemma}

\begin{proof}
C'est immédiat d'après les lemmes \ref{lemma-separated} et
\ref{lemma-Hausdorff}, car le fait que $\Delta(X)$ soit fermé dans
$X \times X$ entraîne que $\Delta(X)$ est fermé dans $X \times_Y X$.
\end{proof}

\begin{lemma}
\label{lemma-base-change-separated}
Soient $f : X \to Y$ et $Y' \to Y$ des applications continues d'espaces
topologiques. Si $f$ est séparée, alors
$f' : Y' \times_Y X \to Y'$ est séparée.
\end{lemma}

\begin{proof}
Cela résulte de la caractérisation (2) du lemme \ref{lemma-separated}.
En effet, si $X' = Y' \times_Y X$, la diagonale $\Delta(X')$ dans le produit
fibré $X' \times_{Y'} X'$ est l'image réciproque de $\Delta(X)$ dans
$X \times_Y X$.
\end{proof}




\section{Bases}
\label{section-bases}

\noindent
Rappels élémentaires sur les bases des espaces topologiques.

\begin{definition}
\label{definition-base}
Soit $X$ un espace topologique. Une famille de parties $\mathcal{B}$ de $X$
est appelée une {\it base de la topologie sur $X$}, ou une {\it base de la
topologie de $X$}, si les conditions suivantes sont satisfaites :
\begin{enumerate}
\item Tout élément $B \in \mathcal{B}$ est ouvert dans $X$.
\item Pour tout ouvert $U \subset X$ et tout $x \in U$, il existe
$B \in \mathcal{B}$ tel que $x \in B \subset U$.
\end{enumerate}
\end{definition}

\noindent
Le lemme suivant sert parfois à définir une topologie.

\begin{lemma}
\label{lemma-make-base}
Soit $X$ un ensemble et soit $\mathcal{B}$ une famille de parties.
Supposons que $X = \bigcup_{B \in \mathcal{B}} B$ et que, pour tous
$x \in B_1 \cap B_2$ avec $B_1, B_2 \in \mathcal{B}$, il existe
$B_3 \in \mathcal{B}$ tel que $x \in B_3 \subset B_1 \cap B_2$.
Il existe alors une unique topologie sur $X$ dont $\mathcal{B}$ est une base.
\end{lemma}

\begin{proof}
Notons $\sigma(\mathcal B)$ l'ensemble des parties de $X$ qui s'écrivent
comme réunions d'éléments de $\mathcal B$. Nous affirmons que
$\sigma(\mathcal{B})$ est une topologie. En effet, l'ensemble vide appartient
à $\sigma(\mathcal{B})$ (en tant que réunion vide), et $X$ appartient à
$\sigma(\mathcal{B})$ (en tant que réunion de tous les éléments de
$\mathcal{B}$). Il est clair que $\sigma(\mathcal{B})$ est stable par réunion.
Enfin, si $U, V \in \sigma(\mathcal{B})$, écrivons
$U = \bigcup_{i \in I} U_i$ et $V = \bigcup_{j \in J} V_j$, où
$U_i, V_j \in \mathcal{B}$ pour tout $i \in I$ et tout $j \in J$.
Alors
$$
U \cap V = \bigcup\nolimits_{i \in I, j \in J}
U_i \cap V_j
$$
L'hypothèse du lemme montre que $U_i \cap V_j \in \sigma(\mathcal{B})$,
et donc que $U \cap V$ y appartient également. Ainsi,
$\sigma(\mathcal{B})$ est une topologie.
Les propriétés (1) et (2) de la définition \ref{definition-base} sont
immédiates pour cette topologie. Pour en montrer l'unicité, soit $\mathcal T$
une topologie sur $X$ dont $\mathcal B$ est une base. Tout élément de
$\mathcal{B}$ appartient évidemment à $\mathcal{T}$ d'après le point (1) de
la définition \ref{definition-base}, et donc
$\sigma(\mathcal{B}) \subset \mathcal{T}$. Réciproquement, le point (2) de
la définition \ref{definition-base} montre que tout élément de $\mathcal{T}$
est une réunion d'éléments de $\mathcal{B}$, c'est-à-dire que
$\mathcal{T} \subset \sigma(\mathcal{B})$. Cela achève la démonstration.
\end{proof}

\begin{lemma}
\label{lemma-refine-covering-basis}
Soit $X$ un espace topologique.
Soit $\mathcal{B}$ une base de la topologie de $X$.
Soit $\mathcal{U} : U = \bigcup_i U_i$ un recouvrement ouvert de
$U \subset X$. Il existe un recouvrement ouvert $U = \bigcup V_j$ qui
raffine $\mathcal{U}$ et tel que chaque $V_j$ appartienne à la base
$\mathcal{B}$.
\end{lemma}

\begin{proof}
Si $ x \in U = \bigcup_{i\in I} U_i $, il existe $ i_x \in I $ tel que
$ x \in U_{i_x} $. Il existe donc $ B_{i_x} \in \mathcal{B}$ vérifiant
$ x \in B_{i_x} \subset U_{i_x}$. Posons
$J = \{i_x | x \in U\}$ et, pour $j = i_x \in J$, posons
$V_j = B_{i_x}$. On obtient ainsi le recouvrement ouvert voulu de $U$ par
$\{V_j\}_{j \in J}$.
\end{proof}

\begin{definition}
\label{definition-subbase}
Soit $X$ un espace topologique. Une famille de parties $\mathcal{B}$ de $X$
est appelée une {\it sous-base de la topologie sur $X$}, ou une {\it sous-base
de la topologie de $X$}, si les intersections finies d'éléments de
$\mathcal{B}$ forment une base de la topologie de $X$.
\end{definition}

\noindent

En particulier, tout élément de $\mathcal{B}$ est ouvert.

\begin{lemma}
\label{lemma-subbase}
Soit $X$ un ensemble. Pour toute famille $\mathcal{B}$ de parties de $X$,
il existe une unique topologie sur $X$ dont $\mathcal{B}$ est une sous-base.
\end{lemma}

\begin{proof}
Par convention, $\bigcap_\emptyset B = X $. Nous pouvons donc appliquer le
lemme \ref{lemma-make-base} à l'ensemble des intersections finies d'éléments
de $\mathcal B$.
\end{proof}

\begin{lemma}
\label{lemma-create-map-from-subcollection}
Soit $X$ un espace topologique. Soit $\mathcal{B}$ une famille d'ouverts de
$X$. Supposons que $X = \bigcup_{U \in \mathcal{B}} U$ et que, pour
$U, V \in \mathcal{B}$, on ait
$U \cap V = \bigcup_{W \in \mathcal{B}, W \subset U \cap V} W$.
Il existe alors une application continue $f : X \to Y$ d'espaces
topologiques telle que
\begin{enumerate}
\item pour $U \in \mathcal{B}$, l'image $f(U)$ soit ouverte,
\item pour $U \in \mathcal{B}$, on ait $f^{-1}(f(U)) = U$, et
\item les ouverts $f(U)$, $U \in \mathcal{B}$, forment une base de la
topologie de $Y$.
\end{enumerate}
\end{lemma}

\begin{proof}
Définissons une relation d'équivalence $\sim$ sur les points de $X$ par
$$
x \sim y \Leftrightarrow
(\forall U \in \mathcal{B} : x \in U \Leftrightarrow y \in U)
$$
Soit $Y$ l'ensemble des classes d'équivalence et soit $f : X \to Y$
l'application naturelle. Le point (2) est satisfait par construction.
Les hypothèses portant sur $\mathcal{B}$ sont exactement celles du
lemme \ref{lemma-make-base} pour l'ensemble des parties $f(U)$,
$U \in \mathcal{B}$. Il existe donc une unique topologie sur $Y$ pour
laquelle (3) soit satisfaite. Le point (1) est alors clair lui aussi.
\end{proof}






\section{Applications submersives}
\label{section-submersive}

\noindent
Si $X$ est un espace topologique et si $E \subset X$ est une partie, nous
munissons habituellement $E$ de la {\it topologie induite}.

\begin{lemma}
\label{lemma-induced}
Soit $X$ un espace topologique. Soit $Y$ un ensemble et soit
$f : Y \to X$ une application injective d'ensembles. La topologie induite sur
$Y$ est la topologie caractérisée par chacune des assertions suivantes :
\begin{enumerate}
\item c'est la topologie la moins fine sur $Y$ pour laquelle $f$ soit continue,
\item les ouverts de $Y$ sont les $f^{-1}(U)$, où $U \subset X$ est ouvert,
\item les fermés de $Y$ sont les ensembles $f^{-1}(Z)$, où $Z \subset X$
est fermé.
\end{enumerate}
\end{lemma}

\begin{proof}
L'ensemble $\mathcal{T} = \{f^{-1}(U) | U \subset X \text{ ouvert}\}$ est
une topologie sur $Y$. Tout d'abord,
$\emptyset = f^{-1}(\emptyset)$ et $f^{-1}(X) = Y$.
Ainsi, $\mathcal{T}$ contient $\emptyset$ et $Y$.

\medskip\noindent
Soit maintenant $\{V_i\}_{i \in I}$ une famille de parties ouvertes avec
$V_i \in \mathcal{T}$, et écrivons $V_i = f^{-1}(U_i)$, où $U_i$ est une
partie ouverte de $X$. Alors
$$
\bigcup\nolimits_{i\in I} V_i =
\bigcup\nolimits_{i\in I} f^{-1}(U_i) =
f^{-1}\left(\bigcup\nolimits_{i\in I} U_i\right)
$$
Ainsi, $\bigcup_{i\in I} V_i \in \mathcal{T}$, puisque
$\bigcup_{i\in I} U_i$ est ouvert dans $X$.
Soient maintenant $V_1, V_2 \in \mathcal{T}$. Il existe des ouverts
$U_1, U_2$ de $X$ tels que $V_1 = f^{-1}(U_1)$ et
$V_2 = f^{-1}(U_2)$. Alors
$$
V_1 \cap V_2 = f^{-1}(U_1) \cap f^{-1}(U_2) = f^{-1}(U_1 \cap U_2)
$$
Ainsi, $V_1 \cap V_2 \in \mathcal{T}$, car $U_1 \cap U_2$ est ouvert dans
$X$.

\medskip\noindent
Toute topologie sur $Y$ pour laquelle $f$ est continue contient
$\mathcal{T}$, par définition d'une application continue. Ainsi,
$\mathcal{T}$ est bien la topologie la moins fine sur $Y$ pour laquelle $f$
soit continue. Cela prouve l'équivalence de (1) et (2).

\medskip\noindent
L'équivalence de (2) et (3) résulte de l'égalité
$f^{-1}(X \setminus E) = Y \setminus f^{-1}(E)$ pour toute partie
$E \subset X$.
\end{proof}

\noindent
Duellement, si $X$ est un espace topologique et si $X \to Y$ est une
surjection d'ensembles, on peut munir $Y$ de la {\it topologie quotient}.

\begin{lemma}
\label{lemma-quotient}
Soit $X$ un espace topologique. Soit $Y$ un ensemble et soit $f : X \to Y$
une application surjective d'ensembles. La topologie quotient sur $Y$ est la
topologie caractérisée par chacune des assertions suivantes :
\begin{enumerate}
\item c'est la topologie la plus fine sur $Y$ pour laquelle $f$ soit continue,
\item une partie $V$ de $Y$ est ouverte si et seulement si $f^{-1}(V)$ est
ouverte,
\item une partie $Z$ de $Y$ est fermée si et seulement si $f^{-1}(Z)$ est
fermée.
\end{enumerate}
\end{lemma}

\begin{proof}
L'ensemble $\mathcal{T} = \{V\subset Y | f^{-1}(V) \text{ est ouvert}\}$ est
une topologie sur $Y$. Tout d'abord,
$\emptyset = f^{-1}(\emptyset)$ et $f^{-1}(Y) = X$.
Ainsi, $\mathcal{T}$ contient $\emptyset$ et $Y$.

\medskip\noindent
Soit $(V_i)_{i \in I}$ une famille d'éléments $V_i \in \mathcal{T}$. Alors
$$
\bigcup\nolimits_{i\in I} f^{-1}(V_i) =
f^{-1}\left(\bigcup\nolimits_{i\in I} V_i\right)
$$
Ainsi, $\bigcup_{i\in I} V_i \in \mathcal{T}$, puisque
$\bigcup_{i\in I}f^{-1}(V_i)$ est ouvert dans $X$.
De plus, si $V_1, V_2 \in \mathcal{T}$, alors
$$
f^{-1}(V_1) \cap f^{-1}(V_2) = f^{-1}(V_1 \cap V_2)
$$
Ainsi, $V_1 \cap V_2 \in \mathcal{T}$, car
$f^{-1}(V_1) \cap f^{-1}(V_2)$ est ouvert dans $X$.

\medskip\noindent
Enfin, toute topologie sur $Y$ pour laquelle $f$ est continue est contenue
dans $\mathcal{T}$, par définition d'une application continue. Par conséquent,
$\mathcal{T}$ est la topologie la plus fine sur $Y$ pour laquelle $f$ soit
continue. Cela prouve l'équivalence de (1) et (2).

\medskip\noindent
Enfin, l'équivalence de (2) et (3) résulte de
$f^{-1}(X\setminus E) = Y \setminus f^{-1}(E)$ pour toute partie
$E \subset X$.
\end{proof}

\noindent
Soit $f : X \to Y$ une application continue d'espaces topologiques.
Nous obtenons alors une factorisation $X \to f(X) \to Y$ d'applications
d'ensembles. Nous pouvons munir $f(X)$ soit de la topologie quotient provenant
de la surjection $X \to f(X)$, soit de la topologie induite provenant de
l'injection $f(X) \to Y$. L'application
$$
(f(X), \text{topologie quotient})
\longrightarrow
(f(X), \text{topologie induite})
$$
est continue.

\begin{definition}
\label{definition-submersive}
Soit $f : X \to Y$ une application continue d'espaces topologiques.
\begin{enumerate}
\item On dit que $f$ est une {\it application stricte d'espaces topologiques}
si la topologie induite et la topologie quotient sur $f(X)$ coïncident
(voir la discussion ci-dessus).
\item On dit que $f$ est {\it submersive}\footnote{Cette notion est très
différente de celle de submersion entre variétés différentielles ! Il est sans
doute préférable d'employer « stricte et surjective » plutôt que
« submersive ».} si $f$ est surjective et stricte.
\end{enumerate}
\end{definition}

\noindent
Ainsi, une application continue $f : X \to Y$ est submersive si $f$ est une
surjection et si, pour toute partie $T \subset Y$, la partie $T$ est ouverte
(respectivement fermée) si et seulement si $f^{-1}(T)$ est ouverte
(respectivement fermée). Autrement dit, $Y$ porte la
topologie quotient relative à la surjection $X \to Y$.

\begin{lemma}
\label{lemma-open-morphism-quotient-topology}
Soit $f : X \to Y$ une application surjective, ouverte et continue d'espaces
topologiques. Soit $T \subset Y$ une partie. Alors
\begin{enumerate}
\item $f^{-1}(\overline{T}) = \overline{f^{-1}(T)}$,
\item $T \subset Y$ est fermé si et seulement si $f^{-1}(T)$ est fermé,
\item $T \subset Y$ est ouvert si et seulement si $f^{-1}(T)$ est ouvert, et
\item $T \subset Y$ est localement fermé si et seulement si $f^{-1}(T)$
est localement fermé.
\end{enumerate}
En particulier, $f$ est submersive.
\end{lemma}

\begin{proof}
Il est clair que $\overline{f^{-1}(T)} \subset f^{-1}(\overline{T})$.
Si $x \in X$ et $x \not \in \overline{f^{-1}(T)}$, il existe un voisinage
ouvert $x \in U \subset X$ tel que $U \cap f^{-1}(T) = \emptyset$.
Comme $f$ est ouverte, $f(U)$ est un voisinage ouvert de $f(x)$ qui ne
rencontre pas $T$. Ainsi, $x \not \in f^{-1}(\overline{T})$. Cela prouve (1).
Le point (2) résulte immédiatement de (1).
Le point (3) est évident puisque $f$ est ouverte et surjective.
Pour (4), si $f^{-1}(T)$ est localement fermé, alors
$f^{-1}(T) \subset \overline{f^{-1}(T)} = f^{-1}(\overline{T})$
est ouvert. En appliquant alors (3) à l'application
$f^{-1}(\overline{T}) \to \overline{T}$, nous voyons que $T$ est ouvert
dans $\overline{T}$, c'est-à-dire que $T$ est localement fermé.
\end{proof}

\begin{lemma}
\label{lemma-closed-morphism-quotient-topology}
Soit $f : X \to Y$ une application surjective, fermée et continue d'espaces
topologiques. Soit $T \subset Y$ une partie. Alors
\begin{enumerate}
\item $\overline{T} = f(\overline{f^{-1}(T)})$,
\item $T \subset Y$ est fermé si et seulement si $f^{-1}(T)$ est fermé,
\item $T \subset Y$ est ouvert si et seulement si $f^{-1}(T)$ est ouvert, et
\item $T \subset Y$ est localement fermé si et seulement si
$f^{-1}(T)$ est localement fermé.
\end{enumerate}
En particulier, $f$ est submersive.
\end{lemma}

\begin{proof}
Il est clair que $\overline{f^{-1}(T)} \subset f^{-1}(\overline{T})$.
Alors $T \subset f(\overline{f^{-1}(T)}) \subset \overline{T}$ est une
partie fermée, d'où (1). Le point (2) est évident puisque $f$ est fermée et
surjective. Le point (3) résulte de (2) appliqué au complémentaire de $T$.
Pour (4), si $f^{-1}(T)$ est localement fermé, alors
$f^{-1}(T) \subset \overline{f^{-1}(T)}$ est ouvert.
Puisque l'application $\overline{f^{-1}(T)} \to \overline{T}$ est surjective
d'après (1), nous pouvons appliquer le point (3) à l'application
$\overline{f^{-1}(T)} \to \overline{T}$ induite par $f$ et conclure que $T$
est ouvert dans $\overline{T}$, c'est-à-dire que $T$ est localement fermé.
\end{proof}








\section{Composantes connexes}
\label{section-connected-components}

\begin{definition}
\label{definition-connected-components}
Soit $X$ un espace topologique.
\begin{enumerate}
\item On dit que $X$ est {\it connexe} si $X$ n'est pas vide et si, chaque
fois que $X = T_1 \amalg T_2$ avec $T_i \subset X$ ouvert et fermé, on a
$T_1 = \emptyset$ ou $T_2 = \emptyset$.
\item On dit que $T \subset X$ est une {\it composante connexe} de $X$ si
$T$ est une partie connexe maximale de $X$.
\end{enumerate}
\end{definition}

\noindent
L'espace vide n'est pas connexe.

\begin{lemma}
\label{lemma-image-connected-space}
Soit $f : X \to Y$ une application continue d'espaces topologiques.
Si $E \subset X$ est une partie connexe, alors $f(E) \subset Y$ est également
connexe.
\end{lemma}

\begin{proof}
Soit $A \subset f(E)$ une partie ouverte et fermée de $f(E)$. Comme $f$ est
continue, $f^{-1}(A)$ est une partie ouverte et fermée de $E$. Puisque $E$ est
connexe, $f^{-1}(A) = \emptyset$ ou $f^{-1}(A) = E$. Or $A\subset f(E)$
entraîne que $A = f(f^{-1}(A))$. En effet, si $x\in f(f^{-1}(A))$, il existe
$y\in f^{-1}(A)$ tel que $f(y) = x$ et, puisque $y\in f^{-1}(A)$, nous avons
$f(y) \in A$, c'est-à-dire $x\in A$. Réciproquement, si $x\in A$, l'inclusion
$A\subset f(E)$ implique qu'il existe $y\in E$ tel que $f(y) = x$. Alors
$y\in f^{-1}(A)$, puis $x\in f(f^{-1}(A))$. Ainsi, $A = \emptyset$ ou
$A = f(E)$, ce qui prouve que $f(E)$ est connexe.
\end{proof}

\begin{lemma}
\label{lemma-connected-components}
Soit $X$ un espace topologique.
\begin{enumerate}
\item Si $T \subset X$ est connexe, son adhérence l'est aussi.
\item Toute composante connexe de $X$ est fermée (mais pas nécessairement
ouverte).
\item Toute partie connexe de $X$ est contenue dans une unique composante
connexe de $X$.
\item Tout point de $X$ appartient à une unique composante connexe ; autrement
dit, $X$ est la réunion disjointe de ses composantes connexes.
\end{enumerate}
\end{lemma}

\begin{proof}
Soit $\overline{T}$ l'adhérence de la partie connexe $T$.
Supposons $\overline{T} = T_1 \amalg T_2$, où
$T_i \subset \overline{T}$ est ouvert et fermé. Alors
$T = (T\cap T_1) \amalg (T \cap T_2)$. Par conséquent, $T$ est égal à l'une
des deux parties, disons $T = T_1 \cap T$. Ainsi,
$\overline{T} \subset T_1$. Cela entraîne (1) et (2).

\medskip\noindent
Soit $A$ un ensemble non vide de parties connexes de $X$ tel que
$\Omega = \bigcap_{T \in A} T$ soit non vide. Nous affirmons que
$E = \bigcup_{T \in A} T$ est connexe. En effet, $E$ est non vide puisqu'il
contient $\Omega$. Écrivons $E = E_1 \amalg E_2$, où $E_i$ est fermé dans
$E$. Nous pouvons supposer que $E_1$ rencontre $\Omega$ (quitte à renuméroter).
Alors tout $T \in A$ rencontre $E_1$ et doit donc être contenu dans $E_1$,
puisque $T$ est connexe. Ainsi $E \subset E_1$, ce qui prouve l'affirmation.

\medskip\noindent
Soit $W \subset X$ une partie connexe non vide. Appliquons le résultat du
paragraphe précédent à l'ensemble de toutes les parties connexes de $X$ qui
contiennent $W$. Nous voyons alors que $E$ est une composante connexe de $X$.
Cela prouve l'existence et l'unicité dans (3).

\medskip\noindent
Soit $x \in X$. En prenant $W = \{x\}$ dans le paragraphe précédent, nous
voyons que $x$ appartient à une unique composante connexe de $X$.
Deux composantes connexes distinctes sont nécessairement disjointes, d'après
le résultat du deuxième paragraphe.

\medskip\noindent
Pour obtenir un exemple où les composantes connexes ne sont pas ouvertes, il
suffit de considérer un produit infini
$\prod_{n \in \mathbf{N}} \{0, 1\}$ muni de la topologie produit. Ses
composantes connexes sont les singletons, qui ne sont pas ouverts.
\end{proof}

\begin{remark}
\label{remark-quasi-components}
\begin{reference}
\cite[exemple 6.1.24]{Engelking}
\end{reference}
Soit $X$ un espace topologique et soit $x \in X$. Soit $Z \subset X$ la
composante connexe de $X$ passant par $x$. Considérons l'intersection $E$ de
toutes les parties ouvertes et fermées de $X$ qui contiennent $x$. Il est clair
que $Z \subset E$. En général, $Z \not = E$. Par exemple, soit
$X = \{x, y, z_1, z_2, \ldots\}$ muni de la topologie ayant pour base
d'ouverts $\{z_n\}$, $\{x, z_n, z_{n + 1}, \ldots\}$ et
$\{y, z_n, z_{n + 1}, \ldots\}$ pour tout $n$. Alors $Z = \{x\}$ et
$E = \{x, y\}$. Nous omettons les détails.
\end{remark}

\begin{lemma}
\label{lemma-connected-fibres-quotient-topology-connected-components}
Soit $f : X \to Y$ une application continue d'espaces topologiques.
Supposons que
\begin{enumerate}
\item toutes les fibres de $f$ soient connexes, et
\item une partie $T \subset Y$ soit fermée si et seulement si $f^{-1}(T)$
est fermée.
\end{enumerate}
Alors $f$ induit une bijection entre les ensembles de composantes connexes de
$X$ et de $Y$.
\end{lemma}

\begin{proof}
Soit $T \subset Y$ une composante connexe. Remarquons que $T$ est fermée ;
voir le lemme \ref{lemma-connected-components}. Le lemme résulte du fait que
$f^{-1}(T)$ est connexe, car toute partie connexe de $X$ s'envoie dans une
composante connexe de $Y$ d'après le lemme
\ref{lemma-image-connected-space}.
Supposons que $f^{-1}(T) = Z_1 \amalg Z_2$, avec $Z_1$, $Z_2$ fermés.
Pour tout $t \in T$, nous avons
$f^{-1}(\{t\}) = Z_1 \cap f^{-1}(\{t\}) \amalg Z_2 \cap f^{-1}(\{t\})$.
D'après (1), $f^{-1}(\{t\})$ est connexe ; nous en concluons que soit
$f^{-1}(\{t\}) \subset Z_1$, soit $f^{-1}(\{t\}) \subset Z_2$.
Autrement dit, $T = T_1 \amalg T_2$, avec $f^{-1}(T_i) = Z_i$.
D'après (2), $T_i$ est fermé dans $Y$. Ainsi, $T_1 = \emptyset$ ou
$T_2 = \emptyset$, comme voulu.
\end{proof}

\begin{lemma}
\label{lemma-connected-fibres-connected-components}
Soit $f : X \to Y$ une application continue d'espaces topologiques.
Supposons que
(a) $f$ soit ouverte,
(b) toutes les fibres de $f$ soient connexes.
Alors $f$ induit une bijection entre les ensembles de composantes connexes de
$X$ et de $Y$.
\end{lemma}

\begin{proof}
C'est un cas particulier du
lemme \ref{lemma-connected-fibres-quotient-topology-connected-components}.
\end{proof}

\begin{lemma}
\label{lemma-finite-fibre-connected-components}
Soit $f : X \to Y$ une application continue d'espaces topologiques non vides.
Supposons que
(a) $Y$ soit connexe,
(b) $f$ soit ouverte et fermée, et
(c) il existe un point $y\in Y$ tel que la fibre $f^{-1}(y)$ soit un ensemble
fini. Alors $X$ possède au plus $|f^{-1}(y)|$ composantes connexes. Par
conséquent, toute composante connexe $T$ de $X$ est ouverte et fermée, et
$f(T)$ est une partie ouverte et fermée non vide de $Y$, donc égale à $Y$.
\end{lemma}

\begin{proof}
Si l'espace topologique $X$ possède au moins $N$ composantes connexes pour un
certain $N \in \mathbf{N}$, nous obtenons par récurrence une décomposition
$X = X_1 \amalg \ldots \amalg X_N$ de $X$ en réunion disjointe de $N$ parties
ouvertes et fermées non vides $X_1, \ldots , X_N$ de $X$. Puisque $f$ est
ouverte et fermée, chaque $f(X_i)$ est une partie ouverte et fermée non vide de
$Y$, donc égale à $Y$. En particulier, l'intersection
$X_i \cap f^{-1}(y)$ est non vide pour tout $1 \leq i \leq N$. Ainsi,
$f^{-1}(y)$ possède au moins $N$ éléments.
\end{proof}

\begin{definition}
\label{definition-totally-disconnected}
Un espace topologique est {\it totalement discontinu} si toutes ses composantes
connexes sont des singletons.
\end{definition}

\noindent
Un espace discret est totalement discontinu.
Un espace totalement discontinu n'est pas nécessairement discret ; par
exemple, $\mathbf{Q} \subset \mathbf{R}$ est totalement discontinu mais non
discret.

\begin{lemma}
\label{lemma-space-connected-components}
Soit $X$ un espace topologique. Soit $\pi_0(X)$ l'ensemble des composantes
connexes de $X$. Soit $X \to \pi_0(X)$ l'application qui envoie $x \in X$
sur la composante connexe de $X$ passant par $x$. Munissons $\pi_0(X)$ de la
topologie quotient. Alors $\pi_0(X)$ est un espace totalement discontinu, et
toute application continue $X \to Y$ de $X$ vers un espace totalement
discontinu $Y$ se factorise par $\pi_0(X)$.
\end{lemma}

\begin{proof}
D'après le lemme
\ref{lemma-connected-fibres-quotient-topology-connected-components},
les composantes connexes de $\pi_0(X)$ sont les singletons.
Nous omettons la démonstration de la seconde assertion.
\end{proof}

\begin{definition}
\label{definition-locally-connected}
Un espace topologique $X$ est dit {\it localement connexe} si tout point
$x \in X$ possède un système fondamental de voisinages connexes.
\end{definition}

\begin{lemma}
\label{lemma-locally-connected}
Soit $X$ un espace topologique. Si $X$ est localement connexe, alors
\begin{enumerate}
\item toute partie ouverte de $X$ est localement connexe, et
\item les composantes connexes de $X$ sont ouvertes.
\end{enumerate}
Les composantes connexes des parties ouvertes de $X$ sont donc elles aussi
ouvertes. En particulier, tout point possède un système fondamental de
voisinages ouverts connexes.
\end{lemma}

\begin{proof}
Pour tout $x\in X$, notons $\mathcal{N}(x)$ le système fondamental de
voisinages connexes de $x$, et soit $U\subset X$ une partie ouverte de $X$.
Pour tout $x\in U$, $U$ est un voisinage de $x$, donc l'ensemble
$\{V \in \mathcal{N}(x) | V\subset U\}$ n'est pas vide et forme un système
fondamental de voisinages connexes de $x$ dans $U$. Ainsi, $U$ est localement
connexe, ce qui prouve (1).

\medskip\noindent
Soit $x \in \mathcal{C} \subset X$, où $\mathcal{C}$ est la composante
connexe de $x$. Comme $X$ est localement connexe, il existe un voisinage
connexe $\mathcal{N}$ de $x$. Par définition d'une composante connexe, nous
avons donc $\mathcal{N} \subset \mathcal{C}$ ; ainsi, $\mathcal{C}$ est un
voisinage de $x$. Cela montre que $\mathcal{C}$ est un voisinage de chacun de
ses points ; autrement dit, $\mathcal{C}$ est ouverte, ce qui prouve (2).
\end{proof}




\section{Composantes irréductibles}
\label{section-irreducible-components}

\begin{definition}
\label{definition-irreducible-components}
Soit $X$ un espace topologique.
\begin{enumerate}
\item On dit que $X$ est {\it irréductible} si $X$ n'est pas vide et si,
chaque fois que $X = Z_1 \cup Z_2$ avec $Z_i$ fermé, on a $X = Z_1$ ou
$X = Z_2$.
\item On dit que $Z \subset X$ est une {\it composante irréductible} de $X$
si $Z$ est une partie irréductible maximale de $X$.
\end{enumerate}
\end{definition}

\noindent
Un espace irréductible est évidemment connexe.

\begin{lemma}
\label{lemma-image-irreducible-space}
Soit $f : X \to Y$ une application continue d'espaces topologiques.
Si $E \subset X$ est une partie irréductible, alors $f(E) \subset Y$ est
également irréductible.
\end{lemma}

\begin{proof}
Nous pouvons évidemment supposer $E = X$ (c'est-à-dire $X$ irréductible) et
$f(E) = Y$ (c'est-à-dire $f$ surjective). Tout d'abord,
$Y \not = \emptyset$, puisque $X \not = \emptyset$. Supposons ensuite
$Y = Y_1 \cup Y_2$, où $Y_1$, $Y_2$ sont fermés. Alors
$X = X_1 \cup X_2$, où $X_i = f^{-1}(Y_i)$ est fermé dans $X$. Par hypothèse
sur $X$, on a nécessairement $X = X_1$ ou $X = X_2$, et donc $Y = Y_1$ ou
$Y = Y_2$, puisque $f$ est surjective.
\end{proof}

\begin{lemma}
\label{lemma-irreducible}
Soit $X$ un espace topologique.
\begin{enumerate}
\item Si $T \subset X$ est irréductible, son adhérence dans $X$ l'est aussi.
\item Toute composante irréductible de $X$ est fermée.
\item Toute partie irréductible de $X$ est contenue dans une composante
irréductible de $X$.
\item Tout point de $X$ appartient à une composante irréductible de $X$ ;
autrement dit, $X$ est la réunion de ses composantes irréductibles.
\end{enumerate}
\end{lemma}

\begin{proof}
Soit $\overline{T}$ l'adhérence de la partie irréductible $T$.
Si $\overline{T} = Z_1 \cup Z_2$, avec $Z_i \subset \overline{T}$ fermé,
alors $T = (T\cap Z_1) \cup (T \cap Z_2)$ ; ainsi, $T$ est égal à l'une des
deux parties, disons $T = Z_1 \cap T$. Il est donc clair que
$\overline{T} \subset Z_1$. Cela prouve (1). Le point (2) résulte aussitôt de
(1) et de la définition des composantes irréductibles.

\medskip\noindent
Soit $T \subset X$ irréductible. Considérons l'ensemble $A$ des parties
irréductibles $T \subset T' \subset X$. Remarquons que $A$ est non vide,
puisque $T \in A$. L'inclusion définit sur $A$ un ordre partiel :
$T' \leq T'' \Leftrightarrow T' \subset T''$.
Choisissons une partie totalement ordonnée maximale $A' \subset A$ et posons
$E = \bigcup_{T' \in A'} T'$. Nous affirmons que $E$ est irréductible.
Supposons en effet que $E =  Z_1 \cup Z_2$ soit la réunion de deux parties
fermées de $E$. Pour tout $T' \in A'$, on a soit $T' \subset Z_1$, soit
$T' \subset Z_2$, par l'irréductibilité de $T'$. Supposons qu'il existe
$T'_0 \in A'$ tel que $T'_0 \not\subset Z_1$ (sinon, nous avons déjà fini).
Puisque $A'$ est totalement ordonné, nous voyons immédiatement que
$T' \subset Z_2$ pour tout $T' \in A'$. Ainsi $E = Z_2$.
Cela prouve (3). Le point (4) résulte immédiatement de (3), puisqu'un espace
réduit à un singleton est irréductible.
\end{proof}

\begin{lemma}
\label{lemma-pick-irreducible-components}
Soit $X$ un espace topologique et supposons
$X = \bigcup_{i = 1, \ldots, n} X_i$, où chaque $X_i$ est une partie fermée
irréductible de $X$ et aucun $X_i$ n'est contenu dans la réunion des autres.
Alors chaque $X_i$ est une composante irréductible de $X$, et toute composante
irréductible de $X$ est l'un des $X_i$.
\end{lemma}

\begin{proof}
Soit $Y$ une composante irréductible de $X$. Écrivons
$Y = \bigcup_{i = 1, \ldots, n} (Y \cap X_i)$ et remarquons que chaque
$Y \cap X_i$ est fermé dans $Y$, puisque $X_i$ est fermé dans $X$.
L'irréductibilité de $Y$ montre que $Y$ est égal à l'un des $Y \cap X_i$,
c'est-à-dire que $Y \subset X_i$. Par maximalité des composantes
irréductibles, nous obtenons $Y = X_i$.

\medskip\noindent
Réciproquement, choisissons l'un des $X_i$. D'après le lemme
\ref{lemma-irreducible}, il est contenu dans une composante irréductible $Y$
de $X$. Le paragraphe précédent donne $Y = X_j$ pour un certain $j$.
Ainsi $X_i \subset X_j$ et, comme la réunion initiale ne contient aucun terme
superflu, $X_i = X_j$, qui est une composante irréductible.
\end{proof}

\begin{lemma}
\label{lemma-surjective-continuous-irreducible-components}
Soit $f : X \to Y$ une application continue et surjective d'espaces
topologiques. Si $X$ possède un nombre fini, disons $n$, de composantes
irréductibles, alors $Y$ possède un nombre de composantes irréductibles qui
est $\leq n$.
\end{lemma}

\begin{proof}
Soient $X_1, \ldots, X_n$ les composantes irréductibles de $X$.
D'après les lemmes \ref{lemma-image-irreducible-space} et
\ref{lemma-irreducible}, l'adhérence $Y_i \subset Y$ de $f(X_i)$ est
irréductible. Comme $f$ est surjective, $Y$ est la réunion des $Y_i$.
Nous pouvons choisir une partie minimale $I \subset \{1, \ldots, n\}$ telle
que $Y = \bigcup_{i \in I} Y_i$. Le lemme
\ref{lemma-pick-irreducible-components} montre alors que les $Y_i$, pour
$i \in I$, sont les composantes irréductibles de $Y$.
\end{proof}

\noindent
Un singleton est irréductible. Ainsi, si $x \in X$ est un point, l'adhérence
$\overline{\{x\}}$ est une partie fermée irréductible de $X$.

\begin{definition}
\label{definition-generic-point}
Soit $X$ un espace topologique.
\begin{enumerate}
\item Soit $Z \subset X$ une partie fermée irréductible. Un
{\it point générique} de $Z$ est un point $\xi \in Z$ tel que
$Z = \overline{\{\xi\}}$.
\item L'espace $X$ est dit {\it de Kolmogorov} si, pour tous
$x, x' \in X$ tels que $x \not = x'$, il existe une partie fermée de $X$ qui
contient exactement l'un des deux points.
\item L'espace $X$ est dit {\it quasi-sobre} si toute partie fermée
irréductible possède un point générique.
\item L'espace $X$ est dit {\it sobre} si toute partie fermée irréductible
possède un unique point générique.
\end{enumerate}
\end{definition}

\noindent
Un espace topologique $X$ est de Kolmogorov, quasi-sobre, resp.\ sobre si et
seulement si l'application $x\mapsto\overline{\{x\}}$ de $X$ vers l'ensemble
des parties fermées irréductibles de $X$ est injective, surjective, resp.\
bijective. Nous voyons donc qu'un espace topologique est sobre si et seulement
s'il est quasi-sobre et de Kolmogorov.

\begin{lemma}
\label{lemma-sober-subspace}
Soit $X$ un espace topologique et soit $Y\subset X$.
\begin{enumerate}
\item Si $X$ est de Kolmogorov, alors $Y$ l'est aussi.
\item Supposons que $Y$ soit localement fermé dans $X$. Si $X$ est
quasi-sobre, alors $Y$ l'est aussi.
\item Supposons que $Y$ soit localement fermé dans $X$. Si $X$ est sobre,
alors $Y$ l'est aussi.
\end{enumerate}
\end{lemma}

\begin{proof}
Démonstration de (1). Supposons que $X$ soit de Kolmogorov. Soient
$x,y\in Y$ avec $x\neq y$. Alors
$\overline{\overline{\{x\}}\cap Y}=\overline{\{x\}}\neq\overline{\{y\}}=
\overline{\overline{\{y\}}\cap Y}$. Par conséquent,
$\overline{\{x\}}\cap Y\neq\overline{\{y\}}\cap Y$. Cela montre que $Y$ est
de Kolmogorov.

\medskip\noindent
Démonstration de (2). Supposons que $X$ soit quasi-sobre. Il suffit de traiter
les cas où $Y$ est fermé et où $Y$ est ouvert. Supposons d'abord $Y$ fermé.
Soit $Z$ une partie fermée irréductible de $Y$. Alors $Z$ est une partie fermée
irréductible de $X$. Il existe donc $x \in Z$ tel que
$\overline{\{x\}} = Z$. Il s'ensuit que $\overline{\{x\}} \cap Y = Z$.
Ainsi, $Y$ est quasi-sobre. Supposons ensuite $Y$ ouvert. Soit $Z$ une partie
fermée irréductible de $Y$. Alors $\overline{Z}$ est une partie fermée
irréductible de $X$. Il existe donc $x \in \overline{Z}$ tel que
$\overline{\{x\}}=\overline{Z}$. Si $x\notin Y$, nous obtenons la
contradiction
$Z=Z\cap Y\subset\overline{Z}\cap Y=\overline{\{x\}}\cap Y=\emptyset$.
Par conséquent, $x\in Y$. Il s'ensuit que
$Z=\overline{Z}\cap Y=\overline{\{x\}}\cap Y$. Ainsi, $Y$ est quasi-sobre.

\medskip\noindent
Démonstration de (3). Elle résulte immédiatement de (1) et (2).
\end{proof}

\begin{lemma}
\label{lemma-sober-local}
Soit $X$ un espace topologique et soit $(X_i)_{i\in I}$ un recouvrement de
$X$.
\begin{enumerate}
\item Supposons $X_i$ localement fermé dans $X$ pour tout $i\in I$. Alors $X$
est de Kolmogorov si et seulement si $X_i$ l'est pour tout $i\in I$.
\item Supposons $X_i$ ouvert dans $X$ pour tout $i\in I$. Alors $X$ est
quasi-sobre si et seulement si $X_i$ l'est pour tout $i\in I$.
\item Supposons $X_i$ ouvert dans $X$ pour tout $i\in I$. Alors $X$ est sobre
si et seulement si $X_i$ l'est pour tout $i\in I$.
\end{enumerate}
\end{lemma}

\begin{proof}
Démonstration de (1). Si $X$ est de Kolmogorov, alors $X_i$ l'est pour tout
$i\in I$ d'après le lemme \ref{lemma-sober-subspace}. Supposons $X_i$ de
Kolmogorov pour tout $i\in I$. Soient $x,y\in X$ tels que
$\overline{\{x\}}=\overline{\{y\}}$. Il existe $i\in I$ tel que $x\in X_i$.
Il existe un ouvert $U\subset X$ tel que $X_i$ soit fermé dans $U$.
Si $y\notin U$, nous obtenons la contradiction
$x\in\overline{\{x\}}\cap U=\overline{\{y\}}\cap U=\emptyset$. Ainsi,
$y\in U$. Il s'ensuit que
$y\in\overline{\{y\}}\cap U=\overline{\{x\}}\cap U\subset X_i$.
Cela montre que $y\in X_i$. Par conséquent,
$\overline{\{x\}}\cap X_i=\overline{\{y\}}\cap X_i$. Comme $X_i$ est de
Kolmogorov, nous obtenons $x=y$. Cela montre que $X$ est de Kolmogorov.

\medskip\noindent
Démonstration de (2). Si $X$ est quasi-sobre, alors $X_i$ l'est pour tout
$i\in I$ d'après le lemme \ref{lemma-sober-subspace}. Supposons $X_i$
quasi-sobre pour tout $i\in I$. Soit $Y$ une partie fermée irréductible de
$X$. Comme $Y\neq\emptyset$, il existe $i\in I$ tel que
$X_i\cap Y\neq\emptyset$. Puisque $X_i$ est ouvert dans $X$, la partie
$X_i\cap Y$ est non vide et ouverte dans $Y$, donc irréductible et dense dans
$Y$. Ainsi, $X_i\cap Y$ est une partie fermée irréductible de $X_i$. Comme
$X_i$ est quasi-sobre, il existe $x\in X_i\cap Y$ tel que
$X_i\cap Y=\overline{\{x\}}\cap X_i\subset\overline{\{x\}}$. Puisque
$X_i\cap Y$ est dense dans $Y$ et que $Y$ est fermé dans $X$, il s'ensuit que
$Y=\overline{X_i\cap Y}\cap Y\subset\overline{X_i\cap Y}\subset
\overline{\{x\}}\subset Y$. Par conséquent,
$Y=\overline{\{x\}}$. Cela montre que $X$ est quasi-sobre.

\medskip\noindent
Démonstration de (3). Elle résulte immédiatement de (1) et (2).
\end{proof}

\begin{example}
\label{example-quasi-sober-not-kolmogorov}
Soit $X$ un espace indiscret de cardinalité au moins $2$. Alors $X$ est
quasi-sobre mais non de Kolmogorov. De plus, la famille de ses singletons est
un recouvrement de $X$ par des espaces discrets, donc de Kolmogorov.
\end{example}

\begin{example}
\label{example-kolmogorov-not-quasi-sober}
Soit $Y$ un ensemble infini, muni de la topologie dont les fermés sont $Y$ et
les parties finies de $Y$. Alors $Y$ est de Kolmogorov mais non quasi-sobre.
Cependant, la famille de ses singletons est un recouvrement par des espaces
discrets, donc sobres.
\end{example}

\begin{example}
\label{example-not-kolmogorov-not-quasi-sober}
Soient $X$ et $Y$ comme dans l'exemple
\ref{example-quasi-sober-not-kolmogorov} et l'exemple
\ref{example-kolmogorov-not-quasi-sober}. Alors $X\amalg Y$ n'est ni de
Kolmogorov ni quasi-sobre.
\end{example}

\begin{example}
\label{example-sober-subspace}
Soit $Z$ un ensemble infini et soit $z\in Z$. Munissons $Z$ de la topologie
dont les fermés sont $Z$ et les parties finies de $Z\setminus\{z\}$. Alors $Z$
est sobre, mais son sous-espace $Z\setminus\{z\}$ n'est pas quasi-sobre.
\end{example}

\begin{example}
\label{example-Hausdorff}
Rappelons qu'un espace topologique $X$ est séparé si et seulement si, pour
toute paire de points distincts $x, y \in X$, il existe des ouverts disjoints
$U, V \subset X$ tels que $x \in U$, $y \in V$. Dans ce cas, $X$ est
irréductible si et seulement si $X$ est réduit à un singleton. De même, toute
partie de $X$ est irréductible si et seulement si elle est réduite à un
singleton. Ainsi, tout espace séparé est sobre.
\end{example}

\begin{lemma}
\label{lemma-irreducible-on-top}
Soit $f : X \to Y$ une application continue d'espaces topologiques.
Supposons que
(a) $Y$ soit irréductible,
(b) $f$ soit ouverte, et
(c) il existe une famille dense de points $y \in Y$ tels que
$f^{-1}(y)$ soit irréductible.
Alors $X$ est irréductible.
\end{lemma}

\begin{proof}
Supposons $X = Z_1 \cup Z_2$, avec $Z_i$ fermé.
Considérons les ouverts $U_1 = Z_1 \setminus Z_2 = X \setminus Z_2$ et
$U_2 = Z_2 \setminus Z_1 = X \setminus Z_1$. Pour obtenir une contradiction,
supposons que $U_1$ et $U_2$ soient tous deux non vides. D'après (b),
$f(U_i)$ est ouvert. D'après (a), $Y$ est irréductible, et donc
$f(U_1) \cap f(U_2) \not = \emptyset$. D'après (c), cet ouvert non vide
contient un point $y$ tel que la fibre
$X_y = f^{-1}(y)$ soit irréductible. Alors $X_y \cap U_1$ et
$X_y \cap U_2$ sont des parties ouvertes non vides et disjointes de $X_y$,
ce qui est contradictoire.
\end{proof}

\begin{lemma}
\label{lemma-irreducible-fibres-irreducible-components}
Soit $f : X \to Y$ une application continue d'espaces topologiques.
Supposons que (a) $f$ soit ouverte et que
(b) pour tout $y \in Y$, la fibre $f^{-1}(y)$ soit irréductible.
Alors $f$ induit une bijection entre les composantes irréductibles.
\end{lemma}

\begin{proof}
Remarquons que l'hypothèse (b) implique que $f$ est surjective (voir la
définition \ref{definition-irreducible-components}).
Soit $T \subset Y$ une composante irréductible.
Remarquons que $T$ est fermée ; voir le lemme \ref{lemma-irreducible}.
Le lemme résulte du fait que $f^{-1}(T)$ est irréductible, car toute partie
irréductible de $X$ s'envoie dans une composante irréductible de $Y$ d'après
le lemme \ref{lemma-image-irreducible-space}.
Remarquons que $f^{-1}(T) \to T$ satisfait aux hypothèses du
lemme \ref{lemma-irreducible-on-top}. Nous avons donc terminé.
\end{proof}

\noindent
La construction du lemme suivant est parfois appelée la « sobrification ».

\begin{lemma}
\label{lemma-make-sober}
\begin{reference}
\cite[p. 68 et suiv.]{EGA1-second}
\end{reference}
Soit $X$ un espace topologique. Il existe une application continue canonique
$$
c : X \longrightarrow X'
$$
de $X$ vers un espace topologique sobre $X'$, universelle parmi les
applications continues de $X$ vers les espaces topologiques sobres.
De plus, l'application $U' \mapsto c^{-1}(U')$ est une bijection entre les
ouverts de $X'$ et ceux de $X$, qui commute aux intersections finies et aux
réunions quelconques. L'image $c(X)$ est un espace topologique de Kolmogorov,
et l'application $c : X \to c(X)$ est universelle parmi les applications de
$X$ vers les espaces de Kolmogorov.
\end{lemma}

\begin{proof}
Soit $X'$ l'ensemble des parties fermées irréductibles de $X$ et soit
$$
c : X \to X', \quad x \mapsto \overline{\{x\}}.
$$
Pour tout ouvert $U \subset X$, notons $U' \subset X'$ l'ensemble des
parties fermées irréductibles de $X$ qui rencontrent $U$.
Alors $c^{-1}(U') = U$. En particulier, si $U_1 \not = U_2$ sont ouverts dans
$X$, alors $U'_1 \not = U_2'$. Ainsi, $c$ induit une bijection entre les
parties de $X'$ de la forme $U'$ et les ouverts de $X$.

\medskip\noindent
Soient $U_1, U_2$ des ouverts de $X$. Supposons que $Z \in U'_1$ et
$Z \in U'_2$. Alors $Z \cap U_1$ et $Z \cap U_2$ sont des parties ouvertes
non vides de l'espace irréductible $Z$ ; par conséquent,
$Z \cap U_1 \cap U_2$ est non vide. Ainsi,
$(U_1 \cap U_2)' = U'_1 \cap U'_2$.
La règle $U \mapsto U'$ est également compatible avec les réunions
quelconques (nous omettons les détails). Il est donc clair que la famille des
$U'$ forme une topologie sur $X'$ et que l'on obtient la bijection annoncée
dans le lemme.

\medskip\noindent
Montrons ensuite que $X'$ est sobre. Soit $T \subset X'$ une partie fermée
irréductible. Soit $U \subset X$ l'ouvert tel que $X' \setminus T = U'$.
Alors $Z = X \setminus U$ est irréductible, en vertu des propriétés de la
bijection du lemme. Nous affirmons que $Z \in T$ est l'unique point générique.
En effet, tout ouvert de la forme $V' \subset X'$ qui ne contient pas $Z$
provient nécessairement d'un ouvert $V \subset X$ qui ne rencontre pas $Z$,
c'est-à-dire qui est contenu dans $U$.

\medskip\noindent
Vérifions enfin la propriété universelle. Soit $f : X \to Y$ une application
continue à valeurs dans un espace topologique sobre. Définissons alors
$f' : X' \to Y$ comme l'application qui envoie la partie fermée irréductible
$Z \subset X$ sur l'unique point générique de $\overline{f(Z)}$.
Il est immédiat que $f' \circ c = f$ comme applications d'ensembles, et les
propriétés de $c$ entraînent que $f'$ est continue. Nous omettons la
vérification de l'unicité de l'application continue $f'$. Nous omettons aussi
la démonstration des assertions concernant les espaces de Kolmogorov.
\end{proof}

\begin{lemma}
\label{lemma-remove-irreducible-connected}
Soit $X$ un espace topologique connexe possédant un nombre fini de composantes
irréductibles $X_1, \ldots, X_n$. Si $n > 1$, il existe
$1 \leq j \leq n$ tel que $X' = \bigcup_{i \not = j} X_i$ soit connexe.
\end{lemma}

\begin{proof}
C'est un problème de théorie des graphes. Soit $\Gamma$ le graphe dont les
sommets sont $V = \{1, \ldots, n\}$ et dans lequel $i$ et $j$ sont reliés par
une arête si et seulement si $X_i \cap X_j$ est non vide. La connexité de $X$
signifie que $\Gamma$ est connexe. Il s'agit de trouver
$1 \leq j \leq n$ tel que $\Gamma \setminus \{j\}$ soit encore connexe.
On peut le faire en choisissant $j, j' \in E$ à distance maximale ; alors $j$
convient (choisir une feuille !). Nous omettons les détails.
\end{proof}

\section{Espaces topologiques noethériens}
\label{section-noetherian}

\begin{definition}
\label{definition-noetherian}
Un espace topologique est dit {\it noethérien} si les parties fermées de $X$
satisfont à la condition de chaîne descendante. Un espace topologique est dit
{\it localement noethérien} si tout point possède un voisinage noethérien.
\end{definition}

\begin{lemma}
\label{lemma-Noetherian}
Soit $X$ un espace topologique noethérien.
\begin{enumerate}
\item Toute partie de $X$, munie de la topologie induite, est noethérienne.
\item L'espace $X$ possède un nombre fini de composantes irréductibles.
\item Chaque composante irréductible de $X$ contient un ouvert non vide de
$X$.
\end{enumerate}
\end{lemma}

\begin{proof}
Soit $T \subset X$ une partie de $X$.
Soit $T_1 \supset T_2 \supset \ldots$ une chaîne descendante de parties
fermées de $T$. Écrivons $T_i =  T \cap Z_i$, avec $Z_i \subset X$ fermé.
Considérons la chaîne descendante de parties fermées
$Z_1 \supset Z_1\cap Z_2 \supset Z_1 \cap Z_2 \cap Z_3 \ldots$
Elle est stationnaire par hypothèse, et la suite initiale des $T_i$ l'est donc
elle aussi. Ainsi $T$ est noethérien.

\medskip\noindent
Soit $A$ l'ensemble des parties fermées de $X$ qui n'ont pas un nombre fini de
composantes irréductibles. Supposons que $A$ soit non vide afin d'aboutir à une
contradiction. L'inclusion définit sur $A$ un ordre partiel :
$\alpha \leq \alpha' \Leftrightarrow Z_{\alpha} \subset Z_{\alpha'}$.
La condition de chaîne descendante permet de trouver un plus petit élément de
$A$, disons $Z$. Comme $Z$ n'est pas une réunion finie de composantes
irréductibles, il n'est pas irréductible. Nous pouvons donc écrire
$Z = Z' \cup Z''$, où les deux termes sont des parties fermées strictement plus
petites. Par construction, $Z' = \bigcup Z'_i$ et
$Z'' = \bigcup Z''_j$ sont des réunions finies de leurs composantes
irréductibles (lemme \ref{lemma-irreducible}). Ainsi,
$Z = \bigcup Z'_i \cup \bigcup Z''_j$ est une réunion finie de parties fermées
irréductibles. Après suppression des termes superflus de cette expression, on
obtient la décomposition de $Z$ en ses composantes irréductibles (lemme
\ref{lemma-pick-irreducible-components}), ce qui est contradictoire.

\medskip\noindent
Soit $Z \subset X$ une composante irréductible de $X$.
Soient $Z_1, \ldots, Z_n$ les autres composantes irréductibles de $X$.
Considérons $U = Z \setminus (Z_1\cup\ldots\cup Z_n)$.
Cette partie n'est pas vide, sinon l'espace irréductible $Z$ serait contenu
dans l'un des autres $Z_i$. Comme
$X = Z \cup Z_1 \cup \ldots Z_n$ (voir le lemme \ref{lemma-irreducible}), on a
également $U = X \setminus (Z_1\cup\ldots\cup Z_n)$, qui est donc ouvert dans
$X$. Ainsi, $Z$ contient un ouvert non vide de $X$.
\end{proof}

\begin{lemma}
\label{lemma-image-Noetherian}
Soit $f : X \to Y$ une application continue d'espaces topologiques.
\begin{enumerate}
\item Si $X$ est noethérien, alors $f(X)$ est noethérien.
\item Si $X$ est localement noethérien et si $f$ est ouverte, alors $f(X)$ est
localement noethérien.
\end{enumerate}
\end{lemma}

\begin{proof}
Dans le cas (1), supposons que
$Z_1 \supset Z_2 \supset Z_3 \supset \ldots$ soit une chaîne descendante de
parties fermées de $f(X)$ (muni, comme d'habitude, de la topologie induite en
tant que partie de $Y$). Alors
$f^{-1}(Z_1) \supset f^{-1}(Z_2) \supset f^{-1}(Z_3) \supset \ldots$ est une
chaîne descendante de parties fermées de $X$. Cette chaîne est donc
stationnaire. Puisque $f(f^{-1}(Z_i)) = Z_i$, nous en concluons que
$Z_1 \supset Z_2 \supset Z_3 \supset \ldots$ est elle aussi stationnaire.
Dans le cas (2), soit $y \in f(X)$. Choisissons $x \in X$ tel que
$f(x) = y$. Par hypothèse, il existe un voisinage $E \subset X$ de $x$ qui est
noethérien. Alors $f(E) \subset f(X)$ est un voisinage noethérien d'après le
point (1).
\end{proof}

\begin{lemma}
\label{lemma-finite-union-Noetherian}
Soit $X$ un espace topologique.
Soient $X_i \subset X$, $i = 1, \ldots, n$, une famille finie de parties.
Si chaque $X_i$ est noethérien (pour la topologie induite), alors
$\bigcup_{i = 1, \ldots, n}  X_i$ est noethérien (pour la topologie induite).
\end{lemma}

\begin{proof}
Soit $\{F_m\}_{m \in \mathbf{N}}$ une suite décroissante de parties fermées de
$X' = \bigcup_{i = 1, \ldots, n}  X_i$ muni de la topologie induite.
On peut trouver une suite décroissante $\{G_m\}_{m \in \mathbf{N}}$ de parties
fermées de $X$ vérifiant $F_m = G_m \cap X'$ pour tout $m$ (nous omettons ce
petit détail). Comme $X_i$ est noethérien et que
$\{G_m \cap X_i\}_{m \in \mathbf{N}}$ est une suite décroissante de parties
fermées de $X_i$, il existe $m_i \in \mathbf{N}$ tel que, pour tout
$m \geq m_i$, on ait $G_m \cap X_i = G_{m_i} \cap X_i$.
Posons $m_0 = \max_{i = 1, \ldots, n} m_i$. Alors, manifestement,
$$
F_m = G_m \cap X' = G_m \cap (X_1 \cup \ldots \cup X_n) =
(G_m \cap X_1) \cup \ldots (G_m \cap X_n)
$$
est stationnaire pour $m \geq m_0$, ce qui achève la démonstration.
\end{proof}

\begin{example}
\label{example-locally-Noetherian-no-closed-point}
Tout espace topologique noethérien non vide et de Kolmogorov possède un point
fermé (combiner les lemmes \ref{lemma-quasi-compact-closed-point} et
\ref{lemma-Noetherian-quasi-compact}).
Soit $X = \{1, 2, 3, \ldots \}$. Définissons sur $X$ la topologie dont les
ouverts sont $\emptyset$, $\{1, 2, \ldots, n\}$, $n \geq 1$, et $X$.
Alors $X$ est un espace topologique localement noethérien sans aucun point
fermé. Cet espace ne peut pas être l'espace topologique sous-jacent à un schéma
localement noethérien ; voir Propriétés, lemme
\ref{properties-lemma-locally-Noetherian-closed-point}.
\end{example}

\begin{lemma}
\label{lemma-locally-Noetherian-locally-connected}
Soit $X$ un espace topologique localement noethérien.
Alors $X$ est localement connexe.
\end{lemma}

\begin{proof}
Soit $x \in X$. Soit $E$ un voisinage de $x$.
Nous devons trouver un voisinage connexe de $x$ contenu dans $E$.
Par hypothèse, il existe un voisinage $E'$ de $x$ qui est noethérien.
Alors $E \cap E'$ est noethérien ; voir le lemme \ref{lemma-Noetherian}.
Soit $E \cap E' = Y_1 \cup \ldots \cup Y_n$ sa décomposition en composantes
irréductibles ; voir le lemme \ref{lemma-Noetherian}.
Soit $E'' = \bigcup_{x \in Y_i} Y_i$. C'est une partie connexe de
$E \cap E'$ qui contient $x$. Elle contient l'ouvert
$E \cap E' \setminus (\bigcup_{x \not \in Y_i} Y_i)$ de $E \cap E'$ ; c'est
donc un voisinage de $x$ dans $X$. Cela prouve le lemme.
\end{proof}




\section{Dimension de Krull}
\label{section-krull-dimension}

\begin{definition}
\label{definition-Krull}
Soit $X$ un espace topologique.
\begin{enumerate}
\item Une {\it chaîne de parties fermées irréductibles} de $X$ est une suite
$Z_0 \subset Z_1 \subset \ldots \subset Z_n \subset X$, où $Z_i$ est fermé
irréductible et $Z_i \not = Z_{i + 1}$ pour $i = 0, \ldots, n - 1$.
\item La {\it longueur} d'une chaîne
$Z_0 \subset Z_1 \subset \ldots \subset Z_n \subset X$ de parties fermées
irréductibles de $X$ est l'entier $n$.
\item La {\it dimension}, ou plus précisément la {\it dimension de Krull},
$\dim(X)$ de $X$, est l'élément de
$\{-\infty, 0, 1, 2, 3, \ldots, \infty\}$ défini par la formule
$$
\dim(X) =
\sup \{\text{longueurs des chaînes de parties fermées irréductibles}\}
$$
Ainsi, $\dim(X) = -\infty$ si et seulement si $X$ est l'espace vide.
\item Soit $x \in X$. La {\it dimension de Krull de $X$ en $x$} est définie
par
$$
\dim_x(X) = \min \{\dim(U), x\in U\subset X\text{ ouvert}\}
$$
c'est-à-dire comme le minimum des $\dim(U)$ lorsque $U$ parcourt les
voisinages ouverts de $x$ dans $X$.
\end{enumerate}
\end{definition}

\noindent
Remarquons que, si $U' \subset U \subset X$ sont ouverts, alors
$\dim(U') \leq \dim(U)$. Ainsi, si $\dim_x(X) = d$, le point $x$ possède un
système fondamental de voisinages ouverts $U$ tels que
$\dim(U) = \dim_x(X)$.

\begin{lemma}
\label{lemma-dimension-supremum-local-dimensions}
Soit $X$ un espace topologique. Alors $\dim(X) = \sup \dim_x(X)$, où la borne
supérieure est prise sur les points $x$ de $X$.
\end{lemma}

\begin{proof}
Il est clair que $\dim(X) \geq \dim_x(X)$ pour tout $x \in X$ (voir la
discussion qui suit la définition \ref{definition-Krull}). Nous obtenons ainsi
l'une des inégalités. Pour la réciproque, soit $n \geq 0$ et supposons que
$\dim(X) \geq n$. On peut alors trouver une chaîne de parties fermées
irréductibles $Z_0 \subset Z_1 \subset \ldots \subset Z_n \subset X$.
Choisissons $x \in Z_0$. Pour tout voisinage ouvert $U$ de $x$, nous obtenons
une chaîne de parties fermées irréductibles
$$
Z_0 \cap U \subset Z_1 \cap U \subset \ldots \subset Z_n \cap U
$$
dans $U$. En effet, les ensembles $U \cap Z_i$ sont fermés irréductibles dans
$U$ et
les inclusions sont strictes (détails omis ; indication : l'adhérence de
$U \cap Z_i$ est $Z_i$).
Nous voyons ainsi que $\dim_x(X) \geq n$, ce qui prouve l'autre inégalité.
\end{proof}

\begin{example}
\label{example-Krull-Rn}
La dimension de Krull de l'espace euclidien usuel $\mathbf{R}^n$ est $0$.
\end{example}

\begin{example}
\label{example-krull-2set}
Soit $X = \{s, \eta\}$, dont les ouverts sont
$\{\emptyset, \{\eta\}, \{s, \eta\}\}$.
Dans ce cas, une chaîne maximale de parties fermées irréductibles est
$\{s\} \subset \{s, \eta\}$. Ainsi, $\dim(X) = 1$. Il est facile de
généraliser cet exemple pour obtenir un espace topologique à $(n + 1)$
éléments et de dimension de Krull $n$.
\end{example}

\begin{definition}
\label{definition-equidimensional}
Soit $X$ un espace topologique.
On dit que $X$ est {\it équidimensionnel} si toutes les composantes
irréductibles de $X$ ont la même dimension.
\end{definition}





\section{Codimension et espaces caténaires}
\label{section-catenary-spaces}

\noindent
Nous ne définissons que la codimension des parties fermées irréductibles.

\begin{definition}
\label{definition-codimension}
Soit $X$ un espace topologique.
Soit $Y \subset X$ une partie fermée irréductible.
La {\it codimension} de $Y$ dans $X$ est la borne supérieure des longueurs
$e$ des chaînes
$$
Y = Y_0 \subset Y_1 \subset \ldots \subset Y_e \subset X
$$
de parties fermées irréductibles de $X$ qui commencent par $Y$.
Nous la noterons $\text{codim}(Y, X)$.
\end{definition}

\noindent
La codimension est un élément de $\{0, 1, 2, \ldots\} \cup \{\infty\}$.
Si $\text{codim}(Y, X) < \infty$, toute chaîne peut être prolongée en une
chaîne maximale (mais celles-ci n'ont pas nécessairement toutes la même
longueur).

\begin{lemma}
\label{lemma-codimension-at-generic-point}
Soit $X$ un espace topologique.
Soit $Y \subset X$ une partie fermée irréductible.
Soit $U \subset X$ un ouvert tel que $Y \cap U$ soit non vide. Alors
$$
\text{codim}(Y, X) = \text{codim}(Y \cap U, U)
$$
\end{lemma}

\begin{proof}
La règle $T \mapsto \overline{T}$ définit une application bijective qui
préserve les inclusions entre les parties fermées irréductibles de $U$ et les
parties fermées irréductibles de $X$ qui rencontrent $U$.
Le lemme en résulte facilement. Nous omettons les détails.
\end{proof}

\begin{example}
\label{example-Noetherian-infinite-codimension}
Soit $X = [0, 1]$ l'intervalle unité muni de la topologie suivante : les
ensembles $[0, 1]$, $(1 - 1/n, 1]$ pour $n \in \mathbf{N}$, et $\emptyset$
sont ouverts. Les fermés sont donc $\emptyset$, $\{0\}$,
$[0, 1 - 1/n]$ pour $n > 1$, et $[0, 1]$.
C'est manifestement un espace topologique noethérien. Mais la partie fermée
irréductible $Y = \{0\}$ est de codimension infinie :
$\text{codim}(Y, X) = \infty$.
Pour le voir, il suffit de remarquer que tous les fermés $[0, 1 - 1/n]$ sont
irréductibles.
\end{example}

\begin{definition}
\label{definition-catenary}
Soit $X$ un espace topologique. On dit que $X$ est {\it caténaire} si, pour
toute paire de parties fermées irréductibles $T \subset T'$, on a
$\text{codim}(T, T') < \infty$ et si toute chaîne maximale de parties fermées
irréductibles
$$
T = T_0 \subset T_1 \subset \ldots \subset T_e = T'
$$
a la même longueur (égale à la codimension).
\end{definition}

\begin{lemma}
\label{lemma-catenary}
Soit $X$ un espace topologique.
Les assertions suivantes sont équivalentes :
\begin{enumerate}
\item $X$ est caténaire,
\item $X$ possède un recouvrement ouvert par des espaces caténaires.
\end{enumerate}
De plus, dans ce cas, tout sous-espace localement fermé de $X$ est caténaire.
\end{lemma}

\begin{proof}
Supposons que $X$ soit caténaire et que $U \subset X$ soit ouvert.
La règle $T \mapsto \overline{T}$ définit une application bijective qui
préserve les inclusions entre les parties fermées irréductibles de $U$ et les
parties fermées irréductibles de $X$ qui rencontrent $U$.
Le lemme en résulte facilement. Nous omettons les détails.
\end{proof}

\begin{lemma}
\label{lemma-catenary-in-codimension}
Soit $X$ un espace topologique. Les assertions suivantes sont équivalentes :
\begin{enumerate}
\item $X$ est caténaire, et
\item pour toute paire de parties fermées irréductibles $Y \subset Y'$, on a
$\text{codim}(Y, Y') < \infty$, et pour tout triplet
$Y \subset Y' \subset Y''$ de parties fermées irréductibles, on a
$$
\text{codim}(Y, Y'') = \text{codim}(Y, Y') + \text{codim}(Y', Y'').
$$
\end{enumerate}
\end{lemma}

\begin{proof}
Supposons que $X$ soit caténaire. D'après la définition
\ref{definition-catenary}, pour toute paire de parties fermées irréductibles
$Y \subset Y'$, on a $\text{codim}(Y,Y') < \infty$.
Soit $Y \subset Y' \subset Y''$ un triplet de parties fermées irréductibles de
$X$. Soit
$$
Y = Y_0 \subset Y_1 \subset ... \subset Y_{e_1} = Y'
$$
une chaîne maximale de parties fermées irréductibles entre $Y$ et $Y'$, où
$e_1 = \text{codim}(Y,Y')$. Soit aussi
$$
Y' = Y_{e_1} \subset Y_{e_1 + 1}\subset ... \subset Y_{e_1 + e_2} = Y''
$$
une chaîne maximale de parties fermées irréductibles entre $Y'$ et $Y''$, où
$e_2 = \text{codim}(Y',Y'')$.
Comme les deux chaînes sont maximales, leur concaténation
$$
Y = Y_0\subset Y_1 \subset ... \subset Y_{e_1} = Y' =
Y_{e_1} \subset Y_{e_1+1}\subset ... \subset Y_{e_1+e_2}=Y''
$$
est elle aussi maximale (entre $Y$ et $Y''$), et sa longueur est
$e_1 + e_2$. Puisque $X$ est caténaire, toute chaîne maximale a la même
longueur, égale à la codimension. Le point (2), à savoir
$\text{codim}(Y,Y'') = e_1 + e_2 = \text{codim}(Y,Y') + \text{codim}(Y',Y'')$,
est donc vérifié.

\medskip\noindent
Réciproquement, une récurrence montre que, si
$Y = Y_1 \subset ... \subset Y_n = Y'$, alors
$ \text{codim}(Y,Y') = \text{codim}(Y_1,Y_2) + ... +
\text{codim}(Y_{n-1},Y_n)$. Cela force les chaînes maximales à avoir la même
longueur.
\end{proof}








\section{Espaces et applications quasi-compacts}
\label{section-quasi-compact}

\noindent
Nous réserverons le mot « compact » aux espaces topologiques séparés. Or de
nombreux espaces qui interviennent en géométrie algébrique ne sont pas séparés.

\begin{definition}
\label{definition-quasi-compact}
Quasi-compacité.
\begin{enumerate}
\item On dit qu'un espace topologique $X$ est {\it quasi-compact} si tout
recouvrement ouvert de $X$ admet un sous-recouvrement fini.
\item On dit qu'une application continue $f : X \to Y$ est
{\it quasi-compacte} si l'image réciproque $f^{-1}(V)$ de tout ouvert
quasi-compact $V \subset Y$ est quasi-compacte.
\item On dit qu'une partie $Z \subset X$ est {\it rétrocompacte} si
l'application d'inclusion $Z \to X$ est quasi-compacte.
\end{enumerate}
\end{definition}

\noindent
Dans de nombreux textes de topologie, un espace est dit {\it compact} s'il est
quasi-compact et séparé ; dans d'autres, la condition de séparation est omise.
Pour éviter toute ambiguïté en géométrie algébrique, nous employons le terme
quasi-compact. La quasi-compacité d'une application est très différente de la
notion d'« application propre » : pour celle-ci, on exige (outre le caractère
fermé et séparé) que l'image réciproque de toute partie quasi-compacte du but
soit quasi-compacte, tandis que la définition ci-dessus ne considère que les
ensembles {\it ouverts} quasi-compacts.

\begin{lemma}
\label{lemma-composition-quasi-compact}
Une composée d'applications quasi-compactes est quasi-compacte.
\end{lemma}

\begin{proof}
C'est immédiat d'après la définition.
\end{proof}

\begin{lemma}
\label{lemma-closed-in-quasi-compact}
Une partie fermée d'un espace topologique quasi-compact est quasi-compacte.
\end{lemma}

\begin{proof}
Soit $E \subset X$ une partie fermée de l'espace quasi-compact $X$.
Soit $E = \bigcup V_j$ un recouvrement ouvert. Choisissons des ouverts
$U_j \subset X$ tels que $V_j = E \cap U_j$. Alors
$X = (X \setminus E) \cup \bigcup U_j$ est un recouvrement ouvert de $X$.
On a donc $X = (X \setminus E) \cup U_{j_1} \cup \ldots \cup U_{j_n}$ pour
un certain $n$ et des indices $j_i$.
Ainsi $E = V_{j_1} \cup \ldots \cup V_{j_n}$, comme voulu.
\end{proof}

\begin{lemma}
\label{lemma-quasi-compact-in-Hausdorff}
Soit $X$ un espace topologique séparé.
\begin{enumerate}
\item Si $E \subset X$ est quasi-compact, alors il est fermé.
\item Si $E_1, E_2 \subset X$ sont deux parties quasi-compactes disjointes,
il existe des ouverts $E_i \subset U_i$ tels que
$U_1 \cap U_2 = \emptyset$.
\end{enumerate}
\end{lemma}

\begin{proof}
Démonstration de (1). Soit $x \in X$, avec $x \not \in E$.
Pour tout $e \in E$, on peut trouver des ouverts disjoints $V_e$ et $U_e$
tels que $e \in V_e$ et $x \in U_e$. Comme $E \subset \bigcup V_e$, il
existe un nombre fini d'éléments $e_1, \ldots, e_n$ tels que
$E \subset V_{e_1} \cup \ldots \cup V_{e_n}$. Alors
$U = U_{e_1} \cap \ldots \cap U_{e_n}$ est un voisinage ouvert de $x$ qui ne
rencontre pas $V_{e_1} \cup \ldots \cup V_{e_n}$. En particulier, il ne
rencontre pas $E$. Ainsi $E$ est fermé.

\medskip\noindent
Démonstration de (2). Dans la démonstration de (1), nous avons vu que, pour
$x \in E_1$, on peut trouver un voisinage ouvert $x \in U_x$ et un ouvert
$E_2 \subset V_x$ tels que $U_x \cap V_x = \emptyset$. Comme $E_1$ est
quasi-compact, il existe un nombre fini de points $x_i \in E_1$ tels que
$E_1 \subset U = U_{x_1} \cup \ldots \cup U_{x_n}$.
Il suffit de prendre $V = V_{x_1} \cap \ldots \cap V_{x_n}$.
\end{proof}

\begin{lemma}
\label{lemma-closed-in-compact}
Soit $X$ un espace séparé quasi-compact. Soit $E \subset X$.
Les assertions suivantes sont équivalentes : (a) $E$ est fermé dans $X$,
(b) $E$ est quasi-compact.
\end{lemma}

\begin{proof}
L'implication (a) $\Rightarrow$ (b) est le
lemme \ref{lemma-closed-in-quasi-compact}.
L'implication (b) $\Rightarrow$ (a) est le
lemme \ref{lemma-quasi-compact-in-Hausdorff}.
\end{proof}

\noindent
L'énoncé suivant n'est qu'une reformulation de la quasi-compacité.

\begin{lemma}
\label{lemma-intersection-closed-in-quasi-compact}
Soit $X$ un espace topologique quasi-compact.
Si $\{Z_\alpha\}_{\alpha \in A}$ est une famille de parties fermées telle que
l'intersection de toute sous-famille finie soit non vide, alors
$\bigcap_{\alpha \in A} Z_\alpha$ est non vide.
\end{lemma}

\begin{proof}
Supposons que $\bigcap_{\alpha \in A} Z_{\alpha} = \emptyset$.
En passant aux complémentaires, on a
$\bigcup_{\alpha \in A} (X\setminus Z_{\alpha})=X$.
Comme les parties $Z_\alpha$ sont fermées,
$\bigcup_{\alpha \in A} (X \setminus Z_{\alpha})$ est un recouvrement ouvert
de l'espace quasi-compact $X$. Il existe donc une partie finie $J \subset A$
telle que $X = \bigcup_{\alpha \in J} (X\setminus Z_{\alpha})$.
Son complémentaire est alors vide, ce qui signifie que
$\bigcap_{\alpha \in J} Z_{\alpha} = \emptyset$. Il existe donc une
sous-famille finie de $\{Z_{\alpha}\}_{\alpha \in J}$ vérifiant
$\bigcap_{\alpha \in J} Z_{\alpha}=\emptyset$, ce qui conclut par
contraposition.
\end{proof}

\begin{lemma}
\label{lemma-image-quasi-compact}
Soit $f : X \to Y$ une application continue d'espaces topologiques.
\begin{enumerate}
\item Si $X$ est quasi-compact, alors $f(X)$ est quasi-compact.
\item Si $f$ est quasi-compacte, alors $f(X)$ est rétrocompact.
\end{enumerate}
\end{lemma}

\begin{proof}
Si $f(X) = \bigcup V_i$ est un recouvrement ouvert, alors
$X = \bigcup f^{-1}(V_i)$ est un recouvrement ouvert. Par conséquent, si $X$
est quasi-compact, on a
$X = f^{-1}(V_{i_1}) \cup \ldots \cup f^{-1}(V_{i_n})$ pour certains
$i_1, \ldots, i_n \in I$, et donc
$f(X) = V_{i_1} \cup \ldots \cup V_{i_n}$. Cela prouve (1).
Supposons $f$ quasi-compacte et soit $V \subset Y$ un ouvert quasi-compact.
Alors $f^{-1}(V)$ est quasi-compact ; d'après (1),
$f(f^{-1}(V)) = f(X) \cap V$ est donc quasi-compact. Ainsi, $f(X)$ est
rétrocompact.
\end{proof}

\begin{lemma}
\label{lemma-quasi-compact-closed-point}
Soit $X$ un espace topologique. Supposons que
\begin{enumerate}
\item $X$ soit non vide,
\item $X$ soit quasi-compact, et
\item $X$ soit de Kolmogorov.
\end{enumerate}
Alors $X$ possède un point fermé.
\end{lemma}

\begin{proof}
Considérons l'ensemble
$$
\mathcal{T} =
\{Z \subset X \mid Z = \overline{\{x\}} \text{ pour un certain }x \in X\}
$$
de toutes les adhérences de singletons de $X$. Il est non vide, puisque $X$
est non vide. Munissons $\mathcal{T}$ de l'ordre partiel défini par
l'inclusion. Supposons que $Z_\alpha$, $\alpha \in A$, soit une partie
totalement ordonnée de $\mathcal{T}$.
D'après le lemme \ref{lemma-intersection-closed-in-quasi-compact},
$\bigcap_{\alpha \in A} Z_\alpha \not = \emptyset$. Il existe donc
$x \in \bigcap_{\alpha \in A} Z_\alpha$, et
$Z = \overline{\{x\}}\in \mathcal{T}$ est un minorant de la famille.
Le lemme de Zorn assure l'existence d'un élément minimal
$Z \in \mathcal{T}$. Comme $X$ est de Kolmogorov, nous en concluons que
$Z = \{x\}$ pour un certain $x$, et $x \in X$ est un point fermé.
\end{proof}

\begin{lemma}
\label{lemma-closed-points-quasi-compact}
Soit $X$ un espace de Kolmogorov quasi-compact. Alors l'ensemble $X_0$ des
points fermés de $X$ est quasi-compact.
\end{lemma}

\begin{proof}
Soit $X_0 = \bigcup U_{i, 0}$ un recouvrement ouvert.
Écrivons $U_{i, 0} = X_0 \cap U_i$ pour un ouvert $U_i \subset X$.
Considérons le complémentaire $Z$ de $\bigcup U_i$. C'est une partie fermée de
$X$, donc quasi-compacte (lemme \ref{lemma-closed-in-quasi-compact})
et de Kolmogorov. D'après le lemme \ref{lemma-quasi-compact-closed-point},
si $Z$ n'était pas vide, il aurait un point
fermé, ce qui contredit le fait que $X_0 \subset \bigcup U_i$.
Ainsi $Z = \emptyset$ et $X = \bigcup U_i$. Comme $X$ est quasi-compact,
ce recouvrement admet un sous-recouvrement fini, ce qui permet de conclure.
\end{proof}

\begin{lemma}
\label{lemma-connected-component-intersection}
Soit $X$ un espace topologique.
Supposons que
\begin{enumerate}
\item $X$ soit quasi-compact,
\item $X$ possède une base de la topologie formée d'ouverts quasi-compacts, et
\item l'intersection de deux ouverts quasi-compacts soit quasi-compacte.
\end{enumerate}
Pour tout $x \in X$, la composante connexe de $X$ contenant
$x$ est l'intersection de toutes les parties à la fois ouvertes et fermées
de $X$ contenant $x$.
\end{lemma}

\begin{proof}
Soit $T$ la composante connexe contenant $x$.
Soit $S = \bigcap_{\alpha \in A} Z_\alpha$ l'intersection de toutes les
parties $Z_\alpha$ de $X$, à la fois ouvertes et fermées, contenant $x$.
Remarquons que $S$ est fermé dans $X$.
Remarquons que toute intersection finie des $Z_\alpha$ est l'un des $Z_\alpha$.
Puisque $T$ est connexe et que $x \in T$, nous avons $T \subset S$.
Il suffit de montrer que $S$ est connexe.
Sinon, il existe une décomposition en réunion disjointe
$S = B \amalg C$ où $B$ et $C$ sont à la fois ouverts et fermés dans $S$.
En particulier, $B$ et $C$ sont fermés dans $X$, et donc quasi-compacts d'après le
lemme \ref{lemma-closed-in-quasi-compact} et l'hypothèse (1).
D'après l'hypothèse (2), il existe des ouverts quasi-compacts
$U, V \subset X$ tels que $B = S \cap U$ et $C = S \cap V$ (détails omis).
Alors $U \cap V \cap S = \emptyset$.
Ainsi $\bigcap_\alpha U \cap V \cap Z_\alpha = \emptyset$.
D'après l'hypothèse (3), l'intersection $U \cap V$ est quasi-compacte.
D'après le lemme \ref{lemma-intersection-closed-in-quasi-compact},
pour un certain $\alpha' \in A$, nous avons $U \cap V \cap Z_{\alpha'} = \emptyset$.
Comme $X \setminus (U \cup V)$ est disjoint de $S$
et fermé dans $X$, donc quasi-compact, nous pouvons appliquer le même lemme
pour voir que $Z_{\alpha''} \subset U \cup V$ pour un certain $\alpha'' \in A$.
Alors $Z_\alpha = Z_{\alpha'} \cap Z_{\alpha''}$ est contenu
dans $U \cup V$ et disjoint de $U \cap V$.
Ainsi $Z_\alpha = U \cap Z_\alpha \amalg V \cap Z_\alpha$
est une décomposition en deux parties ouvertes,
donc $U \cap Z_\alpha$ et $V \cap Z_\alpha$ sont ouverts et fermés dans $X$.
Ainsi, si $x \in B$ par exemple, alors nous voyons que $S \subset U \cap Z_\alpha$,
et nous concluons que $C = \emptyset$.
\end{proof}

\begin{lemma}
\label{lemma-connected-component-intersection-compact-Hausdorff}
Soit $X$ un espace topologique. Supposons que $X$ soit quasi-compact et séparé.
Pour tout $x \in X$, la composante connexe de $X$ contenant
$x$ est l'intersection de toutes les parties à la fois ouvertes et fermées
de $X$ contenant $x$.
\end{lemma}

\begin{proof}
Soit $T$ la composante connexe contenant $x$.
Soit $S = \bigcap_{\alpha \in A} Z_\alpha$ l'intersection de toutes les
parties $Z_\alpha$ de $X$, à la fois ouvertes et fermées, contenant $x$.
Remarquons que $S$ est fermé dans $X$.
Remarquons que toute intersection finie des $Z_\alpha$ est l'un des $Z_\alpha$.
Puisque $T$ est connexe et que $x \in T$, nous avons $T \subset S$.
Il suffit de montrer que $S$ est connexe.
Sinon, il existe une décomposition en réunion disjointe
$S = B \amalg C$ où $B$ et $C$ sont à la fois ouverts et fermés dans $S$.
En particulier, $B$ et $C$ sont fermés dans $X$, et donc quasi-compacts d'après le
lemme \ref{lemma-closed-in-quasi-compact}.
D'après le lemme \ref{lemma-quasi-compact-in-Hausdorff},
il existe des ouverts disjoints $U, V \subset X$ tels que $B \subset U$ et
$C \subset V$. Alors $X \setminus U \cup V$ est fermé dans $X$,
donc quasi-compact (lemme \ref{lemma-closed-in-quasi-compact}).
Il s'ensuit que $(X \setminus U \cup V) \cap Z_\alpha = \emptyset$
pour un certain $\alpha$, d'après le lemme \ref{lemma-intersection-closed-in-quasi-compact}.
Autrement dit, $Z_\alpha \subset U \cup V$. Ainsi,
$Z_\alpha = Z_\alpha \cap V \amalg Z_\alpha \cap U$
est une décomposition en deux parties ouvertes,
donc $U \cap Z_\alpha$ et $V \cap Z_\alpha$ sont ouverts et fermés dans $X$.
Ainsi, si $x \in B$ par exemple, alors nous voyons que $S \subset U \cap Z_\alpha$,
et nous concluons que $C = \emptyset$.
\end{proof}

\begin{lemma}
\label{lemma-closed-union-connected-components}
Soit $X$ un espace topologique.
Supposons que
\begin{enumerate}
\item $X$ soit quasi-compact,
\item $X$ possède une base de la topologie formée d'ouverts quasi-compacts, et
\item l'intersection de deux ouverts quasi-compacts soit quasi-compacte.
\end{enumerate}
Pour une partie $T \subset X$, les assertions suivantes sont équivalentes :
\begin{enumerate}
\item[(a)] $T$ est une intersection de parties à la fois ouvertes et fermées de $X$, et
\item[(b)] $T$ est fermée dans $X$ et est une réunion de composantes connexes de $X$.
\end{enumerate}
\end{lemma}

\begin{proof}
Il est clair que (a) implique (b).
Supposons (b). Soit $x \in X$, $x \not \in T$. Soit $x \in C \subset X$
la composante connexe de $X$ contenant $x$. D'après le
lemme \ref{lemma-connected-component-intersection},
nous voyons que $C = \bigcap V_\alpha$ est l'intersection de toutes les parties à la fois ouvertes et
fermées $V_\alpha$ de $X$ qui contiennent $C$.
En particulier, toute intersection $V_\alpha \cap V_\beta$ de deux de ces parties
est elle-même un $V_\alpha$.
Comme $T$ est une réunion de composantes connexes
de $X$, nous voyons que $C \cap T = \emptyset$. Ainsi,
$T \cap \bigcap V_\alpha = \emptyset$. Puisque $T$ est quasi-compacte en tant que
partie fermée d'un espace quasi-compact (voir le
lemme \ref{lemma-closed-in-quasi-compact}),
nous en déduisons que $T \cap V_\alpha = \emptyset$ pour un certain $\alpha$, voir le
lemme \ref{lemma-intersection-closed-in-quasi-compact}.
Pour cet $\alpha$, nous voyons que $U_\alpha = X \setminus V_\alpha$
est une partie à la fois ouverte et fermée de $X$ qui contient $T$ mais pas $x$.
Le lemme en découle.
\end{proof}

\begin{lemma}
\label{lemma-Noetherian-quasi-compact}
Soit $X$ un espace topologique noethérien.
\begin{enumerate}
\item L'espace $X$ est quasi-compact.
\item Toute partie de $X$ est rétrocompacte.
\end{enumerate}
\end{lemma}

\begin{proof}
Supposons que $X = \bigcup U_i$ soit un recouvrement ouvert de $X$ indexé
par l'ensemble $I$ et n'admettant pas de raffinement par un
recouvrement ouvert fini. Choisissons $i_1, i_2, \ldots $ dans $I$ par récurrence
de la manière suivante : choisissons $i_{n + 1}$ tel que $U_{i_{n + 1}}$
ne soit pas contenu dans $U_{i_1} \cup \ldots \cup U_{i_n}$. Ainsi, nous voyons que
$X \supset (X \setminus U_{i_1}) \supset
(X \setminus U_{i_1} \cup U_{i_2}) \supset \ldots$ est une suite infinie strictement
décroissante de parties fermées. Cela contredit
le fait que $X$ est noethérien. Cela démontre la première assertion.
La seconde assertion est maintenant claire, puisque toute partie de $X$ est noethérienne d'après le
lemme \ref{lemma-Noetherian}.
\end{proof}

\begin{lemma}
\label{lemma-quasi-compact-locally-Noetherian-Noetherian}
Un espace quasi-compact et localement noethérien est noethérien.
\end{lemma}

\begin{proof}
Les hypothèses impliquent immédiatement que $X$ admet un recouvrement fini par
des parties noethériennes, et est donc noethérien d'après le
lemme \ref{lemma-finite-union-Noetherian}.
\end{proof}

\begin{lemma}[Théorème de la sous-base d'Alexander]
\label{lemma-subbase-theorem}
Soit $X$ un espace topologique. Soit $\mathcal{B}$ une sous-base de $X$.
Si tout recouvrement de $X$ par des éléments de $\mathcal{B}$ admet un
raffinement fini, alors $X$ est quasi-compact.
\end{lemma}

\begin{proof}
Supposons qu'il existe un recouvrement ouvert de $X$ qui n'admette aucun
raffinement fini. À l'aide du lemme de Zorn, nous pouvons choisir un recouvrement ouvert maximal
$X = \bigcup_{i \in I} U_i$ qui n'admette aucun raffinement fini
(détails omis).
Autrement dit, si $U \subset X$ est un ouvert quelconque qui ne figure pas
parmi les $U_i$, alors le recouvrement $X = U \cup \bigcup_{i \in I} U_i$
admet un raffinement fini. Soit $I' \subset I$ l'ensemble des indices
tels que $U_i \in \mathcal{B}$. Alors $\bigcup_{i \in I'} U_i \not = X$,
car sinon nous obtiendrions un raffinement fini recouvrant $X$ en vertu de notre
hypothèse sur $\mathcal{B}$. Choisissons $x \in X$,
$x \not \in \bigcup_{i \in I'} U_i$. Choisissons $i \in I$ tel que $x \in U_i$.
Choisissons $V_1, \ldots, V_n \in \mathcal{B}$ tels que
$x \in V_1 \cap \ldots \cap V_n \subset U_i$. Cela est
possible puisque $\mathcal{B}$ est une sous-base de $X$. Remarquons que
$V_j$ ne figure pas parmi les $U_i$. Par maximalité du
recouvrement choisi, nous voyons que pour tout $j$, il existe
$i_{j, 1}, \ldots, i_{j, n_j} \in I$ tels que
$X = V_j \cup U_{i_{j, 1}} \cup \ldots \cup U_{i_{j, n_j}}$.
Puisque $V_1 \cap \ldots \cap V_n \subset U_i$, nous concluons que
$X = U_i \cup \bigcup U_{i_{j, l}}$, ce qui est une contradiction.
\end{proof}







\section{Espaces localement quasi-compacts}
\label{section-locally-quasi-compact}

\noindent
Rappelons qu'un voisinage d'un point n'est pas nécessairement ouvert.

\begin{definition}
\label{definition-locally-quasi-compact}
Un espace topologique $X$ est dit {\it localement
quasi-compact}\footnote{Cette terminologie n'est peut-être pas standard.
On rencontre dans la littérature les deux autres notions suivantes : (1) tout
point possède un voisinage quasi-compact ; (2) tout point possède un voisinage
fermé quasi-compact. Un schéma a la propriété que tout point possède un système
fondamental de voisinages ouverts quasi-compacts.} si tout point possède un
système fondamental de voisinages quasi-compacts.
\end{definition}

\noindent
Dans la littérature, l'expression {\it espace localement compact} désigne
souvent un espace comme dans le lemme suivant.

\begin{lemma}
\label{lemma-locally-quasi-compact-Hausdorff}
Un espace séparé est localement quasi-compact si et seulement si tout point
possède un voisinage quasi-compact.
\end{lemma}

\begin{proof}
Soit $X$ un espace séparé. Soit $x \in X$ et soit
$x \in E \subset X$ un voisinage quasi-compact. Alors $E$ est fermé d'après le
lemme \ref{lemma-quasi-compact-in-Hausdorff}.
Supposons que $x \in U \subset X$ soit un voisinage ouvert de $x$.
Alors $Z = E \setminus U$ est une partie fermée de $E$ qui ne contient pas
$x$. Il existe donc deux ouverts disjoints $W, V \subset E$ de $E$ tels que
$x \in V$ et $Z \subset W$ ; voir le
lemme \ref{lemma-quasi-compact-in-Hausdorff}.
Il s'ensuit que $\overline{V} \subset E$ est un voisinage fermé de $x$ contenu
dans $E \cap U$. De plus, $\overline{V}$ est quasi-compact en tant que partie
fermée de $E$ (lemme \ref{lemma-closed-in-quasi-compact}).
Nous obtenons ainsi un système fondamental de voisinages quasi-compacts de $x$.
\end{proof}

\begin{lemma}[Théorème de Baire]
\label{lemma-baire-category-locally-compact}
Soit $X$ un espace séparé localement quasi-compact.
Soient $U_n \subset X$, $n \geq 1$, des ouverts denses. Alors
$\bigcap_{n \geq 1} U_n$ est dense dans $X$.
\end{lemma}

\begin{proof}
En remplaçant $U_n$ par $\bigcap_{i = 1, \ldots, n} U_i$, nous pouvons
supposer que $U_1 \supset U_2 \supset \ldots$.
Soit $x \in X$. Montrons que $x$ appartient à l'adhérence de
$\bigcap_{n \geq 1} U_n$. Soit donc $E$ un voisinage de $x$.
Pour montrer que $E \cap \bigcap_{n \geq 1} U_n$ est non vide, nous pouvons
remplacer $E$ par un voisinage plus petit. Après avoir remplacé $E$ par un
voisinage plus petit, nous pouvons supposer que $E$ est quasi-compact.

\medskip\noindent
Posons $x_0 = x$ et $E_0 = E$. Nous choisirons par récurrence un point
$x_i \in E_{i - 1} \cap U_i$ et un voisinage quasi-compact $E_i$ de $x_i$ tel
que $E_i \subset E_{i - 1} \cap U_i$.
Comme $X$ est séparé, les parties $E_i \subset X$ sont fermées
(lemme \ref{lemma-quasi-compact-in-Hausdorff}).
Puisque les $E_i$ sont aussi non vides, nous concluons que
$\bigcap_{i \geq 1} E_i$ est non vide
(lemme \ref{lemma-intersection-closed-in-quasi-compact}).
Comme $\bigcap_{i \geq 1} E_i \subset E \cap \bigcap_{n \geq 1} U_n$,
le lemme est démontré.

\medskip\noindent
Le cas initial $i = 0$ a été traité ci-dessus. Passons à l'étape de
récurrence. Puisque $E_{i - 1}$ est un voisinage de $x_{i - 1}$, il existe un
ouvert $x_{i - 1} \in W \subset E_{i - 1}$.
Comme $U_i$ est dense dans $X$, l'intersection $W \cap U_i$ est non vide.
Choisissons un point $x_i \in W \cap U_i$ quelconque.
Par définition des espaces localement quasi-compacts, il existe un voisinage
quasi-compact $E_i$ de $x_i$ contenu dans $W \cap U_i$.
Alors $E_i \subset E_{i - 1} \cap U_i$, comme voulu.
\end{proof}

\begin{lemma}
\label{lemma-relatively-compact-refinement}
Soit $X$ un espace séparé et quasi-compact.
Soit $X = \bigcup_{i \in I} U_i$ un recouvrement ouvert.
Il existe alors un recouvrement ouvert $X = \bigcup_{i \in I} V_i$ tel que
$\overline{V_i} \subset U_i$ pour tout $i$.
\end{lemma}

\begin{proof}
Soit $x \in X$. Choisissons $i(x) \in I$ tel que $x \in U_{i(x)}$.
Comme $X \setminus U_{i(x)}$ et $\{x\}$ sont deux parties fermées disjointes
de $X$, les lemmes \ref{lemma-closed-in-quasi-compact} et
\ref{lemma-quasi-compact-in-Hausdorff} fournissent un voisinage ouvert $U_x$
de $x$ dont l'adhérence ne rencontre pas $X \setminus U_{i(x)}$.
Ainsi, $\overline{U_x} \subset U_{i(x)}$. Comme $X$ est quasi-compact, il
existe une liste finie de points $x_1, \ldots, x_m$ telle que
$X = U_{x_1} \cup \ldots \cup U_{x_m}$. En posant
$V_i = \bigcup_{i = i(x_j)} U_{x_j}$, on achève la démonstration.
\end{proof}

\begin{lemma}
\label{lemma-refine-covering}
Soit $X$ un espace séparé et quasi-compact.
Soit $X = \bigcup_{i \in I} U_i$ un recouvrement ouvert.
Supposons donnés un entier $p \geq 0$ et, pour tout $(p + 1)$-uplet
$i_0, \ldots, i_p$ d'éléments de $I$, un recouvrement ouvert
$U_{i_0} \cap \ldots \cap U_{i_p} = \bigcup W_{i_0 \ldots i_p, k}$.
Il existe alors un recouvrement ouvert $X = \bigcup_{j \in J} V_j$ et une
application $\alpha : J \to I$ tels que $\overline{V_j} \subset U_{\alpha(j)}$
et que chaque $V_{j_0} \cap \ldots \cap V_{j_p}$ soit contenu dans
$W_{\alpha(j_0) \ldots \alpha(j_p), k}$ pour un certain $k$.
\end{lemma}

\begin{proof}
Comme $X$ est quasi-compact, on se ramène au cas où $I$ est fini (détails
omis). Nous démontrons le résultat pour $I$ fini par récurrence sur $p$.
Le cas initial $p = 0$ est immédiat : il suffit de prendre un recouvrement
comme dans le lemme \ref{lemma-relatively-compact-refinement}, qui raffine le
recouvrement ouvert $X = \bigcup W_{i_0, k}$.

\medskip\noindent
Passons à l'étape de récurrence. Supposons le lemme démontré pour $p - 1$.
Pour tout $p$-uplet $i'_0, \ldots, i'_{p - 1}$ d'éléments de $I$, choisissons
$U_{i'_0} \cap \ldots \cap U_{i'_{p - 1}} =
\bigcup W_{i'_0 \ldots i'_{p - 1}, k}$ comme raffinement commun des
recouvrements
$U_{i_0} \cap \ldots \cap U_{i_p} = \bigcup W_{i_0 \ldots i_p, k}$ associés
aux $(p + 1)$-uplets tels que
$\{i'_0, \ldots, i'_{p - 1}\} = \{i_0, \ldots, i_p\}$ (égalité d'ensembles).
(Il n'y en a qu'un nombre fini, puisque $I$ est fini.)
Par récurrence, il existe une solution pour ces ouverts, disons
$X = \bigcup V_j$ et $\alpha : J \to I$.
À ce stade, le recouvrement $X = \bigcup_{j \in J} V_j$ et $\alpha$ vérifient
$\overline{V_j} \subset U_{\alpha(j)}$, et chaque
$V_{j_0} \cap \ldots \cap V_{j_p}$ est contenu dans
$W_{\alpha(j_0) \ldots \alpha(j_p), k}$ pour un certain $k$ dès qu'il y a une
répétition parmi $\alpha(j_0), \ldots, \alpha(j_p)$. Bien entendu, nous pouvons
et allons supposer que $J$ est fini.

\medskip\noindent
Fixons $i_0, \ldots, i_p \in I$ deux à deux distincts. Considérons les
$(p + 1)$-uplets $j_0, \ldots, j_p \in J$ tels que
$i_0 = \alpha(j_0), \ldots, i_p = \alpha(j_p)$ et pour lesquels
$V_{j_0} \cap \ldots \cap V_{j_p}$ n'est {\bf pas} contenu dans
$W_{\alpha(j_0) \ldots \alpha(j_p), k}$, quel que soit $k$.
Soit $N$ le nombre de ces $(p + 1)$-uplets. Montrons comment diminuer $N$.
Comme
$$
\overline{V_{j_0}} \cap \ldots \cap \overline{V_{j_p}} \subset
U_{i_0} \cap \ldots \cap U_{i_p} = \bigcup W_{i_0 \ldots i_p, k}
$$
il existe un ensemble fini $K = \{k_1, \ldots, k_t\}$ tel que le membre de
gauche soit contenu dans $\bigcup_{k \in K} W_{i_0 \ldots i_p, k}$.
Considérons alors le recouvrement ouvert
$$
V_{j_0} =
(V_{j_0} \setminus (\overline{V_{j_1}} \cap \ldots \cap \overline{V_{j_p}}))
\cup (\bigcup\nolimits_{k \in K} V_{j_0} \cap W_{i_0 \ldots i_p, k})
$$
Le premier ouvert du membre de droite a une intersection vide avec
$V_{j_1} \cap \ldots \cap V_{j_p}$, tandis que les
autres ouverts $V_{j_0, k}$ du membre de droite satisfont
$V_{j_0, k} \cap V_{j_1} \ldots \cap V_{j_p} \subset
W_{\alpha(j_0) \ldots \alpha(j_p), k}$.
Posons $J' = J \amalg K$. Pour $j \in J$, posons $V'_j = V_j$ si
$j \not = j_0$, et posons
$V'_{j_0} =
V_{j_0} \setminus (\overline{V_{j_1}} \cap \ldots \cap \overline{V_{j_p}})$.
Pour $k \in K$, posons $V'_k = V_{j_0, k}$. Enfin, l'application
$\alpha' : J' \to I$ est égale à $\alpha$ sur $J$ et envoie tout élément de
$K$ sur $i_0$. Une vérification immédiate montre que ce remplacement diminue
$N$ d'une unité. En répétant cette procédure $N$ fois, nous arrivons à la
situation où $N = 0$.

\medskip\noindent
Pour achever la démonstration, raisonnons par récurrence sur le nombre $M$ de
$(p + 1)$-uplets $i_0, \ldots, i_p \in I$ à entrées deux à deux distinctes
pour lesquels il existe un $(p + 1)$-uplet $j_0, \ldots, j_p \in J$ tel que
$i_0 = \alpha(j_0), \ldots, i_p = \alpha(j_p)$ et tel que
$V_{j_0} \cap \ldots \cap V_{j_p}$ ne soit {\bf pas} contenu dans
$W_{\alpha(j_0) \ldots \alpha(j_p), k}$, quel que soit $k$. À cette fin, nous
affirmons que l'opération du paragraphe précédent n'augmente pas $M$.
Cela résulte formellement du fait que l'application $\alpha' : J' \to I$ se
factorise par une application $\beta : J' \to J$ telle que
$V'_{j'} \subset V_{\beta(j')}$.
\end{proof}

\begin{lemma}
\label{lemma-lift-covering-of-a-closed}
Soit $X$ un espace séparé et localement quasi-compact.
Soit $Z \subset X$ une partie quasi-compacte (donc fermée).
Supposons donnés un entier $p \geq 0$, un ensemble $I$, pour tout $i \in I$
un ouvert $U_i \subset X$, et pour tout $(p + 1)$-uplet
$i_0, \ldots, i_p$ d'éléments de $I$ un ouvert
$W_{i_0 \ldots i_p} \subset U_{i_0} \cap \ldots \cap U_{i_p}$
tel que
\begin{enumerate}
\item $Z \subset \bigcup U_i$, et
\item pour tous $i_0, \ldots, i_p$, on ait
$W_{i_0 \ldots i_p} \cap Z = U_{i_0} \cap \ldots \cap U_{i_p} \cap Z$.
\end{enumerate}
Il existe alors des ouverts $V_i$ de $X$ tels que
$Z \subset \bigcup V_i$,
pour tout $i$, on ait $\overline{V_i} \subset U_i$, et
$V_{i_0} \cap \ldots \cap V_{i_p} \subset W_{i_0 \ldots i_p}$
pour tout $(p + 1)$-uplet $i_0, \ldots, i_p$.
\end{lemma}

\begin{proof}
Puisque $Z$ est quasi-compact, on se ramène au cas où $I$ est fini (détails
omis). Comme $X$ est localement quasi-compact et $Z$ quasi-compact, nous
pouvons trouver un voisinage $Z \subset E$ qui soit quasi-compact,
c'est-à-dire que $E$ est quasi-compact et contient un voisinage ouvert de
$Z$ dans $X$. Si nous démontrons le résultat après avoir remplacé $X$ par
$E$, le résultat s'ensuit. Nous pouvons donc supposer $X$ quasi-compact.

\medskip\noindent
Nous démontrons le résultat lorsque $I$ est fini et $X$ quasi-compact, par
récurrence sur $p$. Le cas initial est $p = 0$. Dans ce cas,
$X = (X \setminus Z) \cup \bigcup W_i$. D'après le
lemme \ref{lemma-relatively-compact-refinement}, il existe un recouvrement
$X = V \cup \bigcup V_i$ par des ouverts $V_i \subset W_i$ et
$V \subset X \setminus Z$, avec $\overline{V_i} \subset W_i$ pour tout $i$.
On voit alors que l'on obtient une solution du problème posé dans le lemme.

\medskip\noindent
Passons à l'étape de récurrence. Supposons le lemme démontré pour $p - 1$.
Posons $W_{j_0 \ldots j_{p - 1}}$ égal à l'intersection de tous les
$W_{i_0 \ldots i_p}$ tels que
$\{j_0, \ldots, j_{p - 1}\} = \{i_0, \ldots, i_p\}$ (égalité d'ensembles).
Par récurrence, il existe une solution pour ces ouverts, disons
$V_i \subset U_i$.
Il résulte de notre choix de $W_{j_0 \ldots j_{p - 1}}$ que
$V_{i_0} \cap \ldots \cap V_{i_p} \subset W_{i_0 \ldots i_p}$
pour tout $(p + 1)$-uplet $i_0, \ldots, i_p$ tel que $i_a = i_b$ pour
certains $0 \leq a < b \leq p$.
Il ne reste donc à modifier notre choix des $V_i$ que si
$V_{i_0} \cap \ldots \cap V_{i_p} \not \subset W_{i_0 \ldots i_p}$
pour un $(p + 1)$-uplet $i_0, \ldots, i_p$ dont les éléments sont deux à
deux distincts. Dans ce cas,
$$
T =
\overline{V_{i_0} \cap \ldots \cap V_{i_p} \setminus W_{i_0 \ldots i_p}}
\subset 
\overline{V_{i_0}} \cap \ldots \cap \overline{V_{i_p}} \setminus
W_{i_0 \ldots i_p}
$$
est un fermé de $X$, contenu dans $U_{i_0} \cap \ldots \cap U_{i_p}$ et
disjoint de $Z$. Nous pouvons donc remplacer $V_{i_0}$ par
$V_{i_0} \setminus T$ pour « corriger » le problème. Après avoir répété cette
opération un nombre fini de fois pour chacun des uplets problématiques, le
lemme est démontré.
\end{proof}

\begin{lemma}
\label{lemma-lift-covering-of-quasi-compact-hausdorff-subset}
Soit $X$ un espace topologique. Soit $Z \subset X$ une partie quasi-compacte
telle que deux points quelconques de $Z$ admettent dans $X$ des voisinages
ouverts disjoints. Supposons donnés un entier $p \geq 0$, un ensemble $I$,
pour tout $i \in I$ un ouvert $U_i \subset X$, et pour tout $(p + 1)$-uplet
$i_0, \ldots, i_p$ d'éléments de $I$ un ouvert
$W_{i_0 \ldots i_p} \subset U_{i_0} \cap \ldots \cap U_{i_p}$
tel que
\begin{enumerate}
\item $Z \subset \bigcup U_i$, et
\item pour tous $i_0, \ldots, i_p$, on ait
$W_{i_0 \ldots i_p} \cap Z = U_{i_0} \cap \ldots \cap U_{i_p} \cap Z$.
\end{enumerate}
Il existe alors des ouverts $V_i$ de $X$ tels que
\begin{enumerate}
\item $Z \subset \bigcup V_i$,
\item $V_i \subset U_i$ pour tout $i$,
\item $\overline{V_i} \cap Z \subset U_i$ pour tout $i$, et
\item $V_{i_0} \cap \ldots \cap V_{i_p} \subset W_{i_0 \ldots i_p}$
pour tout $(p + 1)$-uplet $i_0, \ldots, i_p$.
\end{enumerate}
\end{lemma}

\begin{proof}
Puisque $Z$ est quasi-compact, on se ramène au cas où $I$ est fini (détails
omis). Nous démontrons le résultat lorsque $I$ est fini, par récurrence sur
$p$.

\medskip\noindent
Le cas initial est $p = 0$.
Pour $z \in Z \cap U_i$ et $z' \in Z \setminus U_i$, il existe des ouverts
disjoints $z \in V_{z, z'}$ et $z' \in W_{z, z'}$ de $X$.
Comme $Z \setminus U_i$ est quasi-compact
(lemme \ref{lemma-closed-in-quasi-compact}), nous pouvons choisir un nombre
fini de points $z'_1, \ldots, z'_r$ tels que
$Z \setminus U_i \subset W_{z, z'_1} \cup \ldots \cup W_{z, z'_r}$.
On voit alors que
$V_z = V_{z, z'_1} \cap \ldots \cap V_{z, z'_r} \cap U_i$
est un voisinage ouvert de $z$ contenu dans $U_i$ et vérifiant
$\overline{V_z} \cap Z \subset U_i$.
Comme $z$ et $i$ étaient arbitraires et que $Z$ est quasi-compact, nous
pouvons trouver une liste finie $z_1, i_1, \ldots, z_t, i_t$ et des ouverts
$V_{z_j} \subset U_{i_j}$ tels que
$\overline{V_{z_j}} \cap Z \subset U_{i_j}$
et $Z \subset \bigcup V_{z_j}$.
Nous pouvons alors poser $V_i = W_i \cap (\bigcup_{j : i = i_j} V_{z_j})$
pour résoudre le problème lorsque $p = 0$.

\medskip\noindent
Passons à l'étape de récurrence. Supposons le lemme démontré pour $p - 1$.
Posons $W_{j_0 \ldots j_{p - 1}}$ égal à l'intersection de tous les
$W_{i_0 \ldots i_p}$ tels que
$\{j_0, \ldots, j_{p - 1}\} = \{i_0, \ldots, i_p\}$ (égalité d'ensembles).
Par récurrence, il existe une solution pour ces ouverts, disons
$V_i \subset U_i$.
Il résulte de notre choix de $W_{j_0 \ldots j_{p - 1}}$ que
$V_{i_0} \cap \ldots \cap V_{i_p} \subset W_{i_0 \ldots i_p}$
pour tout $(p + 1)$-uplet $i_0, \ldots, i_p$ tel que $i_a = i_b$ pour
certains $0 \leq a < b \leq p$.
Il ne reste donc à modifier notre choix des $V_i$ que si
$V_{i_0} \cap \ldots \cap V_{i_p} \not \subset W_{i_0 \ldots i_p}$
pour un $(p + 1)$-uplet $i_0, \ldots, i_p$ dont les éléments sont deux à
deux distincts. Dans ce cas,
$$
T =
\overline{V_{i_0} \cap \ldots \cap V_{i_p} \setminus W_{i_0 \ldots i_p}}
\subset
\overline{V_{i_0}} \cap \ldots \cap \overline{V_{i_p}} \setminus
W_{i_0 \ldots i_p}
$$
est un fermé de $X$ disjoint de $Z$, d'après la propriété (3) des ouverts
$V_i$. Nous pouvons donc remplacer $V_{i_0}$ par $V_{i_0} \setminus T$ pour
« corriger » le problème. Après avoir répété cette opération un nombre fini
de fois pour chacun des uplets problématiques, le lemme est démontré.
\end{proof}







\section{Limites d'espaces}
\label{section-limits}

\noindent
La catégorie des espaces topologiques admet les produits. Plus précisément,
si $I$ est un ensemble et si, pour $i \in I$, on se donne un espace
topologique $X_i$, on munit $\prod_{i \in I} X_i$ de la {\it topologie
produit}. On prend pour base de cette topologie les ensembles de la forme
$\prod U_i$, où $U_i \subset X_i$ est ouvert et $U_i = X_i$ pour presque
tout $i$.

\medskip\noindent
La catégorie des espaces topologiques admet les égalisateurs. Plus
précisément, si $a, b : X \to Y$ sont des morphismes d'espaces topologiques,
l'égalisateur de $a$ et $b$ est la partie
$\{x \in X \mid a(x) = b(x)\} \subset X$, munie de la topologie induite.

\begin{lemma}
\label{lemma-limits}
La catégorie des espaces topologiques admet les limites, et le foncteur
d'oubli vers les ensembles commute à celles-ci.
\end{lemma}

\begin{proof}
Cela résulte de la discussion précédente et de Catégories,
lemme \ref{categories-lemma-limits-products-equalizers}.
La description précédente montre que le foncteur d'oubli commute aux
limites. On peut aussi le voir en appliquant Catégories,
lemme \ref{categories-lemma-adjoint-exact}, et en observant que le foncteur
d'oubli admet un adjoint à gauche, à savoir le foncteur qui associe à un
ensemble l'espace topologique discret correspondant.
\end{proof}

\begin{lemma}
\label{lemma-describe-limits}
Soit $\mathcal{I}$ une catégorie cofiltrante. Soit $i \mapsto X_i$ un
diagramme d'espaces topologiques indexé par $\mathcal{I}$. Soit
$X = \lim X_i$ sa limite, avec applications de projection $f_i : X \to X_i$.
\begin{enumerate}
\item Tout ouvert de $X$ est de la forme $\bigcup_{j \in J} f_j^{-1}(U_j)$
pour une partie $J \subset I$ et des ouverts $U_j \subset X_j$.
\item Tout ouvert quasi-compact de $X$ est de la forme
$f_i^{-1}(U_i)$ pour un certain $i$ et un ouvert $U_i \subset X_i$.
\end{enumerate}
\end{lemma}

\begin{proof}
La construction de la limite donnée ci-dessus montre que
$X \subset \prod X_i$, avec la topologie induite. Une base de la topologie de
$\prod X_i$ est donnée par les ouverts $\prod U_i$, où $U_i \subset X_i$ est
ouvert et $U_i = X_i$ pour presque tout $i$. Soient
$i_1, \ldots, i_n \in \Ob(\mathcal{I})$ les objets tels que
$U_{i_j} \not = X_{i_j}$. Alors
$$
X \cap \prod U_i = f_{i_1}^{-1}(U_{i_1}) \cap \ldots \cap f_{i_n}^{-1}(U_{i_n})
$$
Dans le cas d'une limite quelconque d'espaces topologiques, ces ouverts
forment une base de la topologie de $X$. Toutefois, si $\mathcal{I}$ est
cofiltrante comme dans l'énoncé du lemme, nous pouvons choisir un
$j \in \Ob(\mathcal{I})$ et des morphismes $j \to i_l$,
$l = 1, \ldots, n$. Posons
$$
U_j =
(X_j \to X_{i_1})^{-1}(U_{i_1}) \cap \ldots \cap
(X_j \to X_{i_n})^{-1}(U_{i_n})
$$
Il est alors clair que $X \cap \prod U_i = f_j^{-1}(U_j)$. Ainsi, pour tout
ouvert $W$ de $X$, il existe un ensemble $A$, une application
$\alpha : A \to \Ob(\mathcal{I})$ et des ouverts
$U_a \subset X_{\alpha(a)}$ tels que
$W = \bigcup f_{\alpha(a)}^{-1}(U_a)$. Posons $J = \Im(\alpha)$ et, pour
$j \in J$, posons $U_j = \bigcup_{\alpha(a) = j} U_a$ ; on obtient alors
$W = \bigcup_{j \in J} f_j^{-1}(U_j)$.
Ceci démontre (1).

\medskip\noindent
Pour démontrer (2), supposons que
$\bigcup_{j \in J} f_j^{-1}(U_j)$ soit quasi-compact. Cet ouvert est alors
égal à $f_{j_1}^{-1}(U_{j_1}) \cup \ldots \cup f_{j_m}^{-1}(U_{j_m})$
pour certains $j_1, \ldots, j_m \in J$. Puisque $\mathcal{I}$ est
cofiltrante, nous pouvons choisir un $i \in \Ob(\mathcal{I})$ et des
morphismes $i \to j_l$, $l = 1, \ldots, m$. Posons
$$
U_i =
(X_i \to X_{j_1})^{-1}(U_{j_1}) \cup \ldots \cup
(X_i \to X_{j_m})^{-1}(U_{j_m})
$$
Notre ouvert est alors égal à $f_i^{-1}(U_i)$, comme voulu.
\end{proof}

\begin{lemma}
\label{lemma-characterize-limit}
Soit $\mathcal{I}$ une catégorie cofiltrante. Soit $i \mapsto X_i$ un
diagramme d'espaces topologiques indexé par $\mathcal{I}$. Soit $X$ un espace
topologique tel que
\begin{enumerate}
\item $X = \lim X_i$ comme ensemble (on note $f_i$ les applications de
projection),
\item les ensembles $f_i^{-1}(U_i)$, pour $i \in \Ob(\mathcal{I})$ et
$U_i \subset X_i$ ouvert, forment une base de la topologie de $X$.
\end{enumerate}
Alors $X$ est la limite des $X_i$ dans la catégorie des espaces topologiques.
\end{lemma}

\begin{proof}
Cela résulte de la description de la topologie limite donnée dans le
lemme \ref{lemma-describe-limits}.
\end{proof}

\begin{theorem}[Tychonov]
\label{theorem-tychonov}
Tout produit d'espaces quasi-compacts est quasi-compact.
\end{theorem}

\begin{proof}
Soit $I$ un ensemble et, pour $i \in I$, soit $X_i$ un espace topologique
quasi-compact. Posons $X = \prod X_i$. Soit $\mathcal{B}$ l'ensemble des
parties de $X$ de la forme
$U_i \times \prod_{i' \in I, i' \not = i} X_{i'}$, où
$U_i \subset X_i$ est ouvert. Par construction, cette famille est une
sous-base de la topologie de $X$. D'après le
lemme \ref{lemma-subbase-theorem}, il suffit de montrer que tout recouvrement
$X = \bigcup_{j \in J} B_j$ par des éléments $B_j$ de $\mathcal{B}$ admet un
sous-recouvrement fini. Nous pouvons décomposer $J = \coprod J_i$ de telle
sorte que, si $j \in J_i$, alors
$B_j = U_j \times \prod_{i' \not = i} X_{i'}$, avec $U_j \subset X_i$
ouvert. Si $X_i = \bigcup_{j \in J_i} U_j$, il existe un sous-recouvrement
fini et nous concluons que $X = \bigcup_{j \in J} B_j$ admet un
sous-recouvrement fini. Si ce n'est pas le cas, nous pouvons choisir, pour
tout $i$, un point $x_i \in X_i$ qui n'appartient pas à
$\bigcup_{j \in J_i} U_j$. Mais alors le point $x = (x_i)_{i \in I}$ est un
élément de $X$ qui n'appartient pas à $\bigcup_{j \in J} B_j$, contradiction.
\end{proof}

\noindent
Le lemme suivant devient faux si l'on omet l'hypothèse que les espaces $X_i$
sont séparés ; voir Exemples,
section \ref{examples-section-lim-not-quasi-compact}.

\begin{lemma}
\label{lemma-inverse-limit-quasi-compact}
Soit $\mathcal{I}$ une catégorie et soit $i \mapsto X_i$ un diagramme indexé
par $\mathcal{I}$ dans la catégorie des espaces topologiques. Si chaque $X_i$
est quasi-compact et séparé, alors $\lim X_i$ est quasi-compact.
\end{lemma}

\begin{proof}
Rappelons que $\lim X_i$ est un sous-espace de $\prod X_i$. D'après le
théorème \ref{theorem-tychonov}, ce produit est quasi-compact. Il suffit donc
de montrer que $\lim X_i$ est un sous-espace fermé de $\prod X_i$
(lemme \ref{lemma-closed-in-quasi-compact}).
Si $\varphi : j \to k$ est un morphisme de $\mathcal{I}$, notons
$\Gamma_\varphi \subset X_j \times X_k$ le graphe de l'application continue
correspondante $X_j \to X_k$. D'après le lemme \ref{lemma-graph-closed}, ce
graphe est fermé. Il est clair que $\lim X_i$ est l'intersection des fermés
$$
\Gamma_\varphi \times \prod\nolimits_{l \not = j, k} X_l
\subset \prod X_i
$$
Le résultat s'ensuit.
\end{proof}

\noindent
Le lemme suivant généralise Catégories,
lemme \ref{categories-lemma-nonempty-limit}, et généralise partiellement le
lemme \ref{lemma-intersection-closed-in-quasi-compact}.

\begin{lemma}
\label{lemma-nonempty-limit}
Soit $\mathcal{I}$ une catégorie cofiltrante et soit $i \mapsto X_i$ un
diagramme indexé par $\mathcal{I}$ dans la catégorie des espaces
topologiques. Si chaque $X_i$ est quasi-compact, séparé et non vide, alors
$\lim X_i$ est non vide.
\end{lemma}

\begin{proof}
Dans la démonstration du lemme \ref{lemma-inverse-limit-quasi-compact}, nous
avons vu que $X = \lim X_i$ est l'intersection des fermés
$$
Z_\varphi = \Gamma_\varphi \times \prod\nolimits_{l \not = j, k} X_l
$$
dans l'espace quasi-compact $\prod X_i$, où $\varphi : j \to k$ est un
morphisme de $\mathcal{I}$ et où $\Gamma_\varphi \subset X_j \times X_k$
est le graphe du morphisme correspondant $X_j \to X_k$. D'après le
lemme \ref{lemma-intersection-closed-in-quasi-compact}, il suffit donc de
montrer que toute intersection finie de ces parties est non vide.
Soit $\varphi_t : j_t \to k_t$, $t = 1, \ldots, n$, une famille finie de
morphismes de $\mathcal{I}$. Puisque $\mathcal{I}$ est cofiltrante, nous
pouvons choisir un objet $j$ et un morphisme $\psi_t : j \to j_t$ pour tout
$t$. Pour chaque paire $t, t'$ telle que (a) $j_t = j_{t'}$, ou
(b) $j_t = k_{t'}$, ou (c) $k_t = k_{t'}$, nous obtenons deux morphismes
$j \to l$, avec $l = j_t$ dans les cas (a) et (b), ou $l = k_t$ dans le cas
(c). Comme $\mathcal{I}$ est cofiltrante et qu'il n'y a qu'un nombre fini de
paires $(t, t')$, nous pouvons choisir une application $j' \to j$ qui égalise
ces deux morphismes pour toutes ces paires $(t, t')$. Choisissons un élément
$x \in X_{j'}$ et, pour tout $t$, notons $x_{j_t}$, respectivement
$x_{k_t}$, l'image de $x$ par le morphisme $X_{j'} \to X_j \to X_{j_t}$,
respectivement $X_{j'} \to X_j \to X_{j_t} \to X_{k_t}$.
Pour tout indice $l \in \Ob(\mathcal{I})$ qui n'est égal ni à $j_t$ ni à
$k_t$ pour aucun $t$, choisissons un élément arbitraire $x_l \in X_l$ (en
utilisant l'axiome du choix). Alors $(x_i)_{i \in \Ob(\mathcal{I})}$ appartient, par
construction, à l'intersection
$$
Z_{\varphi_1} \cap \ldots \cap Z_{\varphi_n}
$$
et la démonstration est achevée.
\end{proof}


\section{Parties constructibles}
\label{section-constructible}

\begin{definition}
\label{definition-constructible}
Soit $X$ un espace topologique. Soit $E \subset X$ une partie de $X$.
\begin{enumerate}
\item On dit que $E$ est {\it constructible}\footnote{Dans la deuxième édition
d'EGA I \cite{EGA1-second}, cette partie était appelée « globalement
constructible » et la terminologie « constructible » était employée pour ce
que nous appelons une partie localement constructible.}
dans $X$ si $E$ est une réunion finie
de parties de la forme $U \cap V^c$, où $U, V \subset X$ sont des ouverts
rétrocompacts dans $X$.
\item On dit que $E$ est {\it localement constructible} dans $X$ s'il existe un
recouvrement ouvert $X = \bigcup V_i$ tel que chaque $E \cap V_i$ soit
constructible dans $V_i$.
\end{enumerate}
\end{definition}

\begin{lemma}
\label{lemma-constructible}
La collection des parties constructibles est stable par
intersections finies, unions finies et passage au complémentaire.
\end{lemma}

\begin{proof}
Remarquons que si $U_1$, $U_2$ sont des ouverts rétrocompacts dans $X$,
alors $U_1 \cup U_2$ l'est aussi, car la réunion de deux parties
quasi-compactes de $X$ est quasi-compacte. Il est également vrai que
$U_1 \cap U_2$ est rétrocompact. En effet, supposons que $U \subset X$
soit un ouvert quasi-compact ; alors $U_2 \cap U$ est quasi-compact puisque
$U_2$ est rétrocompact dans $X$, et nous en déduisons que
$U_1 \cap (U_2 \cap U)$ est quasi-compact puisque $U_1$ est
rétrocompact dans $X$. À partir de là, il est formel de montrer que
le complémentaire d'une partie constructible est constructible,
que les unions finies de parties constructibles sont constructibles et
que les intersections finies de parties constructibles sont constructibles.
\end{proof}

\begin{lemma}
\label{lemma-inverse-images-constructibles}
Soit $f : X \to Y$ une application continue d'espaces topologiques.
Si l'image réciproque de tout ouvert rétrocompact de $Y$
est rétrocompacte dans $X$, alors les images réciproques des parties
constructibles sont constructibles.
\end{lemma}

\begin{proof}
Cette assertion résulte de $f^{-1}(U \cap V^c) = f^{-1}(U) \cap f^{-1}(V)^c$
et de la définition des parties constructibles.
\end{proof}

\begin{lemma}
\label{lemma-open-immersion-constructible-inverse-image}
Soit $U \subset X$ un ouvert. Pour une partie constructible
$E \subset X$, l'intersection $E \cap U$ est constructible
dans $U$.
\end{lemma}

\begin{proof}
Supposons que $V \subset X$ soit un ouvert rétrocompact dans $X$.
Il suffit de montrer que $V \cap U$ est rétrocompact dans $U$
d'après le lemme \ref{lemma-inverse-images-constructibles}. Pour cela,
soit $W \subset U$ un ouvert quasi-compact. Alors $W$
est ouvert et quasi-compact dans $X$. Par conséquent, $V \cap W = V \cap U \cap W$
est quasi-compact puisque $V$ est rétrocompact dans $X$.
\end{proof}

\begin{lemma}
\label{lemma-quasi-compact-open-immersion-constructible-image}
Soit $U \subset X$ un ouvert rétrocompact. Soit $E \subset U$.
Si $E$ est constructible dans $U$, alors $E$ est constructible dans $X$.
\end{lemma}

\begin{proof}
Supposons que $V, W \subset U$ soient des ouverts rétrocompacts dans $U$.
Alors $V, W$ sont des ouverts rétrocompacts dans $X$
(lemme \ref{lemma-composition-quasi-compact}).
Ainsi, $V \cap (U \setminus W) = V \cap (X \setminus W)$
est constructible dans $X$. On conclut puisque toute partie constructible de $U$
est une réunion finie de parties de la forme $V \cap (U \setminus W)$.
\end{proof}

\begin{lemma}
\label{lemma-collate-constructible}
Soit $X$ un espace topologique. Soit $E \subset X$ une partie.
Soit $X = V_1 \cup \ldots \cup V_m$ un recouvrement fini par des
ouverts rétrocompacts.
Alors $E$ est constructible dans $X$ si et seulement si $E \cap V_j$
est constructible dans $V_j$ pour chaque $j = 1, \ldots, m$.
\end{lemma}

\begin{proof}
Si $E$ est constructible dans $X$, alors le
lemme \ref{lemma-open-immersion-constructible-inverse-image}
montre que $E \cap V_j$ est constructible dans $V_j$ pour tout $j$.
Réciproquement, supposons que $E \cap V_j$
soit constructible dans $V_j$ pour chaque $j = 1, \ldots, m$.
Alors $E = \bigcup E \cap V_j$ est une réunion finie de
parties constructibles d'après le
lemme \ref{lemma-quasi-compact-open-immersion-constructible-image},
et est donc constructible.
\end{proof}

\begin{lemma}
\label{lemma-intersect-constructible-with-closed}
Soit $X$ un espace topologique. Soit $Z \subset X$ une partie fermée
telle que $X \setminus Z$ soit quasi-compact.
Alors, pour toute partie constructible $E \subset X$, l'intersection
$E \cap Z$ est constructible dans $Z$.
\end{lemma}

\begin{proof}
Supposons que $V \subset X$ soit un ouvert rétrocompact dans $X$.
Il suffit de montrer que $V \cap Z$ est rétrocompact dans $Z$
d'après le lemme \ref{lemma-inverse-images-constructibles}. Pour cela,
soit $W \subset Z$ un ouvert quasi-compact. La partie
$W' = W \cup (X \setminus Z)$ est quasi-compacte, ouverte, et $W = Z \cap W'$.
Par conséquent, $V \cap Z \cap W = V \cap Z \cap W'$
est une partie fermée de l'ouvert quasi-compact $V \cap W'$
puisque $V$ est rétrocompact dans $X$. Ainsi, $V \cap Z \cap W$ est quasi-compact
d'après le lemme \ref{lemma-closed-in-quasi-compact}.
\end{proof}

\begin{lemma}
\label{lemma-intersect-constructible-with-retrocompact}
Soit $X$ un espace topologique. Soit $T \subset X$ une partie. Supposons que
\begin{enumerate}
\item $T$ soit rétrocompacte dans $X$,
\item les ouverts quasi-compacts forment une base de la topologie de $X$.
\end{enumerate}
Alors, pour toute partie constructible $E \subset X$, l'intersection $E \cap T$ est
constructible dans $T$.
\end{lemma}

\begin{proof}
Supposons que $V \subset X$ soit un ouvert rétrocompact dans $X$.
Il suffit de montrer que $V \cap T$ est rétrocompact dans $T$
d'après le lemme \ref{lemma-inverse-images-constructibles}. Pour cela,
soit $W \subset T$ un ouvert quasi-compact. Par l'hypothèse (2),
nous pouvons trouver un ouvert quasi-compact $W' \subset X$
tel que $W = T \cap W'$ (détails omis).
Ainsi, $V \cap T \cap W = V \cap T \cap W'$
est l'intersection de $T$ avec  l'ouvert quasi-compact $V \cap W'$
puisque $V$ est rétrocompact dans $X$. Ainsi, $V \cap T \cap W$ est quasi-compact.
\end{proof}

\begin{lemma}
\label{lemma-closed-constructible-image}
Soit $Z \subset X$ une partie fermée dont le complémentaire est un ouvert rétrocompact.
Soit $E \subset Z$. Si $E$ est constructible dans $Z$, alors $E$ est constructible
dans $X$.
\end{lemma}

\begin{proof}
Supposons que $V \subset Z$ soit un ouvert rétrocompact dans $Z$. Considérons l'ouvert
$\tilde V = V \cup (X \setminus Z)$ de $X$. Soit $W \subset X$ un
ouvert quasi-compact. Alors
$$
W \cap \tilde V =
\left(V \cap W\right) \cup \left((X \setminus Z) \cap W\right).
$$
La première partie est quasi-compacte, car $V \cap W = V \cap (Z \cap W)$ et
$(Z \cap W)$ est un ouvert quasi-compact dans $Z$
(lemme \ref{lemma-closed-in-quasi-compact}), et $V$ est rétrocompact dans $Z$.
La seconde partie est quasi-compacte puisque $(X \setminus Z)$ est rétrocompact dans $X$.
Nous voyons ainsi que $\tilde V$ est rétrocompact dans $X$.
Ainsi, si $V_1, V_2 \subset Z$ sont des ouverts rétrocompacts, alors
$$
V_1 \cap (Z \setminus V_2) = \tilde V_1 \cap (X \setminus \tilde V_2)
$$
est constructible dans $X$. On conclut puisque toute partie constructible de $Z$
est une réunion finie de parties de la forme $V_1 \cap (Z \setminus V_2)$.
\end{proof}

\begin{lemma}
\label{lemma-constructible-is-retrocompact}
Soit $X$ un espace topologique. Toute partie constructible
de $X$ est rétrocompacte.
\end{lemma}

\begin{proof}
Soit $E = \bigcup_{i = 1, \ldots, n} U_i \cap V_i^c$, où $U_i, V_i$ sont des
ouverts rétrocompacts dans $X$. Soit $W \subset X$ un ouvert quasi-compact.
Alors $E \cap W = \bigcup_{i = 1, \ldots, n} U_i \cap V_i^c \cap W$.
Il suffit donc de montrer que $U \cap V^c \cap W$ est quasi-compact
si $U, V$ sont des ouverts rétrocompacts et $W$ un ouvert
quasi-compact. Cela est vrai parce que $U \cap V^c \cap W$ est une partie fermée
de la partie quasi-compacte $U \cap W$, si bien que le
lemme \ref{lemma-closed-in-quasi-compact}
s'applique.
\end{proof}

\noindent
Question : le lemme suivant reste-t-il valable si l'on suppose que $X$ est un
espace topologique quasi-compact ? Comparer avec le
lemme \ref{lemma-intersect-constructible-with-closed}.

\begin{lemma}
\label{lemma-intersect-constructible-with-constructible}
Soit $X$ un espace topologique. Supposons que
$X$ possède une base constituée d'ouverts quasi-compacts.
Pour des parties $E, E'$ constructibles dans $X$, l'intersection
$E \cap E'$ est constructible dans $E$.
\end{lemma}

\begin{proof}
Combiner les lemmes \ref{lemma-intersect-constructible-with-retrocompact} et
\ref{lemma-constructible-is-retrocompact}.
\end{proof}

\begin{lemma}
\label{lemma-constructible-in-constructible}
Soit $X$ un espace topologique. Supposons que
$X$ possède une base constituée d'ouverts quasi-compacts.
Soit $E$ une partie constructible dans $X$ et $F \subset E$ une partie constructible dans $E$.
Alors $F$ est constructible dans $X$.
\end{lemma}

\begin{proof}
Observons que toute partie rétrocompacte $T$ de $X$ possède, pour la topologie
induite, une base constituée d'ouverts quasi-compacts. Cela vaut en particulier
pour toute partie constructible
(lemme \ref{lemma-constructible-is-retrocompact}).
Écrivons $E = E_1 \cup \ldots \cup E_n$, avec $E_i = U_i \cap V_i^c$,
où $U_i, V_i \subset X$ sont des ouverts rétrocompacts.
Notons que $E_i = E \cap E_i$ est constructible dans $E$ d'après le
lemme \ref{lemma-intersect-constructible-with-constructible}.
Ainsi, $F \cap E_i$ est constructible dans $E_i$ d'après le
lemme \ref{lemma-intersect-constructible-with-constructible}.
Il suffit donc de démontrer le lemme dans le cas où $E = U \cap V^c$,
où $U, V \subset X$ sont des ouverts rétrocompacts.
Dans ce cas, l'inclusion $E \subset X$ est la composée
$$
E = U \cap V^c \to U \to X
$$
On peut alors appliquer le lemme \ref{lemma-closed-constructible-image}
à la première inclusion et le
lemme \ref{lemma-quasi-compact-open-immersion-constructible-image}
à la seconde.
\end{proof}

\begin{lemma}
\label{lemma-locally-closed-constructible-image}
Soit $X$ un espace topologique quasi-compact possédant une base constituée
d'ouverts quasi-compacts telle que l'intersection de deux
ouverts quasi-compacts quelconques soit quasi-compacte.
Soit $T \subset X$ une partie localement fermée
telle que $T$ soit quasi-compacte et que $T^c$ soit rétrocompacte dans $X$.
Alors $T$ est constructible dans $X$.
\end{lemma}

\begin{proof}
Notons que $T$ est quasi-compacte et ouverte dans $\overline{T}$.
En utilisant notre base d'ouverts quasi-compacts, nous pouvons écrire
$T = U \cap \overline{T}$, où $U$ est un ouvert quasi-compact de $X$.
Alors $V = U \setminus T = U \cap T^c$ est rétrocompact dans $U$, puisque $T^c$
est rétrocompact dans $X$. Par conséquent, $V$ est quasi-compact. Puisque
l'intersection de deux ouverts quasi-compacts quelconques est quasi-compacte,
tout ouvert quasi-compact de $X$ est rétrocompact. Ainsi, $T = U \cap V^c$,
où $U$ et $V = U \setminus T$ sont des ouverts rétrocompacts de $X$.
A fortiori, $T$ est constructible dans $X$.
\end{proof}

\begin{lemma}
\label{lemma-collate-constructible-from-constructible}
Soit $X$ un espace topologique dont la topologie possède une base
constituée d'ouverts quasi-compacts. Soit $E \subset X$ une partie.
Soit $X = E_1 \cup \ldots \cup E_m$ un recouvrement fini par des parties
constructibles. Alors $E$ est constructible dans $X$ si et seulement si $E \cap E_j$
est constructible dans $E_j$ pour chaque $j = 1, \ldots, m$.
\end{lemma}

\begin{proof}
Combiner les
lemmes \ref{lemma-intersect-constructible-with-constructible} et
\ref{lemma-constructible-in-constructible}.
\end{proof}

\begin{lemma}
\label{lemma-generic-point-in-constructible}
Soit $X$ un espace topologique. Supposons que
$Z \subset X$ soit irréductible. Soit $E \subset X$
une réunion finie de parties localement fermées (par exemple,\ $E$
est constructible). Les assertions suivantes sont équivalentes :
\begin{enumerate}
\item L'intersection $E \cap Z$ contient une partie ouverte
dense de $Z$.
\item L'intersection $E \cap Z$ est dense dans $Z$.
\end{enumerate}
Si $Z$ possède un point générique $\xi$, elles sont
également équivalentes à
\begin{enumerate}
\item[(3)] On a $\xi \in E$.
\end{enumerate}
\end{lemma}

\begin{proof}
L'implication (1) $\Rightarrow$ (2) est claire. Supposons (2).
Notons que $E \cap Z$ est une réunion finie de parties localement fermées
$Z_i$ de $Z$. Puisque $Z$ est irréductible, l'une des $Z_i$ doit
être dense dans $Z$. Cette partie $Z_i$ est alors ouverte et dense dans $Z$, car elle est
ouverte dans son adhérence. Par conséquent, (1) est vérifiée.

\medskip\noindent
Supposons que $\xi \in Z$ soit un point générique. Si les conditions équivalentes
(1) et (2) sont vérifiées, alors $\xi \in E$. Réciproquement, si $\xi \in E$, alors
$\xi \in E \cap Z$ et, par conséquent, $E \cap Z$ est dense dans $Z$.
\end{proof}






 

\section{Parties constructibles et espaces noethériens}
\label{section-constructible-Noetherian}

\begin{lemma}
\label{lemma-constructible-Noetherian-space}
Soit $X$ un espace topologique noethérien.
Les parties constructibles de $X$ sont exactement les réunions finies de
parties localement fermées de $X$.
\end{lemma}

\begin{proof}
Cela résulte immédiatement du lemme \ref{lemma-Noetherian-quasi-compact}.
\end{proof}

\begin{lemma}
\label{lemma-constructible-map-Noetherian}
Soit $f : X \to Y$ une application continue entre espaces topologiques
noethériens. Si $E \subset Y$ est constructible dans $Y$, alors $f^{-1}(E)$
est constructible dans $X$.
\end{lemma}

\begin{proof}
Cela résulte immédiatement du
lemme \ref{lemma-constructible-Noetherian-space} et de la définition d'une
application continue.
\end{proof}

\begin{lemma}
\label{lemma-characterize-constructible-Noetherian}
Soit $X$ un espace topologique noethérien.
Soit $E \subset X$ une partie.
Les propriétés suivantes sont équivalentes :
\begin{enumerate}
\item $E$ est constructible dans $X$, et
\item pour tout fermé irréductible $Z \subset X$, l'intersection $E \cap Z$
contient un ouvert non vide de $Z$ ou n'est pas dense dans $Z$.
\end{enumerate}
\end{lemma}

\begin{proof}
Supposons $E$ constructible et $Z \subset X$ fermé irréductible.
Alors $E \cap Z$ est constructible dans $Z$ d'après le
lemme \ref{lemma-constructible-map-Noetherian}.
Ainsi, $E \cap Z$ est une réunion finie de parties localement fermées non
vides $T_i$ de $Z$. Il est clair que, si aucun des $T_i$ n'est ouvert dans
$Z$, alors $E \cap Z$ n'est pas dense dans $Z$. On voit ainsi que (1)
implique (2).

\medskip\noindent
Réciproquement, supposons (2) vérifiée. Considérons l'ensemble $\mathcal{S}$
des fermés $Y$ de $X$ tels que $E \cap Y$ ne soit pas constructible dans $Y$.
Si $\mathcal{S} \not = \emptyset$, il admet un plus petit élément $Y$, car
$X$ est noethérien.
Soit $Y = Y_1 \cup \ldots \cup Y_r$ la décomposition de $Y$ en ses
composantes irréductibles ; voir le lemme \ref{lemma-Noetherian}.
Si $r > 1$, alors chaque $Y_i \cap E$ est constructible dans $Y_i$, donc est
une réunion finie de parties localement fermées de $Y_i$. Par conséquent,
$E \cap Y$ est lui aussi une réunion finie de parties localement fermées de
$Y$, et nous concluons que $E \cap Y$ est constructible dans $Y$ d'après le
lemme \ref{lemma-constructible-Noetherian-space}.
C'est une contradiction ; ainsi $r = 1$. Si $r = 1$, alors $Y$ est
irréductible et, par l'hypothèse (2), on voit que $E \cap Y$ se trouve dans
l'un des deux cas suivants : (a) il contient un ouvert non vide $V$ de $Y$ ; ou
(b) il n'est pas dense dans $Y$.
Dans le cas (a), la minimalité de $Y$ montre que $E \cap (Y \setminus V)$
est une réunion finie de parties localement fermées de $Y \setminus V$.
Ainsi, $E \cap Y$ est une réunion finie de parties localement fermées de $Y$
et est constructible d'après le
lemme \ref{lemma-constructible-Noetherian-space}.
C'est une contradiction ; nous sommes donc nécessairement dans le cas (b).
Dans le cas (b), on a $E \cap Y = E \cap Y'$ pour un fermé propre
$Y' \subset Y$. La minimalité de $Y$ montre que $E \cap Y'$ est une réunion
finie de parties localement fermées de $Y'$, puis que
$E \cap Y' = E \cap Y$ est une réunion finie de parties localement fermées
de $Y$ et est constructible d'après le
lemme \ref{lemma-constructible-Noetherian-space}.
Cette contradiction achève la démonstration du lemme.
\end{proof}

\begin{lemma}
\label{lemma-constructible-neighbourhood-Noetherian}
Soit $X$ un espace topologique noethérien.
Soit $x \in X$.
Soit $E \subset X$ constructible dans $X$.
Les propriétés suivantes sont équivalentes :
\begin{enumerate}
\item $E$ est un voisinage de $x$, et
\item pour tout fermé irréductible $Y$ de $X$ qui contient $x$,
l'intersection $E \cap Y$ est dense dans $Y$.
\end{enumerate}
\end{lemma}

\begin{proof}
Il est clair que (1) implique (2). Supposons (2).
Considérons l'ensemble $\mathcal{S}$ des fermés $Y$ de $X$ qui contiennent
$x$ et tels que $E \cap Y$ ne soit pas un voisinage de $x$ dans $Y$.
Si $\mathcal{S} \not = \emptyset$, il admet un élément minimal $Y$, car $X$
est noethérien. Supposons $Y = Y_1 \cup Y_2$, où $Y_1$, $Y_2$ sont deux
fermés non vides strictement plus petits. Si $x \in Y_i$ pour $i = 1, 2$,
alors $Y_i \cap E$ est un voisinage de $x$ dans $Y_i$, et nous concluons que
$Y \cap E$ est un voisinage de $x$ dans $Y$, contradiction. Si $x \in Y_1$
mais $x \not\in Y_2$ (disons), alors $Y_1 \cap E$ est un voisinage de $x$
dans $Y_1$, donc aussi dans $Y$, ce qui est encore une contradiction.
Nous concluons que $Y$ est un fermé irréductible. Par l'hypothèse (2),
$E \cap Y$ est dense dans $Y$. Ainsi, $E \cap Y$ contient un ouvert $V$ de
$Y$ ; voir le lemme \ref{lemma-characterize-constructible-Noetherian}.
Si $x \in V$, alors $E \cap Y$ est un voisinage de $x$ dans $Y$, ce qui est
une contradiction. Si $x \not \in V$, alors $Y' = Y \setminus V$ est un
fermé propre de $Y$ qui contient $x$. La minimalité de $Y$ montre que
$E \cap Y'$ contient un voisinage ouvert $V' \subset Y'$ de $x$ dans $Y'$.
Mais alors $V' \cup V$ est un voisinage ouvert de $x$ dans $Y$ contenu dans
$E$, contradiction.
Cette contradiction achève la démonstration du lemme.
\end{proof}

\begin{lemma}
\label{lemma-characterize-open-Noetherian}
Soit $X$ un espace topologique noethérien.
Soit $E \subset X$ une partie.
Les propriétés suivantes sont équivalentes :
\begin{enumerate}
\item $E$ est ouvert dans $X$, et
\item pour tout fermé irréductible $Y$ de $X$, l'intersection $E \cap Y$ est
vide ou contient un ouvert non vide de $Y$.
\end{enumerate}
\end{lemma}

\begin{proof}
Cela résulte formellement des
lemmes \ref{lemma-characterize-constructible-Noetherian} et
\ref{lemma-constructible-neighbourhood-Noetherian}.
\end{proof}





\section{Caractérisation des applications propres}
\label{section-proper}

\noindent
Nous consacrons une section à la notion d'application propre en topologie
usuelle. Nous définissons une application continue d'espaces topologiques
comme étant {\it propre} si elle est universellement fermée et séparée.
Bien que cette définition s'accorde avec celle d'un morphisme propre en
géométrie algébrique, elle diffère de la définition de Bourbaki.
Avec notre définition d'une application propre d'espaces topologiques, le
théorème de changement de base propre (Cohomologie,
théorème \ref{cohomology-theorem-proper-base-change}) est valable sans aucune
hypothèse supplémentaire. De plus, pour un morphisme $f : X \to Y$ de schémas
de type fini sur $\mathbf{C}$, on a l'équivalence suivante : $f$ est propre
comme morphisme de schémas si et seulement si l'application continue
$f : X(\mathbf{C}) \to Y(\mathbf{C})$ sur les points à valeurs dans
$\mathbf{C}$, munis de la topologie classique, est propre.
Ceci est expliqué dans \cite[Exp. XII, Prop.\ 3.2(v)]{SGA1}, où une note
signale aussi que la notion topologique de propreté employée est celle de
Bourbaki, à laquelle on ajoute l'hypothèse de séparation.

\medskip\noindent
Il nous paraît utile de nommer trois notions distinctes : être séparée, être
universellement fermée, et posséder ces deux propriétés à la fois
(c'est-à-dire être propre). Pour une application continue $f : X \to Y$
entre espaces localement compacts et séparés, le mot « propre » désigne depuis
longtemps la propriété « $f^{-1}($compact$) =$ compact », laquelle équivaut à
être universellement fermée pour des espaces aussi réguliers.
En fait, nous verrons que, dans le cas général, la condition sur les images
réciproques, formulée par souci de clarté à l'aide du mot « quasi-compact »,
équivaut à être universellement fermée si l'on suppose en outre l'application
fermée. Voir aussi \cite[exercices 22 à 26 du chapitre II]{LangReal}, mais
attention : Lang emploie « propre » comme synonyme d'« universellement
fermée », à l'instar de Bourbaki.

\begin{lemma}[Lemme du tube]
\label{lemma-tube}
Soient $X$ et $Y$ des espaces topologiques.
Soient $A \subset X$ et $B \subset Y$ des parties quasi-compactes.
Soit $A \times B \subset W \subset X \times Y$, où $W$ est ouvert dans
$X \times Y$. Il existe alors des ouverts $A \subset U \subset X$ et
$B \subset V \subset Y$ tels que $U \times V \subset W$.
\end{lemma}

\begin{proof}
Pour tous $a \in A$ et $b \in B$, il existe des ouverts $U_{(a, b)}$ de $X$
et $V_{(a, b)}$ de $Y$ tels que
$(a, b) \in U_{(a, b)} \times V_{(a, b)} \subset W$.
Fixons $b$ ; il existe alors un nombre fini de points $a_1, \ldots, a_n$ tels
que $A \subset U_{(a_1, b)} \cup \ldots \cup U_{(a_n, b)}$.
Par conséquent,
$$
A \times \{b\} \subset
(U_{(a_1, b)} \cup \ldots \cup U_{(a_n, b)}) \times
(V_{(a_1, b)} \cap \ldots \cap V_{(a_n, b)}) \subset W.
$$
Ainsi, pour tout $b \in B$, il existe des ouverts $U_b \subset X$ et
$V_b \subset Y$ tels que $A \times \{b\} \subset U_b \times V_b \subset W$.
Comme précédemment, il existe un nombre fini de points $b_1, \ldots, b_m$
tels que $B \subset V_{b_1} \cup \ldots \cup V_{b_m}$.
Le résultat en découle, car
$A \times B \subset
(U_{b_1} \cap \ldots \cap U_{b_m}) \times
(V_{b_1} \cup \ldots \cup V_{b_m})$.
\end{proof}

\noindent
La terminologie de la définition suivante peut différer légèrement de celle
dont le lecteur a l'habitude.

\begin{definition}
\label{definition-proper-map}
Soit $f : X\to Y$ une application continue entre espaces topologiques.
\begin{enumerate}
\item On dit que l'application $f$ est {\it fermée} si l'image de toute partie
fermée est fermée.
\item On dit que l'application $f$ est {\it propre au sens de
Bourbaki}\footnote{C'est la terminologie employée dans \cite{Bourbaki}.
Cette propriété est parfois appelée « universellement fermée » dans la
littérature.} si l'application $Z \times X\to Z \times Y$ est fermée pour tout
espace topologique $Z$.
\item On dit que l'application $f$ est {\it quasi-propre} si l'image réciproque
$f^{-1}(V)$ de toute partie quasi-compacte $V \subset Y$ est quasi-compacte.
\item On dit que $f$ est {\it universellement fermée} si l'application
$f': Z \times_Y X \to Z$ est fermée pour toute application continue
$g: Z \to Y$.
\item On dit que $f$ est {\it propre} si $f$ est séparée et universellement
fermée.
\end{enumerate}
\end{definition}

\noindent
Le lemme suivant sera utile plus loin.

\begin{lemma}
\label{lemma-characterize-quasi-compact}
\begin{reference}
Combinaison de
\cite[I, p. 75, Lemme 1]{Bourbaki} et de
\cite[I, p. 76, Corollaire 1]{Bourbaki}.
\end{reference}
Un espace topologique $X$ est quasi-compact si et seulement si l'application
de projection $Z \times X \to Z$ est fermée pour tout espace topologique $Z$.
\end{lemma}

\begin{proof}
(Voir aussi la remarque ci-dessous.)
Si $X$ n'est pas quasi-compact, il existe un recouvrement ouvert
$X = \bigcup_{i \in I} U_i$ tel qu'aucun nombre fini des $U_i$ ne recouvre
$X$.
Soit $Z$ la partie de l'ensemble des parties $\mathcal{P}(I)$ de $I$
constituée de $I$ et de toutes les parties finies non vides de $I$.
Définissons sur $Z$ une topologie ayant pour base les ensembles suivants :
\begin{enumerate}
\item Toutes les parties de $Z\setminus\{I\}$.
\item Pour toute partie finie $K$ de $I$, l'ensemble
$U_K := \{J\subset I \mid J \in Z, \ K\subset J \})$.
\end{enumerate}
Nous laissons au lecteur le soin de vérifier qu'il s'agit bien d'une base de
topologie. Considérons la partie de $Z \times X$ définie par la formule
$$
M = \{(J, x) \mid J \in Z, \ x \in \bigcap\nolimits_{i \in J} U_i^c)\}
$$
Si $(J, x) \not \in M$, alors $x \in U_i$ pour un certain $i \in J$.
Ainsi, $U_{\{i\}} \times U_i \subset Z \times X$ est un ouvert contenant
$(J, x)$ et ne rencontrant pas $M$. Par conséquent, $M$ est fermé. La
projection de $M$ sur $Z$ est $Z-\{I\}$, qui n'est pas fermée. L'application
$Z \times X \to Z$ n'est donc pas fermée.

\medskip\noindent
Supposons $X$ quasi-compact. Soit $Z$ un espace topologique.
Soit $M \subset  Z \times X$ fermé. Soit $z \in Z$ un point qui n'appartient
pas à $\text{pr}_1(M)$. D'après le lemme du tube \ref{lemma-tube}, il existe un
ouvert $U \subset Z$ tel que $U \times X$ soit contenu dans le complémentaire
de $M$. Ainsi, $\text{pr}_1(M)$ est fermé.
\end{proof}

\begin{remark}
\label{remark-lemma-literature}
Le lemme \ref{lemma-characterize-quasi-compact} combine
\cite[I, p. 75, Lemme 1]{Bourbaki} et
\cite[I, p. 76, Corollaire 1]{Bourbaki}.
\end{remark}

\begin{theorem}
\label{theorem-characterize-proper}
\begin{reference}
Dans \cite[I, p. 75, théorème 1]{Bourbaki}, on trouve :
(2) $\Leftrightarrow$ (4).
Dans \cite[I, p. 77, Proposition 6]{Bourbaki}, on trouve :
(2) $\Rightarrow$ (1).
\end{reference}
Soit $f: X\to Y$ une application continue entre espaces topologiques.
Les conditions suivantes sont équivalentes :
\begin{enumerate}
\item L'application $f$ est quasi-propre et fermée.
\item L'application $f$ est propre au sens de Bourbaki.
\item L'application $f$ est universellement fermée.
\item L'application $f$ est fermée et $f^{-1}(y)$ est quasi-compact pour tout
$y\in Y$.
\end{enumerate}
\end{theorem}

\begin{proof}
(Voir aussi la remarque ci-dessous.)
Si l'application $f$ vérifie (1), elle vérifie automatiquement (4), car tout
singleton est quasi-compact.

\medskip\noindent
Supposons que l'application $f$ vérifie (4).
Nous allons montrer qu'elle est universellement fermée, c'est-à-dire que (3)
est vérifiée. Soit $g : Z \to Y$ une application continue d'espaces
topologiques et considérons le diagramme
$$
\xymatrix{
Z \times_Y X \ar[r]_{g'} \ar[d]_{f'} & X \ar[d]^f \\
Z \ar[r]^g & Y
}
$$
Dans la démonstration, nous utiliserons le fait que
$Z \times_Y X \to Z \times X$ est un homéomorphisme sur son image,
c'est-à-dire que nous pouvons identifier $Z \times_Y X$ à la partie
correspondante de $Z \times X$, munie de la topologie induite.
L'image de $f' : Z \times_Y X \to Z$ est
$\Im(f') = \{z : g(z) \in f(X)\}$.
Comme $f(X)$ est fermé, on voit que $\Im(f')$ est un sous-espace fermé de $Z$.
Considérons un fermé $P \subset Z \times_Y X$.
Soit $z \in Z$, $z \not \in f'(P)$.
Si $z \not \in \Im(f')$, alors $Z \setminus \Im(f')$ est un voisinage
ouvert qui ne rencontre pas $f'(P)$.
Si $z$ appartient à $\Im(f')$, alors
$(f')^{-1}\{z\} = \{z\} \times f^{-1}\{g(z)\}$, et
$f^{-1}\{g(z)\}$ est quasi-compact par hypothèse. Comme $P$ est un fermé de
$Z \times_Y X$, il existe un fermé $P'$ de $Z \times X$ tel que
$P = P' \cap Z \times_Y X$.
Puisque $(f')^{-1}\{z\}$ est contenu dans
$P^c = P'^c \cup (Z \times_Y X)^c$ et puisque $(f')^{-1}\{z\}$ est disjoint
de $(Z \times_Y X)^c$, on voit que $(f')^{-1}\{z\}$ est contenu dans $P'^c$.
Nous pouvons appliquer le lemme du tube \ref{lemma-tube} à
$(f')^{-1}\{z\} = \{z\} \times f^{-1}\{g(z)\}
\subset (P')^c \subset Z \times X$.
Nous obtenons ainsi $V \times U$ contenant $(f')^{-1}\{z\}$, où $U$ et $V$
sont des ouverts respectivement de $X$ et de $Z$, et tel que
$V \times U$ ne rencontre pas $P'$. Alors l'ensemble
$V \cap g^{-1}(Y-f(U^c))$ est ouvert dans $Z$, puisque $f$ est fermée ; il
contient $z$ et ne rencontre pas l'image de $P$. Ainsi, $f'(P)$ est fermé.
Autrement dit, l'application $f$ est universellement fermée.

\medskip\noindent
L'implication (3) $\Rightarrow$ (2) est immédiate.
En effet, étant donné un espace topologique $Z$, considérons la projection
$g : Z \times Y \to Y$. Il est alors facile de voir que $f'$ est
l'application $Z \times X \to Z \times Y$, autrement dit que
$(Z \times Y) \times_Y X = Z \times X$. (Cette identification est une
propriété purement catégorique, qui n'a rien de spécifiquement topologique.)

\medskip\noindent
Supposons que $f$ vérifie (2). Nous allons montrer qu'elle vérifie (1).
Notons que $f$ est fermée, car $f$ s'identifie à l'application
$\{pt\} \times X \to \{pt\} \times Y$, qui est fermée par hypothèse.
Choisissons une partie quasi-compacte quelconque $K \subset Y$.
Soit $Z$ un espace topologique quelconque.
Comme $Z \times X \to Z \times Y$ est fermée, on voit que l'application
$Z \times f^{-1}(K) \to Z \times K$ est fermée (si $T$ est fermé dans
$Z \times f^{-1}(K)$, écrivons
$T = Z \times f^{-1}(K) \cap T'$ pour un certain fermé
$T' \subset Z \times X$). Comme $K$ est quasi-compact,
$K \times Z\to Z$ est fermée d'après le
lemme \ref{lemma-characterize-quasi-compact}.
Par conséquent, la composée $Z \times f^{-1}(K)\to Z \times K \to Z$ est
fermée ; le lemme \ref{lemma-characterize-quasi-compact} montre alors, une
nouvelle fois, que $f^{-1}(K)$ est quasi-compact.
\end{proof}

\begin{remark}
\label{remark-proof-literature}
Voici quelques références bibliographiques.
Dans \cite[I, p. 75, théorème 1]{Bourbaki}, on trouve :
(2) $\Leftrightarrow$ (4).
Dans \cite[I, p. 77, Proposition 6]{Bourbaki}, on trouve :
(2) $\Rightarrow$ (1).
Bien entendu, on a trivialement (1) $\Rightarrow$ (4).
Ainsi, (1), (2) et (4) sont équivalentes.
L'équivalence de (3) et (4) est \cite[chapitre II, exercice 25]{LangReal}.
\end{remark}

\begin{lemma}
\label{lemma-closed-map}
\begin{slogan}
Une application d'un espace compact vers un espace séparé est
universellement fermée.
\end{slogan}
Soit $f : X \to Y$ une application continue d'espaces topologiques.
Si $X$ est quasi-compact et $Y$ séparé, alors $f$ est universellement fermée.
\end{lemma}

\begin{proof}
Comme tout point de $Y$ est fermé, le
lemme \ref{lemma-closed-in-quasi-compact} montre que le fermé $f^{-1}(y)$ de
$X$ est quasi-compact pour tout $y \in Y$.
Ainsi, d'après le théorème \ref{theorem-characterize-proper}, il suffit de
montrer que $f$ est fermée.
Si $E \subset X$ est fermé, alors il est quasi-compact
(lemme \ref{lemma-closed-in-quasi-compact}), donc $f(E) \subset Y$ est
quasi-compact (lemme \ref{lemma-image-quasi-compact}), donc $f(E)$ est fermé
dans $Y$ (lemme \ref{lemma-quasi-compact-in-Hausdorff}).
\end{proof}

\begin{lemma}
\label{lemma-bijective-map}
Soit $f : X \to Y$ une application continue d'espaces topologiques.
Si $f$ est bijective, si $X$ est quasi-compact et si $Y$ est séparé, alors
$f$ est un homéomorphisme.
\end{lemma}

\begin{proof}
Il suffit de montrer que $f$ est fermée, car il en résulte que $f^{-1}$ est
continue. Si $T \subset X$ est fermé, alors $T$ est quasi-compact d'après le
lemme \ref{lemma-closed-in-quasi-compact}, donc $f(T)$ est quasi-compact
d'après le lemme \ref{lemma-image-quasi-compact}, donc $f(T)$ est fermé
d'après le lemme \ref{lemma-quasi-compact-in-Hausdorff}.
\end{proof}












\section{Espaces de Jacobson}
\label{section-space-jacobson}

\begin{definition}
\label{definition-space-jacobson}
Soit $X$ un espace topologique.
Soit $X_0$ l'ensemble des points fermés de $X$.
On dit que $X$ est {\it de Jacobson} si toute
partie fermée $Z \subset X$ est l'adhérence
de $Z \cap X_0$.
\end{definition}

\noindent
Notons qu'un espace topologique $X$ est de Jacobson si et seulement si
toute partie localement fermée non vide de $X$
possède un point fermé dans $X$.

\medskip\noindent
Soit $X$ un espace de Jacobson et soit $X_0$ l'ensemble
des points fermés de $X$, muni de la topologie induite.
Il est clair que la définition implique que le morphisme
$X_0 \to X$ induit une bijection entre les parties fermées
de $X_0$ et les parties fermées de $X$.
Ainsi, nombre de propriétés de $X$ passent à $X_0$.
Par exemple, les dimensions de Krull de $X$ et de $X_0$
sont égales.

\begin{lemma}
\label{lemma-jacobson-check-irreducible-closed}
Soit $X$ un espace topologique. Soit $X_0$ l'ensemble
des points fermés de $X$.
Supposons que, pour tout point $x\in X$,
l'intersection $X_0 \cap \overline{\{x\}}$ soit dense dans $\overline{\{x\}}$.
Alors $X$ est de Jacobson.
\end{lemma}

\begin{proof}
Soient $Z$ une partie fermée de $X$
et $U$ un ouvert de $X$
tels que $U\cap Z$ soit non vide.
Alors, pour $x\in U\cap Z$, l'ensemble $\overline{\{x\}}\cap U$ est une partie
non vide de $Z\cap U$
et, par hypothèse, il contient un point fermé dans $X$, comme voulu.
\end{proof}

\begin{lemma}
\label{lemma-non-jacobson-Noetherian-characterize}
Soit $X$ un espace topologique de Kolmogorov muni d'une base d'ouverts
quasi-compacts.
Si $X$ n'est pas de Jacobson, alors il existe un point non fermé
$x \in X$ tel que $\{x\}$ soit localement fermé.
\end{lemma}

\begin{proof}
Puisque $X$ n'est pas de Jacobson, il existe une partie fermée $Z$ et un ouvert $U$
de $X$ tels que $Z \cap U$ soit non vide et ne contienne aucun point fermé
dans $X$. Comme $X$ possède une base d'ouverts quasi-compacts, nous pouvons
remplacer $U$
par un voisinage ouvert quasi-compact d'un point de $Z\cap U$ et ainsi supposer
que $U$ est un ouvert quasi-compact. D'après le
lemme \ref{lemma-quasi-compact-closed-point}, il existe un point
$x \in Z \cap U$ fermé dans $Z \cap U$ ;
ainsi, $\{x\}$ est localement fermé mais n'est pas fermé dans $X$.
\end{proof}

\begin{lemma}
\label{lemma-jacobson-local}
Soit $X$ un espace topologique.
Soit $X = \bigcup U_i$ un recouvrement ouvert.
Alors $X$ est de Jacobson si et seulement si chaque $U_i$ est de Jacobson.
De plus, dans ce cas, $X_0 = \bigcup U_{i, 0}$.
\end{lemma}

\begin{proof}
Soit $X$ un espace topologique.
Soit $X_0$ l'ensemble des points fermés de $X$.
Soit $U_{i, 0}$ l'ensemble des points fermés de
$U_i$. Alors $X_0 \cap U_i \subset U_{i, 0}$,
mais l'égalité n'est pas toujours vérifiée.

\medskip\noindent
Supposons d'abord que chaque $U_i$ soit de Jacobson.
Nous affirmons qu'alors $X_0 \cap U_i = U_{i, 0}$.
En effet, supposons que $x \in U_{i, 0}$, c'est-à-dire que $x$ soit fermé dans
$U_i$. Soit $\overline{\{x\}}$ l'adhérence
dans $X$. Considérons $\overline{\{x\}} \cap U_j$.
Si $x \not \in U_j$, alors $\overline{\{x\}} \cap U_j = \emptyset$.
Si $x \in U_j$, alors $U_i \cap U_j \subset U_j$
est un ouvert de $U_j$ contenant $x$.
Soient $T' = U_j \setminus U_i \cap U_j$ et
$T = \{x\} \amalg T'$. Comme $x$ est fermé dans $U_i \cap U_j$
(puisque $x$ est fermé dans $U_i$), on voit que $T$ et $T'$
sont des parties fermées de $U_j$ et que $T$ contient
$x$. Comme $U_j$ est de Jacobson, les points fermés de
$U_j$ sont denses dans $T$. Puisque $T = \{x\} \amalg T'$,
cela n'est possible que si $x$ est fermé dans $U_j$.
Par conséquent, $\overline{\{x\}} \cap U_j = \{x\}$. Nous concluons
que $\overline{\{x\}} = \{ x \}$, comme voulu.

\medskip\noindent
Soit $Z \subset X$ une partie fermée (toujours
en supposant que chaque $U_i$ soit de Jacobson).
Puisque nous savons maintenant que $X_0 \cap Z  \cap U_i
= U_{i, 0} \cap Z$ est dense dans $Z \cap U_i$,
il s'ensuit immédiatement que $X_0 \cap Z$ est
dense dans $Z$.

\medskip\noindent
Réciproquement, supposons que $X$ soit de Jacobson.
Soit $Z \subset U_i$ une partie fermée. Alors
$X_0 \cap \overline{Z}$ est dense dans $\overline{Z}$.
Par conséquent, $X_0 \cap Z$ est également dense dans $Z$, car
$\overline{Z} \setminus Z$ est fermé. Comme $X_0 \cap U_i
\subset U_{i, 0}$, on voit que
$U_{i, 0} \cap Z$ est dense dans $Z$.
Ainsi, $U_i$ est de Jacobson, comme voulu.
\end{proof}

\begin{lemma}
\label{lemma-jacobson-inherited}
Soit $X$ un espace de Jacobson. Les types suivants de parties $T \subset X$
sont de Jacobson :
\begin{enumerate}
\item Les sous-espaces ouverts.
\item Les sous-espaces fermés.
\item Les sous-espaces localement fermés.
\item Les réunions de sous-espaces localement fermés.
\item Les parties constructibles.
\item Toute partie $T \subset X$ qui, localement sur $X$,
est une réunion de parties localement fermées.
\end{enumerate}
Dans chacun de ces cas, les points fermés de $T$ sont
fermés dans $X$.
\end{lemma}

\begin{proof}
Soit $X_0$ l'ensemble des points fermés de $X$. Pour toute partie
$T \subset X$, notons $(*)$ la propriété suivante :
\begin{itemize}
\item[(*)] Toute partie localement fermée non vide de $T$ possède un point
fermé dans $X$.
\end{itemize}
Notons que l'on a toujours $X_0 \cap T \subset T_0$. La propriété $(*)$
implique donc que $T$ est de Jacobson. En outre, elle implique clairement
que tout point fermé de $T$ est fermé dans $X$.

\medskip\noindent
Supposons que $T=\bigcup_i T_i$, où $T_i$ est localement fermé dans $X$.
Prenons une partie $A\subset T$ localement fermée et non vide dans $T$ ;
on peut alors choisir $T_i$ de telle sorte que $A\cap T_i$ soit non vide ; cette intersection est
localement fermée dans $T_i$, donc dans $X$.
Comme $X$ est de Jacobson, $A$ possède un point fermé dans $X$.
\end{proof}

\begin{lemma}
\label{lemma-finite-jacobson}
Un espace de Jacobson fini est discret.
Un espace de Jacobson qui ne possède qu'un nombre fini de points fermés est discret.
\end{lemma}

\begin{proof}
Si $X$ est fini, alors l'ensemble $X_0 \subset X$ de ses points fermés est fini.
Supposons que $X_0$ soit fini et que $X$ soit de Jacobson.
Alors $\overline{X_0} = X$ par la propriété de Jacobson.
Or $X_0 =\{x_1, \ldots, x_n\} = \bigcup_{i = 1, \ldots, n} \{x_i\}$
est une réunion finie de fermés, donc est fermé ; ainsi,
$X = \overline{X_0} = X_0$. Tout point est fermé et, par
finitude, tout point est ouvert.
\end{proof}

\begin{lemma}
\label{lemma-jacobson-equivalent-locally-closed}
\begin{slogan}
Pour les espaces de Jacobson, les points fermés déterminent entièrement la topologie.
\end{slogan}
Supposons que $X$ soit un espace topologique de Jacobson.
Soit $X_0$ l'ensemble des points fermés de $X$.
Il existe une correspondance bijective, qui préserve les inclusions,
$$
\{\text{réunions finies de parties loc. fermées de } X\}
\leftrightarrow
\{\text{réunions finies de parties loc. fermées de } X_0\}
$$
définie par $E \mapsto E \cap X_0$. Cette correspondance préserve
les parties localement fermées, les parties ouvertes et les parties fermées.
\end{lemma}

\begin{proof}
Nous nous contentons de démontrer que la correspondance $E \mapsto E \cap X_0$ est injective.
En effet, si $E\neq E'$, alors, sans perte de généralité, $E\setminus E'$ est
non vide, et c'est une réunion finie de parties localement fermées (détails omis).
Comme $X$ est de Jacobson, on voit que
$(E \setminus E') \cap X_0 = E \cap X_0 \setminus E' \cap X_0$ n'est pas vide.
\end{proof}

\begin{lemma}
\label{lemma-jacobson-equivalent-constructible}
Supposons que $X$ soit un espace topologique de Jacobson.
Soit $X_0$ l'ensemble des points fermés de $X$.
Il existe une correspondance bijective, qui préserve les inclusions,
$$
\{\text{parties constructibles de } X\}
\leftrightarrow
\{\text{parties constructibles de } X_0\}
$$
définie par $E \mapsto E \cap X_0$. Cette correspondance préserve
la sous-collection des ouverts rétrocompacts, ainsi que leurs complémentaires.
\end{lemma}

\begin{proof}
D'après le lemme \ref{lemma-jacobson-equivalent-locally-closed} ci-dessus,
il suffit de voir que, si $U$ est ouvert dans $X$, alors $U\cap X_0$ est
rétrocompact dans $X_0$ si et seulement si $U$ est rétrocompact dans $X$.
Pour cela, il suffit de démontrer que, si $U$ est ouvert dans $X$, alors $U\cap X_0$ est
quasi-compact si et seulement si $U$ est quasi-compact.
D'après le lemme \ref{lemma-jacobson-inherited}, nous pouvons remplacer
$X$ par $U$ et supposer que $U = X$.
Enfin, observons qu'une famille d'ouverts $\mathcal{U}$ de $X$ recouvre
$X$ si et seulement si elle recouvre $X_0$ : cela découle de la propriété
de Jacobson de $X$, appliquée au fermé $X\setminus \bigcup \mathcal{U}$ :
si ce fermé était non vide, il contiendrait un point de $X_0$.
\end{proof}



















\section{Spécialisation}
\label{section-specialization}

\begin{definition}
\label{definition-specialization}
Soit $X$ un espace topologique.
\begin{enumerate}
\item Si $x, x' \in X$, on dit que $x$ est une {\it spécialisation} de $x'$,
ou que $x'$ est une {\it générisation} de $x$, si
$x \in \overline{\{x'\}}$. Notation : $x' \leadsto x$.
\item Une partie $T \subset X$ est {\it stable par spécialisation} si, pour
tout $x' \in T$ et toute spécialisation $x' \leadsto x$, on a $x \in T$.
\item Une partie $T \subset X$ est {\it stable par générisation} si, pour tout
$x \in T$ et toute générisation $x' \leadsto x$, on a $x' \in T$.
\end{enumerate}
\end{definition}

\begin{lemma}
\label{lemma-open-closed-specialization}
Soit $X$ un espace topologique.
\begin{enumerate}
\item Toute partie fermée de $X$ est stable par spécialisation.
\item Toute partie ouverte de $X$ est stable par générisation.
\item Une partie $T \subset X$ est stable par spécialisation si et seulement
si son complémentaire $T^c$ est stable par générisation.
\end{enumerate}
\end{lemma}

\begin{proof}
Soit $F$ une partie fermée de $X$. Si $y\in F$, alors $\{y\} \subset F$, donc
$\overline{\{y\}} \subset \overline{F} = F$, puisque $F$ est fermé. Ainsi,
pour toute spécialisation $x$ de $y$, on a $x\in F$.

\medskip\noindent
Soient $x, y\in X$ tels que $x\in \overline{\{y\}}$, et soit $T$ une partie
de $X$. Dire que $T$ est stable par spécialisation signifie que $y\in T$
implique $x\in T$ ; réciproquement, dire que $T$ est stable par générisation
signifie que $x\in T$ implique $y\in T$. L'assertion (3) résulte donc de la
contraposition.

\medskip\noindent
La deuxième assertion découle de (1) et (3) en passant au complémentaire.
\end{proof}

\begin{lemma}
\label{lemma-stable-specialization}
Soit $T \subset X$ une partie d'un espace topologique $X$.
Les propriétés suivantes sont équivalentes :
\begin{enumerate}
\item $T$ est stable par spécialisation, et
\item $T$ est une réunion (filtrante) de parties fermées de $X$.
\end{enumerate}
\end{lemma}

\begin{proof}
Supposons $T$ stable par spécialisation. Alors, pour tout $y\in T$, on a
$\overline{\{y\}} \subset T$. Ainsi,
$T = \bigcup_{y\in T} \overline{\{y\}}$, qui est une réunion de parties
fermées de $X$. Réciproquement, supposons que
$T = \bigcup_{i\in I}F_i$, où les $F_i$ sont des parties fermées de $X$.
Si $y\in T$, il existe $i\in I$ tel que $y\in F_i$. Comme $F_i$ est fermé,
on a $\overline{\{y\}} \subset F_i \subset T$, ce qui prouve que $T$ est
stable par spécialisation.
\end{proof}

\begin{definition}
\label{definition-lift-specializations}
Soit $f : X \to Y$ une application continue d'espaces topologiques.
\begin{enumerate}
\item On dit que {\it les spécialisations se relèvent le long de $f$}, ou que
$f$ est {\it spécialisante}, si, étant donnés $y' \leadsto y$ dans $Y$ et
$x'\in X$ tel que $f(x') = y'$, il existe une spécialisation
$x' \leadsto x$ de $x'$ dans $X$ telle que $f(x) = y$.
\item On dit que {\it les générisations se relèvent le long de $f$}, ou que
$f$ est {\it générisante}, si, étant donnés $y' \leadsto y$ dans $Y$ et
$x\in X$ tel que $f(x) = y$, il existe une générisation $x' \leadsto x$ de
$x$ dans $X$ telle que $f(x') = y'$.
\end{enumerate}
\end{definition}

\begin{lemma}
\label{lemma-lift-specialization-composition}
Supposons que $f : X \to Y$ et $g : Y \to Z$ soient des applications
continues d'espaces topologiques. Si les spécialisations se relèvent le long
de $f$ et de $g$, elles se relèvent le long de $g \circ f$. Il en va de même
pour le relèvement des générisations.
\end{lemma}

\begin{proof}
Soit $z'\leadsto z$ une spécialisation dans $Z$ et soit $x' \in X$ tel que
$g\circ f (x') = z'$. Puisque les spécialisations se relèvent le long de $g$,
il existe une spécialisation $f(x') \leadsto y$ de $f(x')$ dans $Y$ telle que
$g(y) = z$. De même, puisqu'elles se relèvent le long de $f$, il existe une
spécialisation $x' \leadsto x$ de $x'$ dans $X$ telle que $f(x) = y$. Cela
fournit une spécialisation $x' \leadsto x$ de $x'$ dans $X$ telle que
$g\circ f(x) = z$. Autrement dit, les spécialisations se relèvent le long de
$g\circ f$.
\end{proof}

\begin{lemma}
\label{lemma-lift-specializations-images}
Soit $f : X \to Y$ une application continue d'espaces topologiques.
\begin{enumerate}
\item Si les spécialisations se relèvent le long de $f$ et si $T \subset X$
est stable par spécialisation, alors $f(T) \subset Y$ est stable par
spécialisation.
\item Si les générisations se relèvent le long de $f$ et si $T \subset X$ est
stable par générisation, alors $f(T) \subset Y$ est stable par générisation.
\end{enumerate}
\end{lemma}

\begin{proof}
Soit $y' \leadsto y$ une spécialisation dans $Y$, où $y'\in f(T)$, et soit
$x'\in T$ tel que $f(x') = y'$. Puisque les spécialisations se relèvent le
long de $f$, il existe une spécialisation $x'\leadsto x$ de $x'$ dans $X$
telle que $f(x) = y$. Or $T$ est stable par spécialisation, donc $x\in T$, et
alors $y \in f(T)$. Par conséquent, $f(T)$ est stable par spécialisation.

\medskip\noindent
La démonstration de (2) est identique, en utilisant le relèvement des
générisations le long de $f$.
\end{proof}

\begin{lemma}
\label{lemma-closed-open-map-specialization}
Soit $f : X \to Y$ une application continue d'espaces topologiques.
\begin{enumerate}
\item Si $f$ est fermée, alors les spécialisations se relèvent le long de $f$.
\item Si $f$ est ouverte, si $X$ est un espace topologique noethérien, si tout
fermé irréductible de $X$ possède un point générique et si $Y$ est un espace
de Kolmogorov, alors les générisations se relèvent le long de $f$.
\end{enumerate}
\end{lemma}

\begin{proof}
Supposons $f$ fermée. Soient $y' \leadsto y$ dans $Y$ et $x'\in X$ tel que
$f(x') = y'$. Considérons le fermé $T = \overline{\{x'\}}$ de $X$.
Alors $f(T) \subset Y$ est fermé et $y' \in f(T)$. Par conséquent,
$y \in f(T)$. Ainsi, $y = f(x)$ avec $x \in T$, c'est-à-dire
$x' \leadsto x$.

\medskip\noindent
Supposons $f$ ouverte, $X$ noethérien, tout fermé irréductible de $X$ pourvu
d'un point générique, et $Y$ un espace de Kolmogorov.
Soient $y' \leadsto y$ dans $Y$ et $x \in X$ tel que $f(x) = y$.
Considérons $T = f^{-1}(\{y'\}) \subset X$.
Prenons un voisinage ouvert $x \in U \subset X$ de $x$.
Alors $f(U) \subset Y$ est ouvert et $y \in f(U)$. Par conséquent,
$y' \in f(U)$. Autrement dit, $T \cap U \not = \emptyset$. Ceci prouve que
$x \in \overline{T}$. Puisque $X$ est noethérien, $T$ est noethérien
(lemme \ref{lemma-Noetherian}).
Il admet donc une décomposition $T = T_1 \cup \ldots \cup T_n$ en composantes
irréductibles. On a alors, de même,
$\overline{T} = \overline{T_1} \cup \ldots \cup \overline{T_n}$.
D'après ce qui précède, $x \in \overline{T_i}$ pour un certain $i$. Par
hypothèse, il existe un point générique $x' \in \overline{T_i}$, et l'on voit
que $x' \leadsto x$. Comme $x' \in \overline{T}$, on a
$f(x') \in \overline{\{y'\}}$. Remarquons que
$f(\overline{T_i}) = f(\overline{\{x'\}}) \subset \overline{\{f(x')\}}$.
Si $f(x') \not = y'$, alors, puisque $Y$ est un espace de Kolmogorov, $f(x')$ n'est pas
un point générique du fermé irréductible $\overline{\{y'\}}$, et l'inclusion
$\overline{\{f(x')\}} \subset \overline{\{y'\}}$ est stricte, c'est-à-dire
que $y' \not \in f(\overline{T_i})$.
Ceci contredit l'égalité $f(T_i) = \{y'\}$.
Ainsi, $f(x') = y'$, ce qui conclut.
\end{proof}

\begin{lemma}
\label{lemma-quotient-kolmogorov}
Supposons que $s, t : R \to U$ et $\pi : U \to X$ soient des applications
continues d'espaces topologiques telles que
\begin{enumerate}
\item $\pi$ est ouverte,
\item $U$ est sobre,
\item $s, t$ ont des fibres finies,
\item les générisations se relèvent le long de $s, t$,
\item $(t, s)(R) \subset U \times U$ est une relation d'équivalence sur $U$ et
$X$ est le quotient de $U$ par cette relation d'équivalence (comme ensemble).
\end{enumerate}
Alors $X$ est un espace de Kolmogorov.
\end{lemma}

\begin{proof}
Les propriétés (3) et (5) impliquent qu'un point $x$ correspond à une classe
d'équivalence finie $\{u_1, \ldots, u_n\} \subset U$. Supposons que
$x' \in X$ soit un second point correspondant à la classe d'équivalence
$\{u'_1, \ldots, u'_m\} \subset U$.
Supposons que $u_i \leadsto u'_j$ pour certains $i, j$. Alors, pour tout
$r' \in R$ tel que $s(r') = u'_j$, la propriété (4) fournit
$r \leadsto r'$ tel que $s(r) = u_i$. Par conséquent,
$t(r) \leadsto t(r')$. Puisque
$\{u'_1, \ldots, u'_m\} = t(s^{-1}(\{u'_j\}))$, nous concluons que tout
élément de $\{u'_1, \ldots, u'_m\}$ est une spécialisation d'un élément de
$\{u_1, \ldots, u_n\}$.
Ainsi, $\overline{\{u_1\}} \cup \ldots \cup \overline{\{u_n\}}$ est une
réunion de classes d'équivalence, donc est de la forme $\pi^{-1}(Z)$ pour une
partie $Z \subset X$. La propriété (1) montre que $Z$ est fermé dans $X$ et,
en fait, que $Z = \overline{\{x\}}$, car
$\pi(\overline{\{u_i\}}) \subset \overline{\{x\}}$ pour tout $i$.
Autrement dit, $x \leadsto x'$ si et seulement si un relèvement de $x$ dans
$U$ se spécialise en un relèvement de $x'$ dans $U$, si et seulement si tout
relèvement de $x'$ dans $U$ est une spécialisation d'un relèvement de $x$
dans $U$.

\medskip\noindent
Supposons que l'on ait à la fois $x \leadsto x'$ et $x' \leadsto x$.
Disons que $x$ correspond à $\{u_1, \ldots, u_n\}$ et $x'$ à
$\{u'_1, \ldots, u'_m\}$, comme ci-dessus. D'après les résultats du paragraphe
précédent, nous pouvons alors trouver une suite
$$
\ldots \leadsto u'_{j_3} \leadsto u_{i_3} \leadsto u'_{j_2} \leadsto
u_{i_2} \leadsto u'_{j_1} \leadsto u_{i_1}
$$
qui doit comporter une répétition ; la propriété (2) donne donc
$\{u_1, \ldots, u_n\} = \{u'_1, \ldots, u'_m\}$, c'est-à-dire $x = x'$.
Ainsi, $X$ est un espace de Kolmogorov.
\end{proof}


\begin{lemma}
\label{lemma-dimension-specializations-lift}
Soit $f : X \to Y$ un morphisme d'espaces topologiques.
Supposons que $Y$ soit un espace topologique sobre et que $f$ soit surjectif.
Si les spécialisations ou les générisations se relèvent le long de $f$, alors
$\dim(X) \geq \dim(Y)$.
\end{lemma}

\begin{proof}
Supposons que les spécialisations se relèvent le long de $f$.
Soit $Z_0 \subset Z_1 \subset \ldots Z_e \subset Y$ une chaîne de fermés
irréductibles de $X$. Soit $\xi_e \in X$ un point s'envoyant sur le point
générique de $Z_e$. Par hypothèse, il existe une spécialisation
$\xi_e \leadsto \xi_{e - 1}$ dans $X$ telle que $\xi_{e - 1}$ s'envoie sur le
point générique de $Z_{e - 1}$. En poursuivant ainsi, on obtient une suite de
spécialisations
$$
\xi_e \leadsto \xi_{e - 1} \leadsto \ldots \leadsto \xi_0
$$
où $\xi_i$ s'envoie sur le point générique de $Z_i$.
Cela implique clairement que la suite de fermés irréductibles
$$
\overline{\{\xi_0\}} \subset
\overline{\{\xi_1\}} \subset \ldots
\overline{\{\xi_e\}}
$$
est une chaîne de longueur $e$ dans $X$.
Le cas où les générisations se relèvent le long de $f$ est analogue.
\end{proof}

\begin{lemma}
\label{lemma-characterize-closed-Noetherian}
Soit $X$ un espace topologique noethérien et sobre.
Soit $E \subset X$ une partie de $X$.
\begin{enumerate}
\item Si $E$ est constructible et stable par spécialisation, alors $E$ est
fermé.
\item Si $E$ est constructible et stable par générisation, alors $E$ est
ouvert.
\end{enumerate}
\end{lemma}

\begin{proof}
Soit $E$ constructible et stable par générisation.
Soit $Y \subset X$ un fermé irréductible de point générique $\xi \in Y$.
Si $E \cap Y$ est non vide, alors il contient $\xi$ (par stabilité par
générisation), et est donc dense dans $Y$ ; il contient par conséquent un
ouvert non vide de $Y$ ; voir le
lemme \ref{lemma-characterize-constructible-Noetherian}.
Ainsi, $E$ est ouvert d'après le
lemme \ref{lemma-characterize-open-Noetherian}.
Ceci démontre (2). Pour démontrer (1), appliquons (2) au complémentaire de $E$
dans $X$.
\end{proof}









\section{Fonctions de dimension}
\label{section-dimension-function}

\noindent
Il n'est guère pertinent de considérer des fonctions de dimension à moins que
l'espace considéré ne soit sobre (Définition \ref{definition-generic-point}).
On pourrait donc améliorer la définition ci-dessous en considérant l'espace
topologique sobre associé à $X$. Puisque l'espace topologique sous-jacent à un
schéma est sobre, nous ne nous attarderons pas sur cette amélioration.

\begin{definition}
\label{definition-dimension-function}
Soit $X$ un espace topologique.
\begin{enumerate}
\item Soient $x, y \in X$, $x \not = y$. Supposons que $x \leadsto y$,
c'est-à-dire que $y$ soit une spécialisation de $x$.
On dit que $y$ est une {\it spécialisation immédiate}
de $x$ s'il n'existe aucun
$z \in X \setminus \{x, y\}$ tel que $x \leadsto z$ et $z \leadsto y$.
\item Une application $\delta : X \to \mathbf{Z}$ est appelée une
{\it fonction de dimension}\footnote{Il s'agit vraisemblablement d'une
notation non standard. Cette notion n'est habituellement introduite que
pour les schémas (localement) noethériens, auquel cas la condition (a) résulte
de (b).} si
\begin{enumerate}
\item dès que $x \leadsto y$ et $x \not = y$,
on a $\delta(x) > \delta(y)$, et
\item pour toute spécialisation immédiate $x \leadsto y$ dans $X$,
on a $\delta(x) = \delta(y) + 1$.
\end{enumerate}
\end{enumerate}
\end{definition}

\noindent
Il est clair que si $\delta$ est une fonction de dimension, alors
$\delta + t$ l'est aussi pour tout $t \in \mathbf{Z}$. Voici un petit lemme.

\begin{lemma}
\label{lemma-dimension-function-catenary}
Soit $X$ un espace topologique. Si $X$ est sobre et possède une fonction de
dimension, alors $X$ est caténaire. De plus, pour toute relation de
spécialisation $x \leadsto y$,
on a
$$
\delta(x) - \delta(y) =
\text{codim}\left(\overline{\{y\}}, \ \overline{\{x\}}\right).
$$
\end{lemma}

\begin{proof}
Supposons que $Y \subset Y' \subset X$ soient des fermés irréductibles.
Soient $\xi \in Y$, $\xi' \in Y'$ leurs points génériques.
Les définitions donnent immédiatement
$\text{codim}(Y, Y') \leq \delta(\xi) - \delta(\xi') < \infty$.
En fait, la première inégalité est une égalité. En effet, supposons que
$$
Y = Y_0 \subset Y_1 \subset \ldots \subset Y_e = Y'
$$
soit une chaîne maximale quelconque de fermés irréductibles. Notons
$\xi_i \in Y_i$ le point générique. Alors
$\xi_i \leadsto \xi_{i + 1}$ est une spécialisation immédiate.
Ainsi, $e = \delta(\xi) - \delta(\xi')$, comme voulu.
Cela démontre également la dernière assertion du lemme.
\end{proof}

\begin{lemma}
\label{lemma-dimension-function-unique}
Soit $X$ un espace topologique.
Soient $\delta$, $\delta'$ deux fonctions de dimension sur $X$.
Si $X$ est localement noethérien et sobre, alors $\delta - \delta'$ est
localement constante sur $X$.
\end{lemma}

\begin{proof}
Soit $x \in X$ un point. Nous allons montrer que $\delta - \delta'$ est
constante dans un voisinage de $x$.
Nous pouvons remplacer $X$ par un voisinage ouvert
noethérien de $x$ dans $X$. Nous pouvons donc supposer que $X$ est
noethérien et sobre.
Soient $Z_1, \ldots, Z_r$ les composantes irréductibles de $X$ qui
contiennent $x$. (Il n'y en a qu'un nombre fini puisque
$X$ est noethérien, voir le Lemme \ref{lemma-Noetherian}.)
Soit $\xi_i \in Z_i$ le point générique.
Notons que $Z_1 \cup \ldots \cup Z_r$ est un voisinage de $x$ dans $X$
(pas nécessairement fermé). Nous affirmons que $\delta - \delta'$ est
constante sur $Z_1 \cup \ldots \cup Z_r$. En effet, si $y \in Z_i$,
alors
$$
\delta(x) - \delta(y) = \delta(x) - \delta(\xi_i) + \delta(\xi_i) - \delta(y)
= - \text{codim}(\overline{\{x\}}, Z_i)
+ \text{codim}(\overline{\{y\}}, Z_i)
$$
d'après le Lemme \ref{lemma-dimension-function-catenary}.
Il en va de même pour $\delta'$, d'où le résultat.
\end{proof}

\begin{lemma}
\label{lemma-locally-dimension-function}
Soit $X$ un espace localement noethérien, sobre et caténaire.
Alors tout point possède un voisinage ouvert
$U \subset X$ qui admet une fonction de dimension.
\end{lemma}

\begin{proof}
Nous utiliserons à plusieurs reprises
le fait qu'un sous-espace ouvert d'un espace caténaire est caténaire, voir le
Lemme \ref{lemma-catenary}, et qu'un espace topologique noethérien
n'a qu'un nombre fini de composantes irréductibles, voir le Lemme
\ref{lemma-Noetherian}.
Dans la démonstration du Lemme \ref{lemma-dimension-function-unique}, nous avons
vu comment construire une telle fonction. Tout d'abord, nous remplaçons $X$
par un voisinage ouvert noethérien de $x$. Ensuite, soient
$Z_1, \ldots, Z_r \subset X$ les composantes irréductibles de $X$. Soit
$$
Z_i \cap Z_j = \bigcup Z_{ijk}
$$
la décomposition en composantes irréductibles. Nous remplaçons
$X$ par
$$
X \setminus \left(
\bigcup\nolimits_{x \not \in Z_i} Z_i
\cup
\bigcup\nolimits_{x \not \in Z_{ijk}} Z_{ijk}
\right)
$$
de sorte que nous puissions supposer que $x \in Z_i$ pour tout $i$ et
$x \in Z_{ijk}$ pour tous $i, j, k$. Pour $y \in X$, choisissons un
$i$ quelconque tel que $y \in Z_i$ et posons
$$
\delta(y) = - \text{codim}(\overline{\{y\}}, Z_i)
+ \text{codim}(\overline{\{x\}}, Z_i).
$$
Nous affirmons qu'il s'agit d'une fonction de dimension. Montrons d'abord
qu'elle est bien définie, c'est-à-dire qu'elle ne dépend pas du choix de $i$.
Supposons en effet que $y \in Z_{ijk}$ pour certains $i, j, k$.
Alors nous avons (en utilisant le Lemme
\ref{lemma-catenary-in-codimension})
\begin{align*}
\delta(y) & =
- \text{codim}(\overline{\{y\}}, Z_i)
+ \text{codim}(\overline{\{x\}}, Z_i) \\
& =
- \text{codim}(\overline{\{y\}}, Z_{ijk})
- \text{codim}(Z_{ijk}, Z_i)
+ \text{codim}(\overline{\{x\}}, Z_{ijk})
+ \text{codim}(Z_{ijk}, Z_i) \\
& =
- \text{codim}(\overline{\{y\}}, Z_{ijk})
+ \text{codim}(\overline{\{x\}}, Z_{ijk})
\end{align*}
ce qui est symétrique en $i$ et $j$.
Nous omettons la démonstration du fait qu'il s'agit d'une fonction de dimension.
\end{proof}

\begin{remark}
\label{remark-obstruction-to-dimension-function}
En combinant les Lemmes \ref{lemma-dimension-function-unique} et
\ref{lemma-locally-dimension-function}, nous voyons que, sur un espace
topologique caténaire, localement noethérien et sobre, l'obstruction à
l'existence d'une fonction de dimension est un élément de
$H^1(X, \mathbf{Z})$.
\end{remark}



\section{Ensembles nulle part denses}
\label{section-nowhere-dense}

\begin{definition}
\label{definition-nowhere-dense}
Soit $X$ un espace topologique.
\begin{enumerate}
\item Étant donnée une partie $T \subset X$, l'{\it intérieur} de $T$ est le
plus grand ouvert de $X$ contenu dans $T$.
\item Une partie $T \subset X$ est dite {\it nulle part dense} si l'adhérence
de $T$ a un intérieur vide.
\end{enumerate}
\end{definition}

\begin{lemma}
\label{lemma-nowhere-dense}
Soit $X$ un espace topologique. La réunion d'un nombre fini d'ensembles nulle
part denses est un ensemble nulle part dense.
\end{lemma}

\begin{proof}
Il suffit de démontrer le lemme pour deux ensembles nulle part denses, le cas
général en résultant par récurrence. Soient $A,B \subset X$ deux parties nulle
part denses. On a $\overline{A \cup B} = \overline{A} \cup \overline{B}$.
Ainsi, si $U \subset \overline{A \cup B}$ est un ouvert de $X$, alors
$U \setminus U \cap \overline{B}$ est un ouvert de $U$ et de $X$,
contenu dans $\overline{A}$, et est donc vide. De même,
$U \setminus U \cap \overline{A}$ est vide. Ainsi $U = \emptyset$,
comme voulu.
\end{proof}

\begin{lemma}
\label{lemma-image-nowhere-dense-open}
Soit $X$ un espace topologique.
Soit $U \subset X$ un ouvert.
Soit $T \subset U$ une partie.
Si $T$ est nulle part dense dans $U$, alors $T$ est nulle part dense dans $X$.
\end{lemma}

\begin{proof}
Supposons que $T$ soit nulle part dense dans $U$.
Supposons que $x \in X$ soit un point intérieur à l'adhérence
$\overline{T}$ de $T$ dans $X$. Il existe alors
$x \in V \subset \overline{T}$, où $V \subset X$ est ouvert dans $X$.
Notons que $\overline{T} \cap U$ est
l'adhérence de $T$ dans $U$. Puisque l'intérieur de
$\overline{T} \cap U$ est vide, on a $V \cap U = \emptyset$. Par conséquent,
$x$ ne peut appartenir à l'adhérence de $U$, et a fortiori ne peut appartenir
à l'adhérence de $T$, ce qui est une contradiction.
\end{proof}

\begin{lemma}
\label{lemma-nowhere-dense-local}
Soit $X$ un espace topologique.
Soit $X = \bigcup U_i$ un recouvrement ouvert.
Soit $T \subset X$ une partie.
Si $T \cap U_i$ est nulle part dense dans $U_i$ pour tout $i$,
alors $T$ est nulle part dense dans $X$.
\end{lemma}

\begin{proof}
Notons $\overline{T}_i$ l'adhérence de $T \cap U_i$ dans $U_i$.
On a $\overline{T} \cap U_i = \overline{T}_i$.
La prise de l'intérieur commute avec l'intersection par un ouvert, donc
$$
(\text{intérieur de }\overline{T}) \cap U_i =
\text{intérieur de }(\overline{T} \cap U_i) =
\text{intérieur dans }U_i\text{ de }\overline{T}_i
$$
Par hypothèse, le dernier de ces ensembles est vide. Ainsi,
$T$ est nulle part dense dans $X$.
\end{proof}

\begin{lemma}
\label{lemma-closed-image-nowhere-dense}
Soit $f : X \to Y$ une application continue d'espaces topologiques.
Soit $T \subset X$ une partie.
Si $f$ est un homéomorphisme de $X$ sur une partie fermée de $Y$
et si $T$ est nulle part dense dans $X$, alors $f(T)$ est également nulle
part dense dans $Y$.
\end{lemma}

\begin{proof}
Puisque $f(X)$ est fermé dans $Y$ et que $f$ est un homéomorphisme de $X$
sur $f(X)$, l'adhérence de $f(T)$ dans $Y$ est égale à $f(\overline{T})$.
Ainsi, si $V \subset Y$ est ouvert et contenu dans l'adhérence de $f(T)$,
alors $U = f^{-1}(V)$ est ouvert et contenu dans $\overline{T}$.
Par conséquent, $U = \emptyset$, ce qui entraîne à son tour que
$V = \emptyset$, comme voulu.
\end{proof}

\begin{lemma}
\label{lemma-open-inverse-image-closed-nowhere-dense}
Soit $f : X \to Y$ une application continue d'espaces topologiques.
Soit $T \subset Y$ une partie.
Si $f$ est ouverte et si $T$ est une partie fermée nulle part dense de $Y$,
alors $f^{-1}(T)$ est également une partie fermée nulle part dense de $X$.
Si $f$ est surjective et ouverte, alors
$T$ est fermée et nulle part dense si et seulement
si $f^{-1}(T)$ est fermée et nulle part dense.
\end{lemma}

\begin{proof}
Démonstration omise. (Indication : dans le premier cas, l'intérieur de
$f^{-1}(T)$ s'envoie dans l'intérieur de $T$, tandis que, dans le second cas,
l'intérieur de $f^{-1}(T)$ s'envoie sur l'intérieur de $T$.)
\end{proof}

\begin{lemma}
\label{lemma-dense-in-constructible}
Soit $X$ un espace topologique. Soit $E \subset X$ une réunion finie de
parties localement fermées (par exemple, $E$ est constructible). Alors $E$
contient un ouvert dense de son adhérence.
\end{lemma}

\begin{proof}
Après avoir remplacé $X$ par l'adhérence de $E$, nous pouvons supposer que
$E$ est dense dans $X$ et devons montrer que $E$ contient un ouvert dense.
Le complémentaire $T$ de $E$ est alors une réunion finie de parties localement
fermées, disons $T = T_1 \cup \ldots \cup T_n$, et $T$ est nulle part dense
dans $X$. Puisque $T_i$ est un ouvert dense de son adhérence $Z_i$, nous voyons
que $Z_i$ est également nulle part dense dans $X$. L'adhérence
$Z = Z_1 \cup \ldots \cup Z_n$ de $T$ est donc nulle part dense dans $X$
d'après le Lemme \ref{lemma-nowhere-dense}.
Nous en concluons que $E$ contient l'ouvert dense $X \setminus Z$.
\end{proof}





\section{Espaces profinis}
\label{section-profinite}

\noindent
Voici la définition.

\begin{definition}
\label{definition-profinite}
Un espace topologique est dit {\it profini} s'il est homéomorphe à la limite
d'un diagramme d'espaces discrets finis.
\end{definition}

\noindent
Ce n'est pas la caractérisation la plus commode d'un espace profini.

\begin{lemma}
\label{lemma-profinite}
Soit $X$ un espace topologique.
Les assertions suivantes sont équivalentes :
\begin{enumerate}
\item $X$ est un espace profini, et
\item $X$ est séparé, quasi-compact et totalement discontinu.
\end{enumerate}
Si ces conditions sont satisfaites, alors $X$ est une limite cofiltrante
d'espaces discrets finis.
\end{lemma}

\begin{proof}
Supposons (1). Choisissons un diagramme $i \mapsto X_i$ d'espaces discrets finis
tel que $X = \lim X_i$. Comme chaque $X_i$ est séparé et quasi-compact, le
lemme \ref{lemma-inverse-limit-quasi-compact} montre que $X$ est quasi-compact.
Si $x, x' \in X$ sont des points distincts, alors $x$ et $x'$ ont des images
distinctes dans un certain $X_i$. Par conséquent, $x$ et $x'$ possèdent des
voisinages ouverts disjoints, c'est-à-dire que $X$ est séparé. Exactement de la
même manière, nous voyons que $X$ est totalement discontinu.

\medskip\noindent
Supposons (2). Soit $\mathcal{I}$ l'ensemble des décompositions en réunion
disjointe finie $X = \coprod_{i \in I} U_i$, où $U_i$ est un ouvert non vide
(et fermé) pour tout $i \in I$.
Pour chaque $I \in \mathcal{I}$, il existe une application continue
$X \to I$ qui envoie tout point de $U_i$ sur $i$. Nous définissons un ordre
partiel : $I \leq I'$ pour $I, I' \in \mathcal{I}$ si et seulement si le
recouvrement correspondant à $I'$ raffine le recouvrement correspondant à
$I$. Dans ce cas, nous obtenons une application canonique $I' \to I$. En
d'autres termes, nous obtenons un système projectif d'espaces discrets finis
indexé par $\mathcal{I}$.
Les applications $X \to I$ sont compatibles et fournissent une application
continue
$$
X \longrightarrow \lim_{I \in \mathcal{I}} I
$$
Nous affirmons que cette application est un homéomorphisme, ce qui achève la
démonstration. (L'assertion finale en résulte également, car l'ensemble
partiellement ordonné $\mathcal{I}$ est dirigé : étant données deux
décompositions de $X$ en réunions disjointes, nous pouvons en trouver une
troisième qui raffine les deux.)
En effet, l'application est injective puisque $X$ est totalement discontinu
et que, par conséquent, $\{x\}$ est l'intersection de toutes les parties
ouvertes et fermées de $X$ qui contiennent $x$
(lemme \ref{lemma-connected-component-intersection-compact-Hausdorff}) ;
l'application est surjective d'après le
lemme \ref{lemma-intersection-closed-in-quasi-compact}.
D'après le lemme \ref{lemma-bijective-map}, l'application est un homéomorphisme.
\end{proof}

\begin{lemma}
\label{lemma-directed-inverse-limit-profinite}
Toute limite d'espaces profinis est un espace profini.
\end{lemma}

\begin{proof}
Soit $i \mapsto X_i$ un diagramme d'espaces profinis
indexé par la catégorie $\mathcal{I}$.
Utilisons la caractérisation des espaces profinis donnée par le
lemme \ref{lemma-profinite}. En particulier, chaque $X_i$ est
séparé, quasi-compact et totalement discontinu.
D'après le lemme \ref{lemma-limits}, la limite $X = \lim X_i$ existe.
D'après le lemme \ref{lemma-inverse-limit-quasi-compact},
la limite $X$ est quasi-compacte. Soient $x, x' \in X$ des points distincts.
Il existe alors un $i$ tel que $x$ et $x'$ aient des images distinctes
$x_i$ et $x'_i$ dans $X_i$ par la projection $X \to X_i$.
Les points $x_i$ et $x'_i$ possèdent alors des voisinages ouverts disjoints
dans $X_i$. En prenant les images réciproques de ces ouverts, nous concluons
que $X$ est séparé.
De même, $x_i$ et $x'_i$ appartiennent à des composantes connexes distinctes
de $X_i$ ; nécessairement, $x$ et $x'$ appartiennent donc à des composantes
connexes distinctes de $X$. Ainsi, $X$ est totalement discontinu.
Cela achève la démonstration.
\end{proof}

\begin{lemma}
\label{lemma-profinite-refine-open-covering}
Soit $X$ un espace profini. Tout recouvrement ouvert de $X$ admet un
raffinement par un recouvrement fini $X = \coprod U_i$ dont les $U_i$ sont
ouverts et fermés.
\end{lemma}

\begin{proof}
Écrivons $X = \lim X_i$ comme limite d'un système projectif d'espaces discrets
finis indexé par un ensemble dirigé $I$ (lemme \ref{lemma-profinite}).
Notons $f_i : X \to X_i$ la projection.
Pour tout point $x = (x_i) \in X$, la famille des $f_i^{-1}(\{x_i\})$ forme
un système fondamental de voisinages ouverts. Ainsi, puisque $X$ est
quasi-compact, nous pouvons supposer donné un recouvrement ouvert
$$
X = f_{i_1}^{-1}(\{x_{i_1}\}) \cup \ldots \cup f_{i_n}^{-1}(\{x_{i_n}\})
$$
Choisissons $i \in I$ tel que $i \geq i_j$ pour $j = 1, \ldots, n$ (c'est
possible puisque $I$ est un ensemble dirigé). Nous voyons alors que le
recouvrement
$$
X = \coprod\nolimits_{t \in X_i} f_i^{-1}(\{t\})
$$
raffine le recouvrement donné et possède la forme voulue.
\end{proof}

\begin{lemma}
\label{lemma-pi0-profinite}
Soit $X$ un espace topologique. Si $X$ est quasi-compact et si toute composante
connexe de $X$ est l'intersection des parties ouvertes et fermées qui la
contiennent, alors $\pi_0(X)$ est un espace profini.
\end{lemma}

\begin{proof}
Nous allons appliquer le lemme \ref{lemma-profinite}.
Puisque $\pi_0(X)$ est l'image d'un espace quasi-compact, il est quasi-compact
(lemme \ref{lemma-image-quasi-compact}).
Il est totalement discontinu par construction
(lemme \ref{lemma-space-connected-components}).
Soient $C, D \subset X$ deux composantes connexes distinctes de $X$.
Écrivons $C = \bigcap U_\alpha$ comme intersection des parties ouvertes et
fermées de $X$ qui contiennent $C$. Toute intersection finie de $U_\alpha$
est encore une telle partie. Puisque $\bigcap U_\alpha \cap D = \emptyset$,
nous en concluons que $U_\alpha \cap D = \emptyset$ pour un certain $\alpha$
(utiliser les lemmes \ref{lemma-connected-components},
\ref{lemma-closed-in-quasi-compact} et
\ref{lemma-intersection-closed-in-quasi-compact}).
Comme $U_\alpha$ est ouvert et fermé, il est la réunion des composantes
connexes qu'il contient ; autrement dit, $U_\alpha$ est l'image réciproque
d'une certaine partie ouverte et fermée $V_\alpha \subset \pi_0(X)$.
Cela montre que les points correspondant à $C$ et $D$ sont contenus dans des
ouverts disjoints, c'est-à-dire que $\pi_0(X)$ est séparé.
\end{proof}





\section{Espaces spectraux}
\label{section-spectral}

\noindent
Le contenu de cette section est tiré de \cite{Hochster} et de
\cite{Hochster-thesis}. Dans sa thèse, Hochster démontre (entre autres)
que les espaces spectraux sont exactement les espaces topologiques qui
sont des spectres d'anneaux.

\begin{definition}
\label{definition-spectral-space}
Un espace topologique $X$ est dit {\it spectral} s'il est sobre et
quasi-compact, si l'intersection de deux ouverts quasi-compacts est
quasi-compacte et si l'ensemble des ouverts quasi-compacts forme une
base de la topologie. Une application continue $f : X \to Y$ entre espaces
spectraux est dite {\it spectrale} si l'image réciproque de tout ouvert
quasi-compact est quasi-compacte.
\end{definition}

\noindent
Autrement dit, une application continue entre espaces spectraux est spectrale
si et seulement si elle est quasi-compacte (Définition
\ref{definition-quasi-compact}).

\medskip\noindent
Soit $X$ un espace spectral. La {\it topologie constructible} sur $X$
est la topologie dont une sous-base d'ouverts est formée des ensembles $U$
et $U^c$, où $U$ est un ouvert quasi-compact de $X$. Puisque $X$ est spectral,
un ouvert $U \subset X$ est rétrocompact si et seulement si $U$ est
quasi-compact. La topologie constructible peut donc aussi être caractérisée
comme la topologie la moins fine pour laquelle toute partie constructible de
$X$ est à la fois ouverte et fermée (voir la section
\ref{section-constructible} pour les définitions et les propriétés des parties
constructibles).
Il s'ensuit qu'une partie de $X$ est ouverte, resp.\ fermée, pour la topologie
constructible si et seulement si elle est une réunion, resp.\ une intersection,
de parties constructibles. Comme l'ensemble des ouverts quasi-compacts est une
base de la topologie de $X$, la topologie constructible est plus fine que la
topologie donnée sur $X$.

\begin{lemma}
\label{lemma-constructible-hausdorff-quasi-compact}
Soit $X$ un espace spectral. La topologie constructible est séparée,
totalement discontinue et quasi-compacte.
\end{lemma}

\begin{proof}
Soient $x, y \in X$ avec $x \not = y$. Comme $X$ est sobre, il existe
un ouvert $U$ qui contient exactement l'un des deux points $x, y$.
Supposons que $x \in U$. On peut remplacer $U$ par un voisinage ouvert
quasi-compact de $x$ contenu dans $U$. Alors $U$ et $U^c$ sont ouverts et
fermés pour la topologie constructible. Ainsi, $X$ est séparé pour la
topologie constructible, car $x \in U$ et $y \in U^c$ appartiennent à des
ouverts disjoints de cette topologie. L'existence de $U$ implique aussi que
$x$ et $y$ appartiennent à des composantes connexes distinctes pour la
topologie constructible ; ainsi, $X$ est totalement discontinu pour cette
topologie.

\medskip\noindent
Soit $\mathcal{B}$ l'ensemble des parties $B \subset X$ telles que $B$ soit
un ouvert quasi-compact ou un fermé dont le complémentaire est
quasi-compact. Si $B \in \mathcal{B}$, alors $B^c \in \mathcal{B}$.
Il suffit de montrer que tout recouvrement $X = \bigcup_{i \in I} B_i$
avec $B_i \in \mathcal{B}$ admet un raffinement fini, voir le
Lemme \ref{lemma-subbase-theorem}.
En passant aux complémentaires, on voit qu'il faut montrer que toute famille
$\{B_i\}_{i \in I}$ d'éléments de $\mathcal{B}$ telle que
$B_{i_1} \cap \ldots \cap B_{i_n} \not = \emptyset$
pour tout $n$ et tous $i_1, \ldots, i_n \in I$
possède un point commun. On peut supposer, et l'on suppose,
que $B_i \not = B_{i'}$ si $i \not = i'$.

\medskip\noindent
Raisonnons par l'absurde et supposons que $\{B_i\}_{i \in I}$ soit une famille
d'éléments de $\mathcal{B}$ ayant la propriété d'intersection finie, mais dont
l'intersection est vide. Une application du lemme de Zorn montre que l'on
peut supposer cette famille maximale (les détails sont omis).
Soit $I' \subset I$ l'ensemble des indices pour lesquels $B_i$ est fermé,
et posons $Z = \bigcap_{i \in I'} B_i$.
C'est une partie fermée non vide de $X$, d'après le Lemme
\ref{lemma-intersection-closed-in-quasi-compact}.
Si $Z$ est réductible, on peut écrire $Z = Z' \cup Z''$
comme réunion de deux parties fermées, toutes deux distinctes de $Z$. On peut
donc, en particulier, trouver un ouvert quasi-compact $U' \subset X$ qui
rencontre $Z'$ mais pas $Z''$. De même, on peut trouver un ouvert quasi-compact
$U'' \subset X$ qui rencontre $Z''$ mais pas $Z'$. Posons
$B' = X \setminus U'$ et $B'' = X \setminus U''$. Remarquons que
$Z'' \subset B'$ et $Z' \subset B''$.
S'il existe un nombre fini d'indices $i_1, \ldots, i_n \in I$ tels que
$B' \cap B_{i_1} \cap \ldots \cap B_{i_n} = \emptyset$
ainsi qu'un nombre fini d'indices $j_1, \ldots, j_m \in I$ tels que
$B'' \cap B_{j_1} \cap \ldots \cap B_{j_m} = \emptyset$,
alors on obtient
$Z \cap B_{i_1} \cap \ldots \cap B_{i_n} \cap B_{j_1} \cap \ldots \cap B_{j_m}
= \emptyset$.
Cependant, l'ensemble
$B_{i_1} \cap \ldots \cap B_{i_n} \cap B_{j_1} \cap \ldots \cap B_{j_m}$
est quasi-compact ; il existerait donc un nombre fini d'indices
$i'_1, \ldots, i'_l \in I'$ tels que
$B_{i_1} \cap \ldots \cap B_{i_n} \cap B_{j_1} \cap \ldots \cap
B_{j_m} \cap B_{i'_1} \cap \ldots \cap B_{i'_l} = \emptyset$, ce qui est
contradictoire. On voit donc que l'on peut adjoindre soit $B'$, soit $B''$,
à la famille donnée, contrairement à sa maximalité. On en conclut que $Z$
est irréductible. Mais cela conduit encore à une contradiction, car tout
ouvert non vide $Z \cap B_i$, pour $i \in I \setminus I'$, contient alors
(par le même argument que ci-dessus) l'unique point générique de $Z$.
Cette contradiction démontre le lemme.
\end{proof}

\begin{lemma}
\label{lemma-fibres-spectral-map-quasi-compact}
Soit $f : X \to Y$ une application spectrale entre espaces spectraux. Alors
\begin{enumerate}
\item $f$ est continue pour les topologies constructibles,
\item les fibres de $f$ sont quasi-compactes, et
\item l'image est fermée pour la topologie constructible.
\end{enumerate}
\end{lemma}

\begin{proof}
Notons $X'$ et $Y'$ les espaces $X$ et $Y$ munis de leur topologie
constructible ; ce sont des espaces quasi-compacts séparés d'après le
Lemme \ref{lemma-constructible-hausdorff-quasi-compact}.
L'assertion (1) affirme que $X' \to Y'$ est continue et résulte immédiatement
des définitions. L'assertion (3) résulte du fait que $f(X')$ est une partie
quasi-compacte de l'espace séparé $Y'$, voir le
Lemme \ref{lemma-quasi-compact-in-Hausdorff}.
On a un diagramme commutatif
$$
\xymatrix{
X' \ar[r] \ar[d] & X \ar[d] \\
Y' \ar[r] & Y
}
$$
d'applications continues entre espaces topologiques. Comme $Y'$ est séparé,
les fibres $X'_y$ sont fermées dans $X'$. Puisque $X'$ est quasi-compact,
$X'_y$ est quasi-compact
(Lemme \ref{lemma-closed-in-quasi-compact}).
Comme $X'_y \to X_y$ est une application continue surjective, on en conclut
que $X_y$ est quasi-compact (Lemme \ref{lemma-image-quasi-compact}).
\end{proof}

\begin{lemma}
\label{lemma-spectral-if-continuous-wrt-constructible-top}
Soient $X$ et $Y$ des espaces spectraux et $f : X \to Y$ une application
continue. Alors $f$ est spectrale si et seulement si $f$ est continue pour
les topologies constructibles.
\end{lemma}

\begin{proof}
L'implication directe est le
Lemme \ref{lemma-fibres-spectral-map-quasi-compact}.
Supposons $f$ continue pour les topologies constructibles.
Soit $V \subset Y$ un ouvert quasi-compact.
Alors $V$ est ouvert et fermé pour la topologie constructible.
Ainsi, $f^{-1}(V)$ est ouvert et fermé pour la topologie constructible.
Ainsi, $f^{-1}(V)$ est quasi-compact pour la topologie constructible, puisque
$X$ est quasi-compact pour cette topologie d'après le
Lemme \ref{lemma-constructible-hausdorff-quasi-compact}.
L'application identité $f^{-1}(V) \to f^{-1}(V)$ étant continue et surjective
de la topologie constructible vers la topologie usuelle, on conclut que
$f^{-1}(V)$ est quasi-compact pour la topologie de $X$, d'après le
Lemme \ref{lemma-image-quasi-compact}.
Ceci achève la démonstration.
\end{proof}

\begin{lemma}
\label{lemma-spectral-sub}
Soit $X$ un espace spectral. Soit $E \subset X$ une partie fermée pour la
topologie constructible (par exemple une partie constructible ou fermée).
Alors $E$, muni de la topologie induite, est un espace spectral.
\end{lemma}

\begin{proof}
Soit $Z \subset E$ une partie fermée irréductible. Soit $\eta$ le point
générique de l'adhérence $\overline{Z}$ de $Z$ dans $X$. Pour démontrer que
$E$ est sobre, montrons que $\eta \in E$. Sinon, puisque $E$ est fermé pour
la topologie constructible, il existe une partie constructible $F \subset X$
telle que $\eta \in F$ et $F \cap E = \emptyset$.
D'après le Lemme \ref{lemma-generic-point-in-constructible}, cela implique que
$F \cap \overline{Z}$ contient un ouvert non vide de $\overline{Z}$.
Mais c'est impossible puisque $\overline{Z}$ est l'adhérence de $Z$ et que
$Z \cap F = \emptyset$.

\medskip\noindent
Puisque $E$ est fermé pour la topologie constructible, il est quasi-compact
pour cette topologie
(Lemmes \ref{lemma-closed-in-quasi-compact} et
\ref{lemma-constructible-hausdorff-quasi-compact}). Il est donc, à plus forte
raison, quasi-compact pour la topologie induite par celle de $X$. Si
$U \subset X$ est un ouvert quasi-compact, alors $E \cap U$ est fermé pour la
topologie constructible, donc quasi-compact (comme on vient de le voir). Il en
résulte que les ouverts quasi-compacts de $E$ sont les intersections $E \cap U$
où $U$ est un ouvert quasi-compact de $X$. Ils forment une base de la topologie.
Enfin, étant donnés deux ouverts quasi-compacts $U, U' \subset X$, on a
$(E \cap U) \cap (E \cap U') = E \cap (U \cap U')$, et $U \cap U'$ est
quasi-compact puisque $X$ est spectral. Ceci achève la démonstration.
\end{proof}

\begin{lemma}
\label{lemma-constructible-stable-specialization-closed}
Soit $X$ un espace spectral. Soit $E \subset X$ une partie fermée pour la
topologie constructible (par exemple une partie constructible).
\begin{enumerate}
\item Si $x \in \overline{E}$, alors $x$ est une spécialisation d'un point de
$E$.
\item Si $E$ est stable par spécialisation, alors $E$ est fermée.
\item Si $E' \subset X$ est ouverte pour la topologie constructible
(par exemple constructible) et stable par générisation, alors $E'$ est ouverte.
\end{enumerate}
\end{lemma}

\begin{proof}
Démonstration de (1). Soit $x \in \overline{E}$. Soit $\{U_i\}$ l'ensemble
des voisinages ouverts quasi-compacts de $x$. Toute intersection finie de
$U_i$ est encore un tel voisinage. L'intersection $U_i \cap E$ est non vide
pour tout $i$. Comme les parties $U_i \cap E$ sont fermées pour la topologie
constructible, l'intersection $\bigcap (U_i \cap E)$ est non vide d'après le
Lemme \ref{lemma-constructible-hausdorff-quasi-compact} et le
Lemme \ref{lemma-intersection-closed-in-quasi-compact}.
Puisque $\{U_i\}$ est un système fondamental de voisinages ouverts de $x$,
l'intersection $\bigcap U_i$ est l'ensemble des générisations de $x$.
Ainsi, $x$ est une spécialisation d'un point de $E$.

\medskip\noindent
L'assertion (2) résulte immédiatement de (1).

\medskip\noindent
Démonstration de (3). Supposons que $E'$ vérifie les hypothèses de (3).
Le complémentaire de $E'$ est fermé pour la topologie constructible
(Lemme \ref{lemma-constructible}) et stable par spécialisation
(Lemme \ref{lemma-open-closed-specialization}).
Ce complémentaire est donc fermé d'après (2), c'est-à-dire que $E'$ est ouverte.
\end{proof}

\begin{lemma}
\label{lemma-two-points}
Soit $X$ un espace spectral et soient $x, y \in X$. Alors, ou bien il existe
un troisième point dont $x$ et $y$ sont des spécialisations, ou bien il existe
des voisinages ouverts disjoints de $x$ et de $y$.
\end{lemma}

\begin{proof}
Soit $\{U_i\}$ l'ensemble des voisinages ouverts quasi-compacts de $x$.
Toute intersection finie de $U_i$ est encore un tel voisinage.
Soit $\{V_j\}$ l'ensemble des voisinages ouverts quasi-compacts de $y$.
Toute intersection finie de $V_j$ est encore un tel voisinage.
Si $U_i \cap V_j$ est vide pour certains $i, j$, la conclusion est acquise.
Sinon, $U_i \cap V_j$ est non vide pour tous $i$ et $j$. Les ensembles
$U_i \cap V_j$ sont fermés pour la topologie constructible de $X$. D'après le
Lemme \ref{lemma-constructible-hausdorff-quasi-compact},
l'intersection $\bigcap (U_i \cap V_j)$ est non vide
(Lemme \ref{lemma-intersection-closed-in-quasi-compact}).
Comme $X$ est sobre et que $\{U_i\}$ est un système fondamental de voisinages
ouverts de $x$, l'intersection $\bigcap U_i$ est l'ensemble des générisations
de $x$. De même, $\bigcap V_j$ est l'ensemble des générisations de $y$.
Ainsi, tout élément de $\bigcap (U_i \cap V_j)$ a pour spécialisations à la
fois $x$ et $y$.
\end{proof}

\begin{lemma}
\label{lemma-characterize-profinite-spectral}
Soit $X$ un espace spectral. Les propriétés suivantes sont équivalentes :
\begin{enumerate}
\item $X$ est profini,
\item $X$ est séparé,
\item $X$ est totalement discontinu,
\item tout ouvert quasi-compact est fermé,
\item il n'existe aucune spécialisation non triviale entre points,
\item tout point de $X$ est fermé,
\item tout point de $X$ est le point générique d'une composante irréductible
de $X$,
\item la topologie constructible est égale à la topologie donnée sur $X$, et
\item ajouter quelque chose ici.
\end{enumerate}
\end{lemma}

\begin{proof}
Le Lemme \ref{lemma-profinite} donne l'implication (1) $\Rightarrow$ (3).
Les composantes irréductibles sont fermées ; si $X$ est totalement discontinu,
tout point est donc fermé. Ainsi, (3) implique (6). L'équivalence de (6) et
(5) est immédiate, et (6) $\Leftrightarrow$ (7) parce que $X$ est sobre.
Supposons (5). Alors toutes les parties constructibles de $X$ sont fermées
(Lemme \ref{lemma-constructible-stable-specialization-closed}), et en
particulier tous les ouverts quasi-compacts sont fermés. Donc (5) implique (4).
Comme $X$ est sobre, pour deux points quelconques il existe un ouvert
quasi-compact qui contient exactement l'un des deux ; par conséquent, (4)
implique (2). Les assertions (4) et (8) sont équivalentes par définition de la
topologie constructible.
Il reste à démontrer que (2) implique (1). Supposons $X$ séparé. Tout ouvert
quasi-compact est alors fermé
(Lemme \ref{lemma-quasi-compact-in-Hausdorff}). Cela implique que $X$ est
totalement discontinu. Il est donc profini d'après le
Lemme \ref{lemma-profinite}.
\end{proof}

\begin{lemma}
\label{lemma-spectral-pi0}
Si $X$ est un espace spectral, alors $\pi_0(X)$ est un espace profini.
\end{lemma}

\begin{proof}
Combiner les Lemmes \ref{lemma-connected-component-intersection} et
\ref{lemma-pi0-profinite}.
\end{proof}

\begin{lemma}
\label{lemma-product-spectral-spaces}
Le produit de deux espaces spectraux est spectral.
\end{lemma}

\begin{proof}
Soient $X$ et $Y$ des espaces spectraux. Notons $p : X \times Y \to X$ et
$q : X \times Y \to Y$ les projections. Soit $Z \subset X \times Y$ une
partie fermée irréductible. Alors $p(Z) \subset X$ et $q(Z) \subset Y$ sont
irréductibles. Soit $x \in X$ le point générique de l'adhérence de $p(X)$ et
soit $y \in Y$ le point générique de l'adhérence de $q(Y)$. Si
$(x, y) \not \in Z$, il existe des ouverts $x \in U \subset X$ et
$y \in V \subset Y$ tels que $Z \cap U \times V = \emptyset$. Ainsi, $Z$
est contenu dans
$(X \setminus U) \times Y \cup X \times (Y \setminus V)$.
Comme $Z$ est irréductible, on a soit
$Z \subset (X \setminus U) \times Y$, soit
$Z \subset X \times (Y \setminus V)$.
Dans le premier cas, $p(Z) \subset (X \setminus U)$ et, dans le second,
$q(Z) \subset (Y \setminus V)$. Les deux cas sont absurdes, puisque $x$
appartient à l'adhérence de $p(Z)$ et $y$ à celle de $q(Z)$. On en conclut que
$(x, y) \in Z$, ce qui signifie que $(x, y)$ est le point générique de $Z$.

\medskip\noindent
Une base de la topologie de $X \times Y$ est formée des ouverts de la forme
$U \times V$, où $U \subset X$ et $V \subset Y$ sont des ouverts
quasi-compacts (on utilise ici le fait que $X$ et $Y$ sont spectraux).
Alors $U \times V$ est quasi-compact, car le produit de deux espaces
quasi-compacts est quasi-compact. De plus, tout ouvert quasi-compact de
$X \times Y$ est réunion finie de tels rectangles quasi-compacts $U \times V$.
Il s'ensuit que l'intersection de deux tels ouverts est encore quasi-compacte
(puisque $X$ et $Y$ sont spectraux). Ceci achève la démonstration.
\end{proof}

\begin{lemma}
\label{lemma-spectral-bijective}
Soit $f : X \to Y$ une application continue entre espaces topologiques.
Supposons que
\begin{enumerate}
\item $X$ et $Y$ soient spectraux,
\item $f$ soit spectrale et bijective, et
\item les générisations (resp.\ les spécialisations) se relèvent le long de
$f$.
\end{enumerate}
Alors $f$ est un homéomorphisme.
\end{lemma}

\begin{proof}
Puisque $f$ est spectrale, elle définit une application continue entre $X$ et
$Y$ munis de leurs topologies constructibles. D'après les
Lemmes \ref{lemma-constructible-hausdorff-quasi-compact} et
\ref{lemma-bijective-map}, il en résulte que $X \to Y$ est un homéomorphisme
pour les topologies constructibles. Soit $U \subset X$ un ouvert
quasi-compact. Alors $f(U)$ est constructible dans $Y$. Soit $y \in Y$ une
spécialisation d'un point de $f(U)$. La dernière hypothèse montre que
$f^{-1}(y)$ est une spécialisation d'un point de $U$. Ainsi
$f^{-1}(y) \in U$, et donc $y \in f(U)$.
Il en résulte que $f(U)$ est ouvert, voir le
Lemme \ref{lemma-constructible-stable-specialization-closed}.
Par conséquent, $f$ est un homéomorphisme.
Pour démontrer le lemme lorsque les spécialisations se relèvent le long de
$f$, on montre plutôt que $f(Z)$ est fermé si $X \setminus Z$ est un ouvert
quasi-compact de $X$.
\end{proof}

\begin{lemma}
\label{lemma-directed-inverse-limit-finite-sober-spectral-spaces}
La limite projective d'un système projectif filtrant d'espaces topologiques
finis et sobres est un espace topologique spectral.
\end{lemma}

\begin{proof}
Soit $I$ un ensemble préordonné filtrant. Soit $X_i$ un système projectif
d'espaces finis et sobres indexé par $I$. Soit $X = \lim X_i$, qui existe
d'après le Lemme \ref{lemma-limits}. En tant qu'ensemble, $X = \lim X_i$.
Notons $p_i : X \to X_i$ la projection.
Comme $I$ est filtrant, on peut appliquer le
Lemme \ref{lemma-describe-limits}.
Les ouverts $p_i^{-1}(U_i)$, où $U_i \subset X_i$ est ouvert, forment une base
de la topologie. Comme tout recouvrement ouvert de $p_i^{-1}(U_i)$ est en
particulier un recouvrement ouvert pour la topologie profinie, on en conclut
que $p_i^{-1}(U_i)$ est quasi-compact. Si $U_i \subset X_i$ et
$U_j \subset X_j$, alors
$p_i^{-1}(U_i) \cap p_j^{-1}(U_j) = p_k^{-1}(U_k)$
pour un certain $k \geq i, j$ et un ouvert $U_k \subset X_k$. Enfin, si
$Z \subset X$ est fermée et irréductible, alors $p_i(Z) \subset X_i$ est
irréductible et possède donc un unique point générique $\xi_i$ (puisque $X_i$
est un espace topologique fini et sobre). Alors $\xi = \lim \xi_i$ est un
point générique de $Z$ (c'est un point de $Z$ puisque $Z$ est fermé).
Ceci achève la démonstration.
\end{proof}

\begin{lemma}
\label{lemma-spectral-closed-in-product-two-point-space}
Soit $W$ l'espace topologique à deux points dont l'un est fermé et l'autre ne
l'est pas. Un espace topologique est spectral si et seulement s'il est
homéomorphe à un sous-espace d'un produit de copies de $W$ qui est fermé pour
la topologie constructible.
\end{lemma}

\begin{proof}
Écrivons $W = \{0, 1\}$, où $0$ est une spécialisation de $1$, mais non
réciproquement. Soit $I$ un ensemble. L'espace $\prod_{i \in I} W$ est spectral
d'après le
Lemme \ref{lemma-directed-inverse-limit-finite-sober-spectral-spaces}.
On voit donc qu'un sous-espace de $\prod_{i \in I} W$ fermé pour la topologie
constructible est un espace spectral d'après le
Lemme \ref{lemma-spectral-sub}.

\medskip\noindent
Réciproquement, soit $X$ un espace spectral. Soit $U \subset X$ un ouvert
quasi-compact. Considérons l'application continue
$$
f_U : X \longrightarrow W
$$
qui envoie tout point de $U$ sur $1$ et tout point de $X \setminus U$ sur $0$.
En prenant le produit de ces applications, on obtient une application continue
$$
f = \prod f_U : X \longrightarrow \prod\nolimits_U W
$$
Par construction, l'application $f : X \to Y$ est spectrale. D'après le
Lemme \ref{lemma-fibres-spectral-map-quasi-compact}, l'image de $f$ est fermée
pour la topologie constructible. Si $x', x \in X$ sont distincts, alors, puisque
$X$ est sobre, ou bien $x'$ n'est pas une spécialisation de $x$, ou bien
ce dernier n'est pas une spécialisation du premier. Dans les deux cas (les ouverts
quasi-compacts formant une
base de la topologie de $X$), il existe un ouvert quasi-compact $U \subset X$
tel que $f_U(x') \not = f_U(x)$. Ainsi, $f$ est injective. Soit $Y = f(X)$,
muni de la topologie induite. Soit $y' \leadsto y$ une spécialisation dans $Y$,
et écrivons $f(x') = y'$ et $f(x) = y$. Le même raisonnement montre que
$x' \leadsto x$, car, sinon, il existerait un $U$ tel que $x \in U$ et
$x' \not \in U$, ce qui impliquerait
$f_U(x') \not \leadsto f_U(x)$.
On conclut que $f : X \to Y$ est un homéomorphisme d'après le
Lemme \ref{lemma-spectral-bijective}.
\end{proof}

\begin{lemma}
\label{lemma-spectral-inverse-limit-finite-sober-spaces}
Un espace topologique est spectral si et seulement s'il est limite projective
filtrante d'espaces topologiques finis et sobres.
\end{lemma}

\begin{proof}
Une implication est donnée par le
Lemme \ref{lemma-directed-inverse-limit-finite-sober-spectral-spaces}.
Réciproquement, supposons $X$ spectral. On peut alors supposer que
$X \subset \prod_{i \in I} W$ est une partie fermée pour la topologie
constructible, où $W = \{0, 1\}$ est comme dans le
Lemme \ref{lemma-spectral-closed-in-product-two-point-space}.
On peut écrire
$$
\prod\nolimits_{i \in I} W =
\lim_{J \subset I\text{ fini }} \prod\nolimits_{j \in J} W
$$
comme limite cofiltrante.
Pour chaque $J$, soit $X_J \subset \prod_{j \in J} W$ l'image de $X$.
On voit alors que $X = \lim X_J$ en tant qu'ensembles parce que $X$ est fermé
dans le produit muni de la topologie constructible (détail omis).
Un argument formel (omis) sur les limites montre que $X = \lim X_J$ comme
espaces topologiques.
\end{proof}

\begin{lemma}
\label{lemma-Noetherian-goes-to-spectral}
Soit $X$ un espace topologique et soit $c : X \to X'$ l'application universelle
de $X$ vers un espace topologique sobre, voir le
Lemme \ref{lemma-make-sober}.
\begin{enumerate}
\item Si $X$ est quasi-compact, alors $X'$ l'est aussi.
\item Si $X$ est quasi-compact, possède une base d'ouverts quasi-compacts et si
l'intersection de deux ouverts quasi-compacts est quasi-compacte, alors $X'$ est
spectral.
\item Si $X$ est noethérien, alors $X'$ est un espace spectral noethérien.
\end{enumerate}
\end{lemma}

\begin{proof}
Soit $U \subset X$ un ouvert et soit $U' \subset X'$ l'ouvert correspondant,
c'est-à-dire l'ouvert tel que $c^{-1}(U') = U$.
Alors $U$ est quasi-compact si et seulement si $U'$ est quasi-compact, car
l'image réciproque par $c$ établit une bijection entre les ouverts de $X$ et
ceux de $X'$ qui commute aux réunions. Cela démontre en particulier (1).

\medskip\noindent
Démonstration de (2). Il résulte de ce qui précède que $X'$ possède une base
d'ouverts quasi-compacts. Comme $c^{-1}$ commute aussi aux intersections de
paires d'ouverts, on voit que l'intersection de deux ouverts quasi-compacts de
$X'$ est quasi-compacte. Enfin, $X'$ est quasi-compact d'après (1) et sobre par
construction. Ainsi, $X'$ est spectral.

\medskip\noindent
Démonstration de (3). Il est immédiat que $X'$ est noethérien, puisque cette
propriété se définit par la condition de chaîne ascendante sur les ouverts,
laquelle est satisfaite pour $X$. On a déjà vu en (2) que $X'$ est spectral.
\end{proof}

\section{Limites d'espaces spectraux}
\label{section-limits-spectral}

\noindent
Le Lemme \ref{lemma-spectral-inverse-limit-finite-sober-spaces} nous dit
que tout espace spectral est une limite cofiltrante d'espaces finis
sobres. Tout espace fini sobre est un espace spectral, et toute application
continue entre espaces finis sobres est une application spectrale entre espaces
spectraux. Dans cette section, nous démontrons quelques lemmes concernant les
limites de systèmes d'espaces topologiques spectraux dont les applications de
transition sont spectrales.

\begin{lemma}
\label{lemma-inverse-limit-spectral-spaces-quasi-compact}
Soit $\mathcal{I}$ une catégorie. Soit $i \mapsto X_i$ un diagramme
d'espaces spectraux tel que, pour $a : j \to i$ dans $\mathcal{I}$,
l'application correspondante $f_a : X_j \to X_i$ soit spectrale.
\begin{enumerate}
\item Si l'on se donne des parties $Z_i \subset X_i$ fermées pour la topologie
constructible telles que $f_a(Z_j) \subset Z_i$ pour tout $a : j \to i$ dans
$\mathcal{I}$, alors $\lim Z_i$ est quasi-compact.
\item L'espace $X = \lim X_i$ est quasi-compact.
\end{enumerate}
\end{lemma}

\begin{proof}
La limite $Z = \lim Z_i$ existe d'après le Lemme \ref{lemma-limits}.
Notons $X'_i$ l'espace $X_i$ muni de la topologie constructible
et $Z'_i$ le sous-espace correspondant de $X'_i$.
Soit $a : j \to i$ un morphisme de $\mathcal{I}$. Comme $f_a$ est spectrale,
elle définit une application continue $f_a : X'_j \to X'_i$. Ainsi,
$f_a|_{Z_j} : Z'_j \to Z'_i$ est une application continue entre espaces
quasi-compacts séparés (d'après les
Lemmes \ref{lemma-constructible-hausdorff-quasi-compact} et
\ref{lemma-closed-in-quasi-compact}).
Ainsi, $Z' = \lim Z_i$ est quasi-compact d'après le
Lemme \ref{lemma-inverse-limit-quasi-compact}.
Les applications $Z'_i \to Z_i$ sont continues ; par conséquent,
$Z' \to Z$ est continue et induit une bijection entre les ensembles
sous-jacents. Ainsi,
$Z$ est quasi-compact comme image de l'application continue surjective
$Z' \to Z$ (Lemme \ref{lemma-image-quasi-compact}).
\end{proof}

\begin{lemma}
\label{lemma-inverse-limit-spectral-spaces-nonempty}
Soit $\mathcal{I}$ une catégorie cofiltrante. Soit $i \mapsto X_i$ un diagramme
d'espaces spectraux tel que, pour $a : j \to i$ dans $\mathcal{I}$,
l'application correspondante $f_a : X_j \to X_i$ soit spectrale.
\begin{enumerate}
\item Si l'on se donne des parties non vides $Z_i \subset X_i$ fermées pour
la topologie constructible telles que $f_a(Z_j) \subset Z_i$ pour tout
$a : j \to i$ dans $\mathcal{I}$, alors $\lim Z_i$ est non vide.
\item Si chaque $X_i$ est non vide, alors $X = \lim X_i$ est non vide.
\end{enumerate}
\end{lemma}

\begin{proof}
Notons $X'_i$ l'espace $X_i$ muni de la topologie constructible
et $Z'_i$ le sous-espace correspondant de $X'_i$.
Soit $a : j \to i$ un morphisme de $\mathcal{I}$. Comme $f_a$ est spectrale,
elle définit une application continue $f_a : X'_j \to X'_i$. Ainsi,
$f_a|_{Z_j} : Z'_j \to Z'_i$ est une application continue entre espaces
quasi-compacts séparés (d'après les
Lemmes \ref{lemma-constructible-hausdorff-quasi-compact} et
\ref{lemma-closed-in-quasi-compact}). D'après le
Lemme \ref{lemma-nonempty-limit}, l'espace $\lim Z'_i$ est non vide.
Puisque $\lim Z'_i = \lim Z_i$ comme ensembles, on conclut.
\end{proof}

\begin{lemma}
\label{lemma-inverse-limit-spectral-spaces-equal}
Soit $\mathcal{I}$ une catégorie cofiltrante. Soit $i \mapsto X_i$ un diagramme
d'espaces spectraux tel que, pour $a : j \to i$ dans $\mathcal{I}$,
l'application correspondante $f_a : X_j \to X_i$ soit spectrale. Soit
$X = \lim X_i$, avec pour projections $p_i : X \to X_i$. Soient
$i \in \Ob(\mathcal{I})$ et $E, F \subset X_i$ des parties telles que $E$ soit
fermée et $F$ ouverte pour la topologie constructible.
Alors $p_i^{-1}(E) \subset p_i^{-1}(F)$ si et seulement s'il existe un morphisme
$a : j \to i$ dans $\mathcal{I}$ tel que
$f_a^{-1}(E) \subset f_a^{-1}(F)$.
\end{lemma}

\begin{proof}
Observons que
$$
p_i^{-1}(E) \setminus p_i^{-1}(F) =
\lim_{a : j \to i} f_a^{-1}(E) \setminus f_a^{-1}(F)
$$
Puisque $f_a$ est une application spectrale, elle est continue pour la
topologie constructible ; la partie
$f_a^{-1}(E) \setminus f_a^{-1}(F)$ est donc fermée pour la topologie
constructible. Le
Lemme \ref{lemma-inverse-limit-spectral-spaces-nonempty}
montre alors que le membre de gauche est non vide si et seulement si chacun
des espaces du membre de droite est non vide.
\end{proof}

\begin{lemma}
\label{lemma-inverse-limit-spectral-spaces-constructible}
Soit $\mathcal{I}$ une catégorie cofiltrante. Soit $i \mapsto X_i$ un diagramme
d'espaces spectraux tel que, pour $a : j \to i$ dans $\mathcal{I}$,
l'application correspondante $f_a : X_j \to X_i$ soit spectrale. Soit
$X = \lim X_i$, avec pour projections $p_i : X \to X_i$. Soit
$E \subset X$ une partie constructible. Il existe alors un
$i \in \Ob(\mathcal{I})$ et une partie constructible $E_i \subset X_i$ tels que
$p_i^{-1}(E_i) = E$. Si $E$ est ouverte, resp.\ fermée, on peut choisir $E_i$
ouverte, resp.\ fermée.
\end{lemma}

\begin{proof}
Supposons que $E$ soit un ouvert quasi-compact de $X$. D'après le
Lemme \ref{lemma-describe-limits}, on peut écrire $E = p_i^{-1}(U_i)$
pour un certain $i$ et un certain ouvert $U_i \subset X_i$.
Écrivons $U_i = \bigcup U_{i, \alpha}$ comme réunion d'ouverts quasi-compacts.
Comme $E$ est quasi-compact, on peut trouver $\alpha_1, \ldots, \alpha_n$
tels que $E = p_i^{-1}(U_{i, \alpha_1} \cup \ldots \cup U_{i, \alpha_n})$.
Ainsi, $E_i = U_{i, \alpha_1} \cup \ldots \cup U_{i, \alpha_n}$ convient.

\medskip\noindent
Supposons que $E$ soit une partie constructible fermée. Alors $E^c$ est un
ouvert quasi-compact. Ainsi, $E^c = p_i^{-1}(F_i)$ pour un certain $i$ et
un ouvert quasi-compact $F_i \subset X_i$, d'après le résultat du paragraphe
précédent. Alors $E = p_i^{-1}(F_i^c)$, comme voulu.

\medskip\noindent
Si $E$ est quelconque, on peut écrire
$E = \bigcup_{l = 1, \ldots, n} U_l \cap Z_l$, où $U_l$ est un ouvert
constructible et $Z_l$ un fermé constructible. D'après les paragraphes
précédents, on peut écrire $U_l = p_{i_l}^{-1}(U_{l, i_l})$
et $Z_l = p_{j_l}^{-1}(Z_{l, j_l})$,
où $U_{l, i_l} \subset X_{i_l}$ est un ouvert constructible et
$Z_{l, j_l} \subset X_{j_l}$ un fermé constructible. Comme $\mathcal{I}$ est
cofiltrante, on peut choisir un objet $k$ de $\mathcal{I}$ et des morphismes
$a_l : k \to i_l$ et $b_l : k \to j_l$. En prenant alors
$E_k = \bigcup_{l = 1, \ldots, n}
f_{a_l}^{-1}(U_{l, i_l}) \cap f_{b_l}^{-1}(Z_{l, j_l})$,
on obtient une partie constructible de $X_k$ dont l'image réciproque dans $X$
est $E$.
\end{proof}

\begin{lemma}
\label{lemma-directed-inverse-limit-spectral-spaces}
Soit $\mathcal{I}$ une catégorie d'indices cofiltrante.
Soit $i \mapsto X_i$ un diagramme d'espaces spectraux tel
que, pour $a : j \to i$ dans $\mathcal{I}$, l'application correspondante
$f_a : X_j \to X_i$ soit spectrale. Alors la limite projective
$X = \lim X_i$ est un espace topologique spectral, et les applications de
projection $p_i : X \to X_i$ sont spectrales.
\end{lemma}

\begin{proof}
La limite $X = \lim X_i$ existe (Lemme \ref{lemma-limits})
et est quasi-compacte d'après le
Lemme \ref{lemma-inverse-limit-spectral-spaces-quasi-compact}.

\medskip\noindent
Notons $p_i : X \to X_i$ la projection.
Comme $\mathcal{I}$ est cofiltrante, on peut appliquer le
Lemme \ref{lemma-describe-limits}.
Une base de la topologie de $X$ est donc donnée par les ouverts
$p_i^{-1}(U_i)$, où $U_i \subset X_i$ est ouvert. Comme une base de
la topologie de $X_i$ est donnée par les ouverts quasi-compacts, on conclut
qu'une base de la topologie de $X$ est donnée par les $p_i^{-1}(U_i)$,
où $U_i \subset X_i$ est un ouvert quasi-compact. Un argument formel
montre que
$$
p_i^{-1}(U_i) = \lim_{a : j \to i} f_a^{-1}(U_i)
$$
comme espaces topologiques. Puisque chaque $f_a$ est spectrale, les parties
$f_a^{-1}(U_i)$ sont fermées pour la topologie constructible de $X_j$ ;
par conséquent, $p_i^{-1}(U_i)$ est quasi-compact d'après le
Lemme \ref{lemma-inverse-limit-spectral-spaces-quasi-compact}.
Ainsi, $X$ possède une base de sa topologie formée d'ouverts quasi-compacts.

\medskip\noindent
Tout ouvert quasi-compact $U$ de $X$ est de la forme $U = p_i^{-1}(U_i)$
pour un certain $i$ et un certain ouvert quasi-compact $U_i \subset X_i$
(voir le Lemme \ref{lemma-inverse-limit-spectral-spaces-constructible}).
Étant donnés des ouverts quasi-compacts $U_i \subset X_i$ et
$U_j \subset X_j$, on a
$p_i^{-1}(U_i) \cap p_j^{-1}(U_j) = p_k^{-1}(U_k)$ pour un certain $k$
et un certain ouvert quasi-compact $U_k \subset X_k$. En effet, choisissons
$k$ et des morphismes $k \to i$ et $k \to j$, et prenons pour $U_k$
l'intersection des images réciproques de $U_i$ et de $U_j$ dans $X_k$.
On voit ainsi que l'intersection de deux ouverts quasi-compacts de $X$ est
un ouvert quasi-compact.

\medskip\noindent
Enfin, soit $Z \subset X$ une partie fermée irréductible. Alors
$p_i(Z) \subset X_i$ est irréductible, et par conséquent
$Z_i = \overline{p_i(Z)}$ possède un unique point générique $\xi_i$
(puisque $X_i$ est un espace spectral). On a alors
$f_a(\xi_j) = \xi_i$ pour $a : j \to i$ dans $\mathcal{I}$, car
$\overline{f_a(Z_j)} = Z_i$. Ainsi, $\xi = \lim \xi_i$ est un point de $X$.
Nous affirmons que $\xi \in Z$. Sinon, on pourrait trouver un ouvert
quasi-compact contenant $\xi$ et disjoint de $Z$. Il serait de la forme
$p_i^{-1}(U_i)$ pour un certain $i$ et un certain ouvert quasi-compact
$U_i \subset X_i$. On aurait alors $\xi_i \in U_i$, mais
$p_i(Z) \cap U_i = \emptyset$, ce qui contredit
$\xi_i \in \overline{p_i(Z)}$. Ainsi, $\xi \in Z$, et donc
$\overline{\{\xi\}} \subset Z$. Réciproquement, tout $z \in Z$ appartient
à l'adhérence de $\xi$. En effet, étant donné un voisinage ouvert quasi-compact
$U$ de $z$, écrivons $U = p_i^{-1}(U_i)$ pour un certain $i$
et un certain ouvert quasi-compact $U_i \subset X_i$. On voit que
$p_i(z) \in U_i$, donc que $\xi_i \in U_i$, et donc que $\xi \in U$. Ainsi,
$\xi$ est un point générique de $Z$. Nous omettons de démontrer que $\xi$ est
l'unique point générique de $Z$ (indication : montrer qu'un second point
générique est nécessairement égal à $\xi$ en montrant qu'il a pour image
$\xi_i$ dans $X_i$, puisque $X_i$ est spectral et que l'irréductible $Z_i$
possède un unique point générique). Cela achève la démonstration.
\end{proof}

\begin{lemma}
\label{lemma-descend-opens}
Soit $\mathcal{I}$ une catégorie d'indices cofiltrante.
Soit $i \mapsto X_i$ un diagramme d'espaces spectraux tel
que, pour $a : j \to i$ dans $\mathcal{I}$, l'application correspondante
$f_a : X_j \to X_i$ soit spectrale. Posons $X = \lim X_i$ et notons
$p_i : X \to X_i$ la projection.
\begin{enumerate}
\item Pour tout ouvert quasi-compact $U \subset X$,
il existe un $i \in \Ob(\mathcal{I})$ et un ouvert quasi-compact
$U_i \subset X_i$ tels que $p_i^{-1}(U_i) = U$.
\item Étant donnés des ouverts quasi-compacts $U_i \subset X_i$ et
$U_j \subset X_j$ tels que $p_i^{-1}(U_i) \subset p_j^{-1}(U_j)$,
il existe $k \in \Ob(\mathcal{I})$ et des morphismes
$a : k \to i$ et $b : k \to j$ tels que
$f_a^{-1}(U_i) \subset f_b^{-1}(U_j)$.
\item Si $U_i, U_{1, i}, \ldots, U_{n, i} \subset X_i$ sont des ouverts
quasi-compacts et si
$p_i^{-1}(U_i) = p_i^{-1}(U_{1, i}) \cup \ldots \cup p_i^{-1}(U_{n, i})$,
alors
$f_a^{-1}(U_i) = f_a^{-1}(U_{1, i}) \cup \ldots \cup f_a^{-1}(U_{n, i})$
pour un certain morphisme $a : j \to i$ dans $\mathcal{I}$.
\item Même énoncé qu'en (3), mais pour les intersections.
\end{enumerate}
\end{lemma}

\begin{proof}
L'assertion (1) est un cas particulier du
Lemme \ref{lemma-inverse-limit-spectral-spaces-constructible}.
L'assertion (2) est un cas particulier du
Lemme \ref{lemma-inverse-limit-spectral-spaces-equal},
car les ouverts quasi-compacts sont à la fois ouverts et fermés pour la
topologie constructible. Les assertions (3) et (4) résultent formellement de
(1) et (2), ainsi que du fait que la formation des images réciproques
commute aux réunions et aux intersections.
\end{proof}

\begin{lemma}
\label{lemma-make-spectral-space}
Soit $W$ une partie d'un espace spectral $X$. Les conditions suivantes sont
équivalentes :
\begin{enumerate}
\item $W$ est une intersection de parties constructibles et est stable par
générisation,
\item $W$ est quasi-compacte et stable par générisation,
\item il existe une partie quasi-compacte $E \subset X$ telle que
$W$ soit l'ensemble des points qui se spécialisent en un point de $E$,
\item $W$ est une intersection d'ouverts quasi-compacts,
\item
\label{item-intersection-quasi-compact-open}
il existe un ensemble non vide $I$ et des ouverts quasi-compacts
$U_i \subset X$, $i \in I$,
tels que $W = \bigcap U_i$ et que, pour tous $i, j \in I$, il existe un
$k \in I$ tel que $U_k \subset U_i \cap U_j$.
\end{enumerate}
Dans ce cas, (a) $W$ est un espace spectral, (b) $W = \lim U_i$
comme espaces topologiques et (c) pour tout ouvert $U$ contenant $W$,
il existe un $i$ tel que $U_i \subset U$.
\end{lemma}

\begin{proof}
Soit $W \subset X$ vérifiant (1). Alors $W$ est fermée pour la topologie
constructible, donc quasi-compacte pour la topologie constructible (d'après les
Lemmes \ref{lemma-constructible-hausdorff-quasi-compact} et
\ref{lemma-closed-in-quasi-compact}), donc quasi-compacte pour la topologie
de $X$ (car les ouverts de $X$ sont ouverts pour la topologie constructible).
Ainsi, (2) est satisfaite.

\medskip\noindent
Il est clair que (2) implique (3) en prenant $E = W$.

\medskip\noindent
Soit $X$ un espace spectral et soit $E \subset W$ comme en (3).
Puisque tout point de $W$ se spécialise en un point de $E$, on voit qu'un
ouvert de $W$ qui contient $E$ est égal à $W$. Par conséquent, puisque $E$
est quasi-compacte, $W$ l'est également.
Si $x \in X$, $x \not \in W$, alors $Z = \overline{\{x\}}$ est
disjoint de $W$. Comme $W$ est quasi-compacte, on peut trouver un
ouvert quasi-compact $U$ tel que $W \subset U$ et $U \cap Z = \emptyset$.
On conclut que (4) est satisfaite.

\medskip\noindent
Si $W = \bigcap_{j \in J} U_j$, alors, en prenant pour $I$ l'ensemble des
parties finies de $J$ et $U_i = U_{j_1} \cap \ldots \cap U_{j_r}$
pour $i = \{j_1, \ldots, j_r\}$, on voit que (4) implique (5). Il est immédiat
que (5) implique (1).

\medskip\noindent
Soient $I$ et $U_i$ comme en (5).
Puisque $W = \bigcap U_i$, on a $W = \lim U_i$ par la propriété universelle
des limites. Alors $W$ est un espace spectral d'après le
Lemme \ref{lemma-directed-inverse-limit-spectral-spaces}.
Soit $U \subset X$ un voisinage ouvert de $W$.
Alors les $E_i = U_i \cap (X \setminus U)$ forment une famille de parties
constructibles de l'espace spectral $Z = X \setminus U$
d'intersection vide. En utilisant le fait que la topologie spectrale sur $Z$
est quasi-compacte (Lemme \ref{lemma-constructible-hausdorff-quasi-compact}),
on déduit du
Lemme \ref{lemma-intersection-closed-in-quasi-compact}
que $E_i = \emptyset$ pour un certain $i$.
\end{proof}

\begin{lemma}
\label{lemma-make-spectral-space-minus}
Soit $X$ un espace spectral. Soit $E \subset X$ une partie constructible.
Soit $W \subset X$ l'ensemble des points de $X$ qui se spécialisent
en un point de $E$. Alors $W \setminus E$ est un espace spectral.
Si $W = \bigcap U_i$, avec les $U_i$ comme dans le
Lemme \ref{lemma-make-spectral-space}
(\ref{item-intersection-quasi-compact-open}),
alors $W \setminus E = \lim (U_i \setminus E)$.
\end{lemma}

\begin{proof}
Puisque $E$ est constructible, elle est quasi-compacte, et le
Lemme \ref{lemma-make-spectral-space} s'applique donc à $W$.
Si $E$ est constructible, alors $E$ est constructible dans $U_i$
pour tout $i \in I$. Ainsi, $U_i \setminus E$
est spectral d'après le Lemme \ref{lemma-spectral-sub}.
Puisque $W \setminus E = \bigcap (U_i \setminus E)$, on a
$W \setminus E = \lim U_i \setminus E$ par la propriété universelle des limites.
Alors $W \setminus E$ est un espace spectral d'après le
Lemme \ref{lemma-directed-inverse-limit-spectral-spaces}.
\end{proof}






\section{Compactification de Stone-{\v C}ech}
\label{section-stone-cech}

\noindent
La compactification de Stone-{\v C}ech d'un espace topologique $X$
est une application $X \to \beta(X)$ de $X$ vers un espace séparé
quasi-compact $\beta(X)$, universelle parmi de telles applications. Nous
démontrons son existence par un argument standard utilisant le lemme élémentaire
suivant.

\begin{lemma}
\label{lemma-dense-image}
Soit $f : X \to Y$ une application continue entre espaces topologiques. Supposons
que $f(X)$ soit dense dans $Y$ et que $Y$ soit séparé. Alors le cardinal de
$Y$ est inférieur ou égal à celui de $P(P(X))$, où $P$ désigne l'opération
qui associe à un ensemble l'ensemble de ses parties.
\end{lemma}

\begin{proof}
Soit $S = f(X) \subset Y$. Soit $\mathcal{D}$ l'ensemble de tous les
fermés réguliers de $Y$, c'est-à-dire des parties $D \subset Y$ qui
sont égales à l'adhérence de leur intérieur. Notons que l'adhérence d'une
partie ouverte de $Y$ est un fermé régulier. Pour $y \in Y$, considérons
l'ensemble
$$
I_y = \{T \subset S \mid \text{ il existe }
D \in \mathcal{D}\text{ tel que }T = S \cap D\text{ et }y \in D\}.
$$
Comme $S$ est dense dans $Y$, pour tout fermé régulier $D$, on voit
que $S \cap D$ est dense dans $D$. Par conséquent, si
$D \cap S = D' \cap S$ pour $D, D' \in \mathcal{D}$, alors
$D = D'$. Ainsi, $I_y = I_{y'}$ implique $y = y'$, car la condition de
séparation assure que l'on peut trouver un fermé régulier contenant $y$
mais pas $y'$. Le résultat s'ensuit.
\end{proof}

\noindent
Soit $X$ un espace topologique. D'après le Lemme \ref{lemma-dense-image}, il
existe un ensemble $I$
de classes d'isomorphisme d'applications continues $f : X \to Y$ à image
dense, où $Y$ est séparé et quasi-compact. Pour $i \in I$, choisissons
un représentant $f_i : X \to Y_i$. Considérons l'application
$$
\prod f_i : X \longrightarrow \prod\nolimits_{i \in I} Y_i
$$
et notons $\beta(X)$ l'adhérence de son image. Puisque chaque $Y_i$ est
séparé, $\beta(X)$ l'est aussi. Puisque chaque $Y_i$ est quasi-compact,
$\beta(X)$ l'est aussi (utiliser le Théorème \ref{theorem-tychonov} et le
Lemme \ref{lemma-closed-in-quasi-compact}).

\medskip\noindent
Montrons que l'application canonique $X \to \beta(X)$ satisfait à la propriété
universelle relativement aux applications vers des espaces séparés
quasi-compacts. Soit donc $f : X \to Y$ un tel morphisme. Soit
$Z \subset Y$ l'adhérence de $f(X)$. Alors $X \to Z$ est isomorphe à l'une des
applications $f_i : X \to Y_i$, disons $f_{i_0} : X \to Y_{i_0}$. Ainsi, $f$
se factorise sous la forme
$X \to \beta(X) \to \prod Y_i \to Y_{i_0} \cong Z \to Y$, comme voulu.

\begin{lemma}
\label{lemma-one-point-compactification}
Soit $X$ un espace séparé et localement quasi-compact.
Il existe une application $X \to X^*$ qui identifie $X$ à un sous-espace
ouvert d'un espace séparé quasi-compact $X^*$ tel que
$X^* \setminus X$ soit un singleton (compactification par adjonction d'un point).
En particulier, l'application $X \to \beta(X)$ identifie $X$
à un sous-espace ouvert de $\beta(X)$.
\end{lemma}

\begin{proof}
Posons $X^* = X \amalg \{\infty\}$. Déclarons qu'une partie $V$ de $X^*$ est
ouverte si $V \subset X$ est ouverte dans $X$, ou bien si $\infty \in V$ et si
$U = V \cap X$ est un ouvert de $X$ tel que $X \setminus U$ soit quasi-compact.
Nous omettons de vérifier que cela définit une topologie. Il est clair
que $X \to X^*$ identifie $X$ à un sous-espace ouvert de $X$.

\medskip\noindent
Puisque $X$ est localement quasi-compact, tout point $x \in X$ possède un
voisinage quasi-compact $x \in E \subset X$. Alors $E$
est fermé (Lemme \ref{lemma-quasi-compact-in-Hausdorff}, partie (1)), et
$V = (X \setminus E) \amalg \{\infty\}$ est un voisinage ouvert
de $\infty$ disjoint de l'intérieur de $E$. Ainsi, $X^*$ est séparé.

\medskip\noindent
Soit $X^* = \bigcup V_i$ un recouvrement ouvert. Alors, pour un certain $i$,
disons $i_0$, on a $\infty \in V_{i_0}$. Par construction,
$Z = X^* \setminus V_{i_0}$ est quasi-compact. Le recouvrement
$Z \subset \bigcup_{i \not = i_0} Z \cap V_i$ admet donc un raffinement fini,
ce qui implique que le recouvrement donné de $X^*$ admet un raffinement fini.
Ainsi, $X^*$ est quasi-compact.

\medskip\noindent
L'application $X \to X^*$ se factorise sous la forme
$X \to \beta(X) \to X^*$ par la propriété universelle de la compactification
de Stone-{\v C}ech. Soit
$\varphi : \beta(X) \to X^*$ cette factorisation.
Alors $X \to \varphi^{-1}(X)$ est une section de
$\varphi^{-1}(X) \to X$, et possède donc une image fermée
(Lemme \ref{lemma-section-closed}).
Comme l'image de $X \to \beta(X)$ est dense, on conclut que
$X = \varphi^{-1}(X)$.
\end{proof}








\section{Espaces extrêmement discontinus}
\label{section-extremally-disconnected}

\noindent
Le contenu de cette section est tiré de \cite{Gleason}
(avec une légère modification comme dans \cite{Rainwater}).
Dans l'article de Gleason, il est démontré que, dans la catégorie des espaces
séparés quasi-compacts, les « objets projectifs » sont exactement les espaces
extrêmement discontinus.

\begin{definition}
\label{definition-extremally-disconnected}
On dit qu'un espace topologique $X$ est {\it extrêmement discontinu}
si l'adhérence de toute partie ouverte de $X$ est ouverte.
\end{definition}

\noindent
Si $X$ est séparé et extrêmement discontinu, alors $X$ est totalement
discontinu (cela est faux en général). Si $X$ est quasi-compact, séparé et
extrêmement discontinu, alors $X$ est profini d'après le
lemme \ref{lemma-profinite}, mais la réciproque est fausse en général.
Par exemple, l'anneau des entiers $p$-adiques
$\mathbf{Z}_p = \lim \mathbf{Z}/p^n\mathbf{Z}$ est un espace profini
qui n'est pas extrêmement discontinu. En effet, si $U \subset \mathbf{Z}_p$
est l'ensemble des éléments non nuls dont la valuation est paire, alors $U$ est
ouvert, mais son adhérence est $U \cup \{0\}$, qui n'est pas ouvert.

\begin{lemma}
\label{lemma-image-open-technical}
Soit $f : X \to Y$ une application continue entre espaces topologiques.
Supposons que $f$ soit surjective et que $f(E) \not = Y$ pour toute partie
fermée propre $E \subset X$. Alors, pour tout ouvert $U \subset X$, la partie
$f(U)$ est contenue dans l'adhérence de $Y \setminus f(X \setminus U)$.
\end{lemma}

\begin{proof}
Choisissons $y \in f(U)$ et un voisinage ouvert quelconque $V \subset Y$ de
$y$. Nous allons montrer que $V$ rencontre $Y \setminus f(X \setminus U)$.
Remarquons que $W = U \cap f^{-1}(V)$ est un ouvert non vide de $X$ ; ainsi,
$f(X \setminus W) \not = Y$. Prenons $y' \in Y$ tel que
$y' \not \in f(X \setminus W)$. Il est élémentaire de vérifier que $y' \in V$
et que $y' \in Y \setminus f(X \setminus U)$.
\end{proof}

\begin{lemma}
\label{lemma-intersection-empty}
Soit $X$ un espace extrêmement discontinu.
Si $U, V \subset X$ sont des ouverts disjoints, alors
$\overline{U}$ et $\overline{V}$ sont eux aussi disjoints.
\end{lemma}

\begin{proof}
Par hypothèse, $\overline{U}$ est ouvert ; ainsi,
$V \cap \overline{U}$ est ouvert et disjoint de $U$, donc vide, puisque
$\overline{U}$ est l'intersection de toutes les parties fermées de $X$ qui
contiennent $U$. Cela signifie que l'ouvert
$\overline{V} \cap \overline{U}$ ne rencontre pas $V$ ; il est donc vide par
le même argument.
\end{proof}

\begin{lemma}
\label{lemma-isomorphism}
Soit $f : X \to Y$ une application continue entre espaces topologiques
séparés quasi-compacts. Si $Y$ est extrêmement discontinu, si $f$ est
surjective et si $f(Z) \not = Y$ pour toute partie fermée propre $Z$ de $X$,
alors $f$ est un homéomorphisme.
\end{lemma}

\begin{proof}
D'après le lemme \ref{lemma-bijective-map}, il suffit de montrer que $f$ est
injective. Supposons que $x, x' \in X$ soient des points distincts tels que
$y = f(x) = f(x')$. Choisissons des voisinages ouverts disjoints
$U, U' \subset X$ respectifs de $x, x'$. Remarquons que $f$ est fermée
(lemme \ref{lemma-closed-map}) ; ainsi, $T = f(X \setminus U)$ et
$T' = f(X \setminus U')$ sont fermés dans $Y$. Puisque $X$ est la réunion de
$X \setminus U$ et de $X \setminus U'$, on voit que $Y = T \cup T'$.
D'après le lemme \ref{lemma-image-open-technical}, on voit que $y$ appartient
à l'adhérence de $Y \setminus T$ et à celle de $Y \setminus T'$. D'autre part,
d'après le lemme \ref{lemma-intersection-empty}, leur intersection est vide.
Nous obtenons ainsi la contradiction souhaitée.
\end{proof}

\begin{lemma}
\label{lemma-find-compact-subset}
Soit $f : X \to Y$ une application continue surjective entre espaces
topologiques séparés quasi-compacts. Il existe une partie quasi-compacte
$E \subset X$ telle que $f(E) = Y$, mais $f(E') \not = Y$ pour toute partie
fermée propre $E' \subset E$.
\end{lemma}

\begin{proof}
Nous utiliserons sans plus le mentionner le fait que les parties quasi-compactes
de $X$ sont exactement les parties fermées
(lemme \ref{lemma-closed-in-compact}).
Considérons la collection $\mathcal{E}$ de toutes les parties quasi-compactes
$E \subset X$ telles que $f(E) = Y$, ordonnée par inclusion. Nous allons
utiliser le lemme de Zorn pour montrer que $\mathcal{E}$ possède un élément
minimal. Pour cela, il suffit de montrer que, pour toute famille totalement
ordonnée $E_\lambda$ d'éléments de $\mathcal{E}$, l'intersection
$\bigcap E_\lambda$ appartient à $\mathcal{E}$. Elle est quasi-compacte car
elle est fermée. Pour tout $y \in Y$, les parties
$E_\lambda \cap f^{-1}(\{y\})$ sont non vides et fermées ; ainsi, l'intersection
$\bigcap E_\lambda \cap f^{-1}(\{y\}) = \bigcap (E_\lambda \cap f^{-1}(\{y\}))$
est non vide d'après le
lemme \ref{lemma-intersection-closed-in-quasi-compact}.
Cela achève la démonstration.
\end{proof}

\begin{proposition}
\label{proposition-projective-in-category-hausdorff-qc}
Soit $X$ un espace topologique séparé et quasi-compact.
Les assertions suivantes sont équivalentes :
\begin{enumerate}
\item $X$ est extrêmement discontinu,
\item pour toute application continue surjective $f : Y \to X$, où $Y$ est
séparé et quasi-compact, il existe une section continue, et
\item pour tout diagramme commutatif en traits pleins
$$
\xymatrix{
& Y \ar[d] \\
X \ar@{..>}[ru] \ar[r] & Z
}
$$
formé d'applications continues entre espaces séparés quasi-compacts, dans lequel
$Y \to Z$ est surjective, il existe dans la catégorie des espaces topologiques
une flèche pointillée qui rend le diagramme commutatif.
\end{enumerate}
\end{proposition}

\begin{proof}
Il est clair que (3) implique (2). D'autre part, si (2) est vérifiée et si
$X \to Z$ et $Y \to Z$ sont comme en (3), alors (2) assure l'existence d'une
section de la projection $X \times_Z Y \to X$, d'où l'existence d'une flèche
pointillée convenable (détails omis). Ainsi, (3) est équivalente à (2).

\medskip\noindent
Supposons $X$ extrêmement discontinu et soit $f : Y \to X$ comme en (2).
D'après le lemme \ref{lemma-find-compact-subset}, il existe une partie
quasi-compacte $E \subset Y$ telle que $f(E) = X$, mais
$f(E') \not = X$ pour toute partie fermée propre $E' \subset E$. D'après le
lemme \ref{lemma-isomorphism}, $f|_E : E \to X$ est un homéomorphisme, dont
l'inverse donne la section souhaitée.

\medskip\noindent
Supposons (2). Soit $U \subset X$ un ouvert de complémentaire $Z$.
Considérons la surjection continue $f : \overline{U} \amalg Z \to X$.
Soit $\sigma$ une section. Alors $\overline{U} = \sigma^{-1}(\overline{U})$
est ouverte. Ainsi, $X$ est extrêmement discontinu.
\end{proof}

\begin{lemma}
\label{lemma-rainwater}
Soit $f : X \to X$ une application continue surjective d'un espace
topologique séparé dans lui-même. Si $f$ n'est pas $\text{id}_X$, alors il
existe une partie fermée propre $E \subset X$ telle que $X = E \cup f(E)$.
\end{lemma}

\begin{proof}
Choisissons $p \in X$ tel que $f(p) \not = p$. Choisissons des voisinages
ouverts disjoints $p \in U$, $f(p) \in V$ et posons
$E = X \setminus U \cap f^{-1}(V)$. Alors $p \not \in E$ ; ainsi, $E$ est une
partie fermée propre. Si $x \in X$, alors soit $x \in E$, soit, dans le cas
contraire, $x \in U \cap f^{-1}(V)$. Écrivons $x = f(y)$, ce qui est possible
puisque $f$ est surjective. Si $y \in U \cap f^{-1}(V)$, alors on aurait
$x = f(y) \in V$, ce qui contredit $x \in U$. Ainsi, $y \in E$ et
$x \in f(E)$.
\end{proof}

\begin{example}
\label{example-stone-Cech-discrete}
La proposition \ref{proposition-projective-in-category-hausdorff-qc} permet de
voir que la compactification de Stone-{\v C}ech $\beta(X)$ d'un espace discret
$X$ est extrêmement discontinue. En effet, soit $f : Y \to \beta(X)$ une
surjection continue, où $Y$ est quasi-compact et séparé. Nous pouvons alors
relever l'application $X \to \beta(X)$ en une application continue (!)
$X \to Y$, puisque $X$ est discret. Par la propriété universelle de la
compactification de Stone-{\v C}ech, nous obtenons une factorisation
$X \to \beta(X) \to Y$. Puisque $\beta(X) \to Y \to \beta(X)$ est égale à
l'identité sur la partie dense $X$, nous concluons que nous avons obtenu une
section.
En particulier, nous concluons que la compactification de Stone-{\v C}ech d'un
espace discret est totalement discontinue, donc profinie
(voir la discussion qui suit la
définition \ref{definition-extremally-disconnected} et le
lemme \ref{lemma-profinite}).
\end{example}

\noindent
Grâce aux espaces extrêmement discontinus fournis par
l'exemple \ref{example-stone-Cech-discrete}, nous pouvons démontrer que tout
espace séparé quasi-compact possède une « enveloppe projective » dans la
catégorie des espaces séparés quasi-compacts.

\begin{lemma}
\label{lemma-existence-projective-cover}
\begin{slogan}
Tout espace séparé quasi-compact possède une enveloppe
extrêmement discontinue canonique
\end{slogan}
Soit $X$ un espace séparé quasi-compact.
Il existe une surjection continue $X' \to X$ telle que $X'$ soit
quasi-compact, séparé et extrêmement discontinu.
Si l'on exige qu'aucune partie fermée propre de $X'$ n'ait pour image $X$,
alors $X'$ est unique à isomorphisme près.
\end{lemma}

\begin{proof}
Soit $Y = X$, mais muni de la topologie discrète. Soit $X' = \beta(Y)$.
L'application continue $Y \to X$ se factorise en $Y \to X' \to X$. Cela
démontre la première assertion du lemme d'après
l'exemple \ref{example-stone-Cech-discrete}.

\medskip\noindent
D'après le lemme \ref{lemma-find-compact-subset}, nous pouvons trouver une
partie quasi-compacte $E \subset X'$ dont l'image est $X$ et telle qu'aucune
partie fermée propre de $E$ n'ait pour image $X$.
Puisque $X'$ est extrêmement discontinu, il existe une application continue
$f : X' \to E$ au-dessus de $X$
(proposition \ref{proposition-projective-in-category-hausdorff-qc}).
En composant $f$ avec l'application $E \to X'$, on obtient l'endomorphisme
continu $f|_E : E \to E$. Remarquons que $f|_E$ doit être surjectif, car sinon
son image serait une partie fermée propre d'image $X$.
Ainsi, $f|_E$ doit être $\text{id}_E$, car sinon le
lemme \ref{lemma-rainwater} montrerait que $E$ n'est pas minimal.
Par conséquent, $\text{id}_E$ se factorise par l'espace extrêmement discontinu
$X'$. Un argument catégorique formel (utilisant la caractérisation de la
proposition \ref{proposition-projective-in-category-hausdorff-qc}) montre que
$E$ est extrêmement discontinu.

\medskip\noindent
Pour démontrer l'unicité, supposons donnée une seconde enveloppe minimale
$X'' \to X$. Par la propriété de relèvement démontrée dans la
proposition \ref{proposition-projective-in-category-hausdorff-qc}, nous pouvons
trouver une application continue $g : X' \to X''$ au-dessus de $X$.
Remarquons que $g$ est une application fermée
(lemme \ref{lemma-closed-map}). Ainsi, $g(X') \subset X''$ est une partie
fermée d'image $X$, et nous concluons que
$g(X') = X''$ par minimalité de $X''$. D'autre part, si $E \subset X'$ est une
partie fermée
propre, alors $g(E) \not = X''$, car $E$ n'a pas pour image $X$ par minimalité
de $X'$. D'après le lemme \ref{lemma-isomorphism}, $g$ est un
isomorphisme.
\end{proof}

\begin{remark}
\label{remark-size-projective-cover}
Soit $X$ un espace séparé quasi-compact. Soit $\kappa$ un cardinal infini
supérieur ou égal au cardinal de $X$. Alors le cardinal de l'enveloppe minimale
séparée, quasi-compacte et extrêmement discontinue
$X' \to X$ (lemme \ref{lemma-existence-projective-cover})
est au plus $2^{2^\kappa}$. En effet, choisissons une partie $S \subset X'$
qui s'envoie bijectivement sur $X$. Par minimalité de $X'$, la partie
$S$ est dense dans $X'$. Ainsi, $|X'| \leq 2^{2^\kappa}$ d'après le
lemme \ref{lemma-dense-image}.
\end{remark}






\section{Divers}
\label{section-miscellany}


\noindent
Le lemme suivant s'applique à l'espace topologique sous-jacent à un schéma
quasi-séparé.

\begin{lemma}
\label{lemma-topology-quasi-separated-scheme}
Soit $X$ un espace topologique qui
\begin{enumerate}
\item possède une base de la topologie formée d'ouverts quasi-compacts, et
\item est tel que l'intersection de deux ouverts quasi-compacts
quelconques est quasi-compacte.
\end{enumerate}
Alors
\begin{enumerate}
\item $X$ est localement quasi-compact,
\item tout ouvert quasi-compact $U \subset X$ est rétrocompact,
\item tout ouvert quasi-compact $U \subset X$ possède un système cofinal de
recouvrements ouverts $\mathcal{U} : U = \bigcup_{j\in J} U_j$ où $J$ est fini
et où tous les $U_j$ et $U_j \cap U_{j'}$ sont quasi-compacts,
\item en ajouter ici.
\end{enumerate}
\end{lemma}

\begin{proof}
La démonstration est omise.
\end{proof}

\begin{definition}
\label{definition-isolated-point}
Soit $X$ un espace topologique. On dit que $x \in X$ est un
{\it point isolé} de $X$ si $\{x\}$ est ouvert dans $X$.
\end{definition}

\section{Partitions et stratifications}
\label{section-stratifications}

\noindent
Les stratifications peuvent être définies de bien des manières. Toute remarque
sur le choix des définitions adoptées dans cette section est la bienvenue.

\begin{definition}
\label{definition-paritition}
Soit $X$ un espace topologique. Une {\it partition} de $X$ est une
décomposition $X = \coprod X_i$ en parties localement fermées $X_i$.
Les $X_i$ sont appelées les {\it parties} de la partition.
Étant donné deux partitions de $X$, nous disons que l'une {\it raffine}
l'autre si les parties de l'une sont des réunions de parties de l'autre.
\end{definition}

\noindent
Tout espace topologique $X$ admet une partition en composantes connexes.
Si $X$ possède un nombre fini de composantes irréductibles
$Z_1, \ldots, Z_r$, il existe une partition qui raffine la partition en
composantes connexes et dont les parties
$X_I = \bigcap_{i \in I} Z_i \setminus (\bigcup_{i \not \in I} Z_i)$
sont indexées par les sous-ensembles $I \subset \{1, \ldots, r\}$.

\begin{definition}
\label{definition-good-stratification}
Soit $X$ un espace topologique. Une {\it bonne stratification}
de $X$ est une partition $X = \coprod X_i$ telle que, pour tous
$i, j \in I$, on ait
$$
X_i \cap \overline{X_j} \not = \emptyset
\Rightarrow
X_i \subset \overline{X_j}.
$$
\end{definition}

\noindent
Une bonne stratification $X = \coprod_{i \in I} X_i$ définit
un ordre partiel sur $I$ en posant $i \leq j$ si et seulement si
$X_i \subset \overline{X_j}$. On voit alors que
$$
\overline{X_j} = \bigcup\nolimits_{i \leq j} X_i
$$
Cependant, en géométrie algébrique, il arrive souvent que l'on sache seulement
que le membre de gauche est contenu dans le membre de droite dans la dernière
formule affichée. Cela conduit à la définition suivante.

\begin{definition}
\label{definition-stratification}
Soit $X$ un espace topologique. Une {\it stratification} de $X$ est
la donnée d'une partition $X = \coprod_{i \in I} X_i$ et d'un ordre partiel
sur $I$ tels que, pour tout $j \in I$, on ait
$$
\overline{X_j} \subset \bigcup\nolimits_{i \leq j} X_i
$$
Les parties $X_i$ sont appelées les {\it strates} de la stratification.
\end{definition}

\noindent
Nous imposons souvent des conditions supplémentaires à la stratification.
Par exemple, les stratifications sont particulièrement agréables lorsqu'elles
sont {\it localement finies}, ce qui signifie que tout point possède un
voisinage ne rencontrant qu'un nombre fini de strates. Plus généralement,
nous introduisons la définition suivante.

\begin{definition}
\label{definition-locally-finite}
Soit $X$ un espace topologique. Soit $I$ un ensemble et, pour $i \in I$,
soit $E_i \subset X$ une partie. Nous disons que la famille
$\{E_i\}_{i \in I}$ est {\it localement finie} si, pour tout $x \in X$,
il existe un voisinage ouvert $U$ de $x$ tel que
$\{i \in I | E_i \cap U \not = \emptyset\}$ soit fini.
\end{definition}


\begin{remark}
\label{remark-locally-finite-stratification}
Étant donné une stratification localement finie $X = \coprod X_i$ d'un
espace topologique $X$, on obtient une famille de parties fermées
$Z_i = \bigcup_{j \leq i} X_j$ de $X$, indexée par $I$, telle que
$$
Z_i \cap Z_j = \bigcup\nolimits_{k \leq i, j} Z_k
$$
Réciproquement, supposons données des parties fermées $Z_i \subset X$
indexées par un ensemble ordonné $I$, telles que $X = \bigcup Z_i$, que tout
point possède un voisinage ne rencontrant qu'un nombre fini de $Z_i$ et que la
formule affichée soit vérifiée. On obtient alors une stratification localement
finie de $X$ en posant $X_i = Z_i \setminus \bigcup_{j < i} Z_j$.
\end{remark}

\begin{lemma}
\label{lemma-partition-refined-by-stratification}
Soit $X$ un espace topologique et soit $X = \coprod X_i$ une partition finie
de $X$. Il existe alors une stratification finie de $X$ qui la raffine.
\end{lemma}

\begin{proof}
Posons $T_i = \overline{X_i}$ et $\Delta_i = T_i \setminus X_i$.
Soit $S$ l'ensemble de toutes les intersections de $T_i$ et de $\Delta_i$.
(Par exemple, $T_1 \cap T_2 \cap \Delta_4$ est un élément de $S$.)
Alors $S = \{Z_s\}$ est une famille finie de parties fermées de $X$ telle que
$Z_s \cap Z_{s'} \in S$ pour tous $s, s' \in S$. Munissons $S$ de l'ordre
partiel donné par l'inclusion. En posant ensuite
$Y_s = Z_s \setminus \bigcup_{s' < s} Z_{s'}$, on obtient la stratification
cherchée.
\end{proof}

\begin{lemma}
\label{lemma-constructible-partition-refined-by-stratification}
Soit $X$ un espace topologique. Supposons que
$X = T_1 \cup \ldots \cup T_n$ soit écrit comme réunion de parties
constructibles. Il existe une stratification finie
$X = \coprod X_i$, dont chaque $X_i$ est constructible, telle que chaque
$T_k$ soit une réunion de strates.
\end{lemma}

\begin{proof}
Par définition des parties constructibles, chaque $T_i$ s'écrit comme réunion
finie de parties $U \cap V^c$, où $U, V \subset X$ sont des ouverts
rétrocompacts. Nous pouvons donc supposer que $T_i = U_i \cap V_i^c$,
où $U_i, V_i \subset X$ sont des ouverts rétrocompacts. Soit $S$ l'ensemble
fini de parties fermées de $X$ constitué de
$\emptyset, X, U_i^c, V_i^c$ et de leurs intersections finies.
Si $Z \in S$, alors $Z$ est constructible dans $X$
(Lemme \ref{lemma-constructible}).
De plus, $Z \cap Z' \in S$ pour tous $Z, Z' \in S$.
Munissons $S$ de l'ordre partiel donné par l'inclusion. Pour $Z \in S$, posons
$X_Z = Z \setminus \bigcup_{Z' < Z,\ Z' \in S} Z'$.
On obtient ainsi une stratification $X = \coprod_{Z \in S} X_Z$
qui possède les propriétés énoncées dans le lemme.
\end{proof}

\begin{lemma}
\label{lemma-noetherian-partition-refined-by-stratification}
Soit $X$ un espace topologique noethérien. Toute partition finie
de $X$ peut être raffinée par une bonne stratification finie.
\end{lemma}

\begin{proof}
Soit $X = \coprod X_i$ une partition finie de $X$.
Soit $Z$ une composante irréductible de $X$. Puisque
$X = \bigcup \overline{X_i}$ et que l'ensemble d'indices est fini, il existe
un $i$ tel que $Z \subset \overline{X_i}$. Comme $X_i$ est localement fermé,
il s'ensuit que $Z \cap X_i$ contient un ouvert de $Z$. Ainsi,
$Z \cap X_i$ contient un ouvert $U$ de $X$ (Lemme \ref{lemma-Noetherian}).
Écrivons $X_i = U \amalg X_i^1 \amalg X_i^2$, où
$X_i^1 = (X_i \setminus U) \cap \overline{U}$ et
$X_i^2 = (X_i \setminus U) \cap \overline{U}^c$.
Pour $i' \not = i$, posons
$X_{i'}^1 = X_{i'} \cap \overline{U}$ et
$X_{i'}^2 = X_{i'} \cap \overline{U}^c$.
Alors
$$
X \setminus U = \coprod X^k_l
$$
est une partition telle que $\overline{U} \setminus U = \bigcup X_l^1$.
Remarquons que $X \setminus U$ est une partie fermée strictement contenue
dans $X$.
Par récurrence noethérienne, nous pouvons raffiner cette partition
par une bonne stratification finie
$X \setminus U = \coprod_{\alpha \in A} T_\alpha$.
Alors $X = U \amalg \coprod_{\alpha \in A} T_\alpha$ est une bonne
stratification finie de $X$ qui raffine la partition initiale.
\end{proof}







\section{Colimites d'espaces}
\label{section-colimits}

\noindent
La catégorie des espaces topologiques possède des coproduits. En effet, si
$I$ est un ensemble et si, pour $i \in I$, on se donne un espace topologique
$X_i$, on munit l'ensemble $\coprod_{i \in I} X_i$ de la {\it topologie
somme}. Une base de cette topologie est formée des parties $U_i$, où
$U_i \subset X_i$ est ouvert.

\medskip\noindent
La catégorie des espaces topologiques possède des coégalisateurs. En effet,
si $a, b : X \to Y$ sont des morphismes d'espaces topologiques, alors le
coégalisateur de $a$ et $b$ est le coégalisateur $Y/\sim$ dans la catégorie
des ensembles, muni de la topologie quotient (section \ref{section-submersive}).

\begin{lemma}
\label{lemma-colimits}
La catégorie des espaces topologiques possède des colimites, et le foncteur
d'oubli vers la catégorie des ensembles commute aux colimites.
\end{lemma}

\begin{proof}
Cela résulte de la discussion précédente et de Catégories, Lemme
\ref{categories-lemma-colimits-coproducts-coequalizers}. Une autre preuve de
l'existence des colimites est esquissée dans Catégories, Remarque
\ref{categories-remark-how-to-use-it}. Il résulte de ce qui précède que le
foncteur d'oubli commute aux colimites. On peut également le voir en utilisant
Catégories, Lemme \ref{categories-lemma-adjoint-exact}, et le fait que le
foncteur d'oubli possède un adjoint à droite, à savoir le foncteur qui associe
à un ensemble l'espace topologique chaotique (ou indiscret) correspondant.
\end{proof}





\section{Groupes, anneaux et modules topologiques}
\label{section-topological-groups}

\noindent
Cette courte section ne contient que des définitions et des propriétés
élémentaires.

\begin{definition}
\label{definition-topological-group}
Un {\it groupe topologique} est un groupe $G$ muni d'une topologie telle que
la multiplication $G \times G \to G$, $(x, y) \mapsto xy$, et l'inversion
$G \to G$, $x \mapsto x^{-1}$, soient continues.
Un {\it homomorphisme de groupes topologiques} est un homomorphisme de groupes
qui est continu.
\end{definition}

\noindent
Si $G$ est un groupe topologique et $H \subset G$ un sous-groupe, alors $H$,
muni de la topologie induite, est un groupe topologique. Si $G$ est un groupe
topologique et $G \to H$ une surjection de groupes, alors $H$, muni de la
topologie quotient, est un groupe topologique.

\begin{example}
\label{example-automorphisms-of-a-set}
Soit $E$ un ensemble. On peut munir l'ensemble de ses applications dans
lui-même, noté $\text{Map}(E, E)$, de la topologie compacte-ouverte,
c'est-à-dire de la topologie pour laquelle, étant donné $f : E \to E$, un
système fondamental de voisinages de $f$ est formé des ensembles
$U_S(f) = \{f' : E \to E \mid f'|_S = f|_S\}$, où $S \subset E$ est fini.
Pour cette topologie, l'action
$$
\text{Map}(E, E) \times E \longrightarrow E,\quad
(f, e) \longmapsto f(e)
$$
est continue lorsque $E$ est muni de la topologie discrète. Si $X$ est un
espace topologique et si $X \times E \to E$ est une application continue,
alors l'application $X \to \text{Map}(E, E)$ est continue. Autrement dit, la
topologie compacte-ouverte est la topologie la moins fine pour laquelle
l'application d'« action » affichée ci-dessus est continue. La composition
$$
\text{Map}(E, E) \times \text{Map}(E, E) \to \text{Map}(E, E)
$$
est elle aussi continue (comme on le vérifie aisément à l'aide de la
description précédente des voisinages). Enfin, si
$\text{Aut}(E) \subset \text{Map}(E, E)$ est la partie des applications
bijectives, alors l'inversion $i : \text{Aut}(E) \to \text{Aut}(E)$,
$f \mapsto f^{-1}$, est également continue. En effet, si $S \subset E$ est
fini, alors $i^{-1}(U_S(f^{-1})) = U_{f^{-1}(S)}(f)$. Ainsi,
$\text{Aut}(E)$ est un groupe topologique au sens de la
Définition \ref{definition-topological-group}.
\end{example}

\begin{lemma}
\label{lemma-topological-group-limits}
La catégorie des groupes topologiques possède des limites, et les foncteurs
d'oubli vers (a) la catégorie des espaces topologiques et (b) la catégorie des
groupes commutent aux limites.
\end{lemma}

\begin{proof}
Il suffit d'établir ces assertions pour les produits et les égalisateurs ;
voir Catégories, Lemme
\ref{categories-lemma-limits-products-equalizers}. Soient $G_i$, $i \in I$,
une famille de groupes topologiques. Munissons le produit usuel
$G = \prod G_i$ de la topologie produit. Puisque
$G \times G = \prod (G_i \times G_i)$ comme espace topologique (les produits
commutent entre eux dans toute catégorie), la multiplication sur $G$ est
continue. Il en va de même de l'inversion. Soient
$a, b : G \to H$ deux homomorphismes de groupes topologiques. Pour
l'égalisateur, il suffit alors de prendre l'égalisateur de $a$ et $b$ dans la
catégorie des espaces topologiques : il s'agit aussi de leur égalisateur dans
la catégorie des groupes, muni de la topologie induite.
\end{proof}

\begin{lemma}
\label{lemma-profinite-group}
Soit $G$ un groupe topologique. Les conditions suivantes sont équivalentes :
\begin{enumerate}
\item $G$, considéré comme espace topologique, est profini,
\item $G$ est la limite d'un diagramme de groupes topologiques discrets finis,
\item $G$ est une limite cofiltrante de groupes topologiques discrets finis.
\end{enumerate}
\end{lemma}

\begin{proof}
Nous disposons du résultat correspondant pour les espaces topologiques ; voir
le Lemme \ref{lemma-profinite}. Combiné au
Lemme \ref{lemma-topological-group-limits}, ce résultat montre qu'il suffit de
prouver que (1) implique (3).

\medskip\noindent
Montrons d'abord que tout voisinage $E$ de l'élément neutre $e$ contient un
sous-groupe ouvert. En effet, puisque $G$ est une limite cofiltrante d'espaces
topologiques discrets finis (Lemme \ref{lemma-profinite}), on peut choisir une
application continue $f : G \to T$ vers un espace discret fini $T$ telle que
$f^{-1}(f(\{e\})) \subset E$. Considérons
$$
H = \{g \in G \mid f(gg') = f(g')\text{ pour tout }g' \in G\}
$$
Il s'agit d'un sous-groupe de $G$ contenu dans $E$. Il suffit donc de montrer
que $H$ est ouvert. Fixons $t \in T$ et posons $W = f^{-1}(\{t\})$.
La partie $W \subset G$ est ouverte et fermée, donc en particulier
quasi-compacte. Pour tout $w \in W$, il existe des voisinages ouverts
$e \in U_w \subset G$ et $w \in U'_w \subset W$ tels que
$U_wU'_w \subset W$. Par quasi-compacité, on peut trouver
$w_1, \ldots, w_n$ tels que $W = \bigcup U'_{w_i}$. Alors
$U_t = U_{w_1} \cap \ldots \cap U_{w_n}$ est un voisinage ouvert de $e$ tel
que $f(gw) = t$ pour tout $w \in W$. Comme $T$ est fini,
$\bigcap_{t \in T} U_t \subset H$ est un voisinage ouvert de $e$. Puisque
$H \subset G$ est un sous-groupe, il s'ensuit que $H$ est ouvert.

\medskip\noindent
Supposons que $H \subset G$ soit un sous-groupe ouvert. Puisque $G$ est
quasi-compact, l'indice de $H$ dans $G$ est fini. Écrivons
$G = Hg_1 \cup \ldots \cup Hg_n$. Alors
$N = \bigcap_{i = 1, \ldots, n} g_iHg_i^{-1}$ est un sous-groupe normal
ouvert contenu dans $H$. Comme $N$ est lui aussi
d'indice fini, $G \to G/N$ est une surjection vers un groupe topologique
discret fini.

\medskip\noindent
Considérons l'application
$$
G \longrightarrow \lim_{N \subset G\text{ ouvert et normal}} G/N
$$
Nous affirmons que cette application est un isomorphisme de groupes
topologiques. Cela achève la preuve du lemme, car le système qui figure à
droite est cofiltrant (l'intersection de deux sous-groupes normaux ouverts est
ouverte et normale). L'application est continue puisque chaque $G \to G/N$
l'est.
Elle est injective parce que $G$ est séparé et que tout voisinage de $e$
contient un tel $N$, d'après les arguments précédents. Elle est surjective par
le Lemme \ref{lemma-intersection-closed-in-quasi-compact}. Enfin, le
Lemme \ref{lemma-bijective-map} montre que c'est un homéomorphisme.
\end{proof}

\begin{definition}
\label{definition-profinite-group}
Un groupe topologique est appelé un {\it groupe profini} s'il satisfait aux
conditions équivalentes du Lemme \ref{lemma-profinite-group}.
\end{definition}

\noindent
Si $G_1 \to G_2 \to G_3 \to \ldots$ est un système de groupes topologiques,
alors sa colimite $G = \colim G_n$ dans la catégorie des groupes topologiques
(Lemme \ref{lemma-topological-group-colimits}) diffère en général de sa
colimite dans la catégorie des espaces topologiques
(Lemme \ref{lemma-colimits}), bien que ces deux colimites aient le même
ensemble sous-jacent. Voir Exemples,
section \ref{examples-section-colimit-topology}.

\begin{lemma}
\label{lemma-topological-group-colimits}
La catégorie des groupes topologiques possède des colimites, et le foncteur
d'oubli vers la catégorie des groupes commute aux colimites.
\end{lemma}

\begin{proof}
Pour démontrer l'existence des colimites, nous utiliserons l'argument de
Catégories, Remarque \ref{categories-remark-how-to-use-it}. Supposons donc que
$\mathcal{I} \to \textit{Top}$, $i \mapsto G_i$, soit un foncteur à valeurs
dans la catégorie $\textit{TopGroup}$ des groupes topologiques. Considérons
$$
F : \textit{TopGroup} \longrightarrow \textit{Ensembles},\quad
H \longmapsto \lim_\mathcal{I} \Mor_{\textit{TopGroup}}(G_i, H)
$$
Ce foncteur commute aux limites. De plus, étant donné un groupe topologique
$H$ et un élément $(\varphi_i : G_i \to H)$ de $F(H)$, il existe un
sous-groupe $H' \subset H$ de cardinal au plus $|\coprod G_i|$ (le coproduit
étant pris dans la catégorie des groupes, c'est-à-dire le produit libre des
$G_i$) tel que les morphismes $\varphi_i$ soient à valeurs dans $H'$. On peut
en effet choisir le sous-groupe engendré par les images des $\varphi_i$, muni
de la topologie induite. Les hypothèses de Catégories, Lemme
\ref{categories-lemma-a-version-of-brown}, sont donc satisfaites, et l'on
obtient un groupe topologique $G$ qui représente le foncteur $F$, ce qui
signifie précisément que $G$ est la colimite du diagramme $i \mapsto G_i$.

\medskip\noindent
Pour montrer que le foncteur d'oubli vers les groupes commute aux colimites,
nous utilisons Catégories, Lemme \ref{categories-lemma-adjoint-exact}. En
effet, le foncteur d'oubli possède un adjoint à droite : le foncteur qui associe
à un groupe le groupe topologique chaotique (ou indiscret) correspondant.
\end{proof}

\begin{definition}
\label{definition-topological-ring}
Un {\it anneau topologique} est un anneau $R$ muni d'une topologie telle que
l'addition $R \times R \to R$, $(x, y) \mapsto x + y$, et la multiplication
$R \times R \to R$, $(x, y) \mapsto xy$, soient continues.
Un {\it homomorphisme d'anneaux topologiques} est un homomorphisme d'anneaux
qui est continu.
\end{definition}

\noindent
Dans le projet Champs, les anneaux sont commutatifs et possèdent une unité,
notée $1$. Si $R$ est un anneau topologique, alors $(R, +)$ est un groupe
topologique, puisque
$x \mapsto -x$ est continue. Si $R$ est un anneau topologique et
$R' \subset R$ un sous-anneau, alors $R'$, muni de la topologie induite, est
un anneau topologique. Si $R$ est un anneau topologique et $R \to R'$ une
surjection d'anneaux, alors $R'$, muni de la topologie quotient, est un anneau
topologique.

\begin{lemma}
\label{lemma-topological-ring-limits}
La catégorie des anneaux topologiques possède des limites, et les foncteurs
d'oubli vers (a) la catégorie des espaces topologiques et (b) la catégorie des
anneaux commutent aux limites.
\end{lemma}

\begin{proof}
Il suffit d'établir ces assertions pour les produits et les égalisateurs ;
voir Catégories, Lemme
\ref{categories-lemma-limits-products-equalizers}. Soient $R_i$, $i \in I$,
une famille d'anneaux topologiques. Munissons le produit usuel
$R = \prod R_i$ de la topologie produit. Puisque
$R \times R = \prod (R_i \times R_i)$ comme espace topologique (les produits
commutent entre eux dans toute catégorie), l'addition et la multiplication
sur $R$ sont continues. Soient $a, b : R \to R'$ deux homomorphismes
d'anneaux topologiques. Pour l'égalisateur, il suffit alors de prendre
l'égalisateur de $a$ et $b$ dans la catégorie des espaces topologiques : il
s'agit aussi de leur égalisateur dans la catégorie des anneaux, muni de la
topologie induite.
\end{proof}

\begin{lemma}
\label{lemma-topological-ring-colimits}
La catégorie des anneaux topologiques possède des colimites, et le foncteur
d'oubli vers la catégorie des anneaux commute aux colimites.
\end{lemma}

\begin{proof}
Le même argument que dans la preuve du
Lemme \ref{lemma-topological-group-colimits} démontre l'existence des
colimites. Pour montrer que le foncteur d'oubli vers les anneaux commute aux
colimites, nous utilisons Catégories, Lemme
\ref{categories-lemma-adjoint-exact}. En effet, le foncteur d'oubli possède un
adjoint à droite : le foncteur qui associe à un anneau l'anneau topologique
chaotique (ou indiscret) correspondant.
\end{proof}

\begin{definition}
\label{definition-topological-module}
Soit $R$ un anneau topologique. Un {\it module topologique} est un $R$-module
$M$ muni d'une topologie telle que l'addition $M \times M \to M$ et la
multiplication par les scalaires $R \times M \to M$ soient continues.
Un {\it homomorphisme de modules topologiques} est un homomorphisme de modules
qui est continu.
\end{definition}

\noindent
Si $R$ est un anneau topologique et $M$ un module topologique, alors
$(M, +)$ est un groupe topologique, puisque $x \mapsto -x$ est continue.
Si $R$ est un anneau topologique, $M$ un module topologique et
$M' \subset M$ un sous-module, alors $M'$, muni de la topologie induite, est
un module topologique. Si $R$ est un anneau topologique, $M$ un module
topologique et $M \to M'$ une surjection de modules, alors $M'$, muni de la
topologie quotient, est un module topologique.

\begin{lemma}
\label{lemma-topological-module-limits}
Soit $R$ un anneau topologique. La catégorie des modules topologiques sur $R$
possède des limites, et les foncteurs d'oubli vers (a) la catégorie des espaces
topologiques et (b) la catégorie des $R$-modules commutent aux limites.
\end{lemma}

\begin{proof}
Il suffit d'établir ces assertions pour les produits et les égalisateurs ;
voir Catégories, Lemme
\ref{categories-lemma-limits-products-equalizers}. Soient $M_i$, $i \in I$,
une famille de modules topologiques sur $R$. Munissons le produit usuel
$M = \prod M_i$ de la topologie produit. Puisque
$M \times M = \prod (M_i \times M_i)$ comme espace topologique (les produits
commutent entre eux dans toute catégorie), l'addition sur $M$ est continue. Il
en va de même de la multiplication par les scalaires $R \times M \to M$.
Soient $a, b : M \to M'$ deux homomorphismes de modules topologiques sur $R$.
Pour l'égalisateur, il suffit alors de prendre l'égalisateur de $a$ et $b$
dans la catégorie des espaces topologiques : il s'agit aussi de leur
égalisateur dans la catégorie des modules, muni de la topologie induite.
\end{proof}

\begin{lemma}
\label{lemma-topological-module-colimits}
Soit $R$ un anneau topologique. La catégorie des modules topologiques sur $R$
possède des colimites, et le foncteur d'oubli vers la catégorie des modules sur
$R$ commute aux colimites.
\end{lemma}

\begin{proof}
Le même argument que dans la preuve du
Lemme \ref{lemma-topological-group-colimits} démontre l'existence des
colimites. Pour montrer que le foncteur d'oubli vers les $R$-modules commute
aux colimites, nous utilisons Catégories, Lemme
\ref{categories-lemma-adjoint-exact}. En effet, le foncteur d'oubli possède un
adjoint à droite : le foncteur qui associe à un module le module topologique
chaotique (ou indiscret) correspondant.
\end{proof}





\begin{multicols}{2}[\section{Autres chapitres}]
\noindent
Préliminaires
\begin{enumerate}
\item \hyperref[introduction-section-phantom]{Introduction}
\item \hyperref[conventions-section-phantom]{Conventions}
\item \hyperref[sets-section-phantom]{Théorie des ensembles}
\item \hyperref[categories-section-phantom]{Catégories}
\item \hyperref[topology-section-phantom]{Topologie}
\item \hyperref[sheaves-section-phantom]{Faisceaux sur les espaces}
\item \hyperref[sites-section-phantom]{Sites et faisceaux}
\item \hyperref[stacks-section-phantom]{Champs}
\item \hyperref[fields-section-phantom]{Corps}
\item \hyperref[algebra-section-phantom]{Algèbre commutative}
\item \hyperref[brauer-section-phantom]{Groupes de Brauer}
\item \hyperref[homology-section-phantom]{Algèbre homologique}
\item \hyperref[derived-section-phantom]{Catégories dérivées}
\item \hyperref[simplicial-section-phantom]{Méthodes simpliciales}
\item \hyperref[more-algebra-section-phantom]{Compléments d'algèbre}
\item \hyperref[smoothing-section-phantom]{Lissage des morphismes d'anneaux}
\item \hyperref[modules-section-phantom]{Faisceaux de modules}
\item \hyperref[sites-modules-section-phantom]{Modules sur les sites}
\item \hyperref[injectives-section-phantom]{Objets injectifs}
\item \hyperref[cohomology-section-phantom]{Cohomologie des faisceaux}
\item \hyperref[sites-cohomology-section-phantom]{Cohomologie sur les sites}
\item \hyperref[dga-section-phantom]{Algèbre différentielle graduée}
\item \hyperref[dpa-section-phantom]{Algèbre à puissances divisées}
\item \hyperref[sdga-section-phantom]{Faisceaux différentiels gradués}
\item \hyperref[hypercovering-section-phantom]{Hyperrecouvrements}
\end{enumerate}
Schémas
\begin{enumerate}
\setcounter{enumi}{25}
\item \hyperref[schemes-section-phantom]{Schémas}
\item \hyperref[constructions-section-phantom]{Constructions de schémas}
\item \hyperref[properties-section-phantom]{Propriétés des schémas}
\item \hyperref[morphisms-section-phantom]{Morphismes de schémas}
\item \hyperref[coherent-section-phantom]{Cohomologie des schémas}
\item \hyperref[divisors-section-phantom]{Diviseurs}
\item \hyperref[limits-section-phantom]{Limites de schémas}
\item \hyperref[varieties-section-phantom]{Variétés}
\item \hyperref[topologies-section-phantom]{Topologies sur les schémas}
\item \hyperref[descent-section-phantom]{Descente}
\item \hyperref[perfect-section-phantom]{Catégories dérivées de schémas}
\item \hyperref[more-morphisms-section-phantom]{Compléments sur les morphismes}
\item \hyperref[flat-section-phantom]{Compléments sur la platitude}
\item \hyperref[groupoids-section-phantom]{Schémas en groupoïdes}
\item \hyperref[more-groupoids-section-phantom]{Compléments sur les schémas en groupoïdes}
\item \hyperref[etale-section-phantom]{Morphismes étales de schémas}
\end{enumerate}
Thèmes de théorie des schémas
\begin{enumerate}
\setcounter{enumi}{41}
\item \hyperref[chow-section-phantom]{Homologie de Chow}
\item \hyperref[intersection-section-phantom]{Théorie de l'intersection}
\item \hyperref[pic-section-phantom]{Schémas de Picard des courbes}
\item \hyperref[weil-section-phantom]{Théories cohomologiques de Weil}
\item \hyperref[adequate-section-phantom]{Modules adéquats}
\item \hyperref[dualizing-section-phantom]{Complexes dualisants}
\item \hyperref[duality-section-phantom]{Dualité pour les schémas}
\item \hyperref[discriminant-section-phantom]{Discriminants et différentes}
\item \hyperref[derham-section-phantom]{Cohomologie de de Rham}
\item \hyperref[local-cohomology-section-phantom]{Cohomologie locale}
\item \hyperref[algebraization-section-phantom]{Géométrie algébrique et formelle}
\item \hyperref[curves-section-phantom]{Courbes algébriques}
\item \hyperref[resolve-section-phantom]{Résolution des surfaces}
\item \hyperref[models-section-phantom]{Réduction semi-stable}
\item \hyperref[functors-section-phantom]{Foncteurs et morphismes}
\item \hyperref[equiv-section-phantom]{Catégories dérivées de variétés}
\item \hyperref[pione-section-phantom]{Groupes fondamentaux de schémas}
\item \hyperref[etale-cohomology-section-phantom]{Cohomologie étale}
\item \hyperref[crystalline-section-phantom]{Cohomologie cristalline}
\item \hyperref[proetale-section-phantom]{Cohomologie pro-étale}
\item \hyperref[relative-cycles-section-phantom]{Cycles relatifs}
\item \hyperref[more-etale-section-phantom]{Compléments de cohomologie étale}
\item \hyperref[trace-section-phantom]{La formule des traces}
\end{enumerate}
Espaces algébriques
\begin{enumerate}
\setcounter{enumi}{64}
\item \hyperref[spaces-section-phantom]{Espaces algébriques}
\item \hyperref[spaces-properties-section-phantom]{Propriétés des espaces algébriques}
\item \hyperref[spaces-morphisms-section-phantom]{Morphismes d'espaces algébriques}
\item \hyperref[decent-spaces-section-phantom]{Espaces algébriques décents}
\item \hyperref[spaces-cohomology-section-phantom]{Cohomologie des espaces algébriques}
\item \hyperref[spaces-limits-section-phantom]{Limites d'espaces algébriques}
\item \hyperref[spaces-divisors-section-phantom]{Diviseurs sur les espaces algébriques}
\item \hyperref[spaces-over-fields-section-phantom]{Espaces algébriques sur des corps}
\item \hyperref[spaces-topologies-section-phantom]{Topologies sur les espaces algébriques}
\item \hyperref[spaces-descent-section-phantom]{Descente et espaces algébriques}
\item \hyperref[spaces-perfect-section-phantom]{Catégories dérivées d'espaces}
\item \hyperref[spaces-more-morphisms-section-phantom]{Compléments sur les morphismes d'espaces}
\item \hyperref[spaces-flat-section-phantom]{Platitude sur les espaces algébriques}
\item \hyperref[spaces-groupoids-section-phantom]{Groupoïdes dans les espaces algébriques}
\item \hyperref[spaces-more-groupoids-section-phantom]{Compléments sur les groupoïdes dans les espaces}
\item \hyperref[bootstrap-section-phantom]{Amorçage}
\item \hyperref[spaces-pushouts-section-phantom]{Sommes amalgamées d'espaces algébriques}
\end{enumerate}
Thèmes de géométrie
\begin{enumerate}
\setcounter{enumi}{81}
\item \hyperref[spaces-chow-section-phantom]{Groupes de Chow des espaces}
\item \hyperref[groupoids-quotients-section-phantom]{Quotients de groupoïdes}
\item \hyperref[spaces-more-cohomology-section-phantom]{Compléments sur la cohomologie des espaces}
\item \hyperref[spaces-simplicial-section-phantom]{Espaces simpliciaux}
\item \hyperref[spaces-duality-section-phantom]{Dualité pour les espaces}
\item \hyperref[formal-spaces-section-phantom]{Espaces algébriques formels}
\item \hyperref[restricted-section-phantom]{Algébrisation des espaces formels}
\item \hyperref[spaces-resolve-section-phantom]{Résolution des surfaces revisitée}
\end{enumerate}
Théorie des déformations
\begin{enumerate}
\setcounter{enumi}{89}
\item \hyperref[formal-defos-section-phantom]{Théorie formelle des déformations}
\item \hyperref[defos-section-phantom]{Théorie des déformations}
\item \hyperref[cotangent-section-phantom]{Le complexe cotangent}
\item \hyperref[examples-defos-section-phantom]{Problèmes de déformation}
\end{enumerate}
Champs algébriques
\begin{enumerate}
\setcounter{enumi}{93}
\item \hyperref[algebraic-section-phantom]{Champs algébriques}
\item \hyperref[examples-stacks-section-phantom]{Exemples de champs}
\item \hyperref[stacks-sheaves-section-phantom]{Faisceaux sur les champs algébriques}
\item \hyperref[criteria-section-phantom]{Critères de représentabilité}
\item \hyperref[artin-section-phantom]{Axiomes d'Artin}
\item \hyperref[quot-section-phantom]{Espaces de Quot et de Hilbert}
\item \hyperref[stacks-properties-section-phantom]{Propriétés des champs algébriques}
\item \hyperref[stacks-morphisms-section-phantom]{Morphismes de champs algébriques}
\item \hyperref[stacks-limits-section-phantom]{Limites de champs algébriques}
\item \hyperref[stacks-cohomology-section-phantom]{Cohomologie des champs algébriques}
\item \hyperref[stacks-perfect-section-phantom]{Catégories dérivées de champs}
\item \hyperref[stacks-introduction-section-phantom]{Introduction aux champs algébriques}
\item \hyperref[stacks-more-morphisms-section-phantom]{Compléments sur les morphismes de champs}
\item \hyperref[stacks-geometry-section-phantom]{La géométrie des champs}
\end{enumerate}
Thèmes de théorie des espaces de modules
\begin{enumerate}
\setcounter{enumi}{107}
\item \hyperref[moduli-section-phantom]{Champs de modules}
\item \hyperref[moduli-curves-section-phantom]{Espaces de modules de courbes}
\end{enumerate}
Divers
\begin{enumerate}
\setcounter{enumi}{109}
\item \hyperref[examples-section-phantom]{Exemples}
\item \hyperref[exercises-section-phantom]{Exercices}
\item \hyperref[guide-section-phantom]{Guide bibliographique}
\item \hyperref[desirables-section-phantom]{Desiderata}
\item \hyperref[coding-section-phantom]{Style de programmation}
\item \hyperref[obsolete-section-phantom]{Obsolète}
\item \hyperref[fdl-section-phantom]{Licence de documentation libre GNU}
\item \hyperref[index-section-phantom]{Index généré automatiquement}
\end{enumerate}
\end{multicols}


\bibliography{my}
\bibliographystyle{amsalpha}

\end{document}
